\documentclass[11pt,a4paper]{article}
\usepackage[margin=24mm,headheight=15pt]{geometry}
\usepackage[T1]{fontenc}
\usepackage[utf8]{inputenc}
\usepackage{lmodern,microtype}
\usepackage{amsmath,amssymb,amsthm,mathtools}
\usepackage{booktabs,longtable,array,enumitem,fancyvrb}
\usepackage{xcolor,fancyhdr}
\usepackage[colorlinks=true,linkcolor=blue!40!black,citecolor=blue!40!black,urlcolor=blue!50!black,bookmarksdepth=1]{hyperref}
\newtheorem{theorem}{Theorem}[section]
\newtheorem{proposition}[theorem]{Proposition}
\newtheorem{lemma}[theorem]{Lemma}
\newtheorem{corollary}[theorem]{Corollary}
\theoremstyle{definition}
\newtheorem{definition}[theorem]{Definition}
\newtheorem{example}[theorem]{Example}
\theoremstyle{remark}
\newtheorem{remark}[theorem]{Remark}
\numberwithin{equation}{section}
\newcommand{\Qbar}{\overline{\mathbb Q}}
\newcommand{\Qb}{\overline{\mathbb Q}}
\newcommand{\Res}{\operatorname{Res}}
\newcommand{\F}[2]{{}_{#1}F_{#2}}
\newcommand{\Ein}{\operatorname{Ein}}
\newcommand{\Hom}{\operatorname{Hom}}
\newcommand{\trdeg}{\operatorname{trdeg}}
\setlength{\emergencystretch}{2em}
\setlength{\parskip}{3pt}
\setlength{\parindent}{1.3em}
\allowdisplaybreaks[2]
\setcounter{tocdepth}{1}
\pagestyle{fancy}\fancyhf{}
\fancyhead[L]{\small Beta periods and Gamma transfers}
\fancyhead[R]{\small KEIO}
\fancyfoot[C]{\thepage}
\hypersetup{pdftitle={Explicit Transcendence Criteria for Beta-Type Integrals and Gamma Transfers},pdfauthor={KEIO},pdfsubject={A consolidated research manuscript with complete proofs and explicit open conditions}}
\title{\bfseries Explicit Transcendence Criteria for Beta-Type Integrals\\[5pt]
\large Integral Quotients, Cyclotomic Fields, Statistical Quantiles,\\and Transfers to the Gamma Function}
\author{KEIO}
\date{11 October 2026}
\begin{document}
\maketitle
\begin{abstract}
We consolidate a sequence of explicit transcendence arguments for algebraic
integrals, normalized hypergeometric values, Gamma quotients, and inverse
distribution functions. The arithmetic inputs are classical theorems on
logarithms, beta values, algebraic groups, and $1$-periods. Applications include
all forty-four selected entries of the DBKZ Log2 catalogue, twenty-two
additional Log1 entries, six double-integral constants, and seven
triple-integral quotients. Uniform families are obtained from generalized
Jacobians, cyclotomic eigenspaces, and torsion-packet classifications.
Polynomially weighted incomplete beta integrals admit an exact exceptional
hyperplane, while related arguments classify several statistical quantiles
and probability intervals. Balanced Gamma quotients at denominators $21$
and $33$ yield explicit number fields and counterexamples to proposed
dimension bounds. We then audit transfers to individual Gamma values and
Euler's constant. A positive-polynomial kernel yields an explicit fourth-power
degree obstruction, and a rapidly convergent algebraically independent
sequence shows why finite-point transcendence does not pass to a limit.
The remaining exponential-period and mixed-Gevrey assertions are stated
separately as unproved conditions. Neither the irrationality of Euler's
constant nor the transcendence of $\Gamma(1/5)$ is proved.
\end{abstract}
\noindent\textbf{Keywords.} Transcendence; beta integrals; generalized Jacobians;
hypergeometric functions; torsion; Gamma values; quantiles; Euler's constant.

\section*{Scope and organization}
\addcontentsline{toc}{section}{Scope and organization}
This manuscript organizes the mathematical content of a research dialogue
into one proof document. Earlier special cases are consolidated under later
uniform results where possible. Deep external theorems are cited as inputs;
their application here is not presented as their invention. A public output
labelled ``unidentified'' does not establish a literature-wide open problem.
No historical-priority or external-peer-review claim is made.

Counts in the catalogues refer to exact parameter tuples, unless explicitly
stated otherwise. They do not assert that all resulting real numbers are
distinct or algebraically independent. Algebraic cases, endpoint exceptions,
and hypotheses on positivity and denominator structure are retained.
The quantitative irrationality measures and divisibility conjectures attached
to the original DBKZ constructions are not established by qualitative
transcendence alone.

All real powers in real integrals are taken on their positive branches.
The symbols $\Gamma(z)$, $\gamma=-\Gamma'(1)$, and
$\delta=\int_0^\infty e^{-t}/(1+t)\,dt$ denote different objects.
The functional transcendence of $\Gamma$ must not be confused with arithmetic
transcendence of one value or derivative.


\clearpage
\begingroup
\small
\setlength{\parskip}{0pt}
\tableofcontents
\endgroup
\clearpage

\part{One-dimensional integral quotients}
\section{The one-dimensional DBKZ family and all forty-four entries}
\label{dbkz-sec:family}
For rational $a,b>-1$, real algebraic $c>-1$, and rational $d$, put
\[
 C(a,b,c,d)=\frac{\int_0^1x^a(1-x)^b(1+cx)^{-d-1}\,dx}
 {\int_0^1x^a(1-x)^b(1+cx)^{-d}\,dx}.
\]
Positive real branches are understood. Both integrals converge and the
denominator is positive. The final unidentified block of the DBKZ
\texttt{Log2} output contains the forty-four tuples in
Table~\ref{dbkz-tab:44} \cite{dbkz-source}. In the proofs below every one
of these forty-four values is transcendental. ``Unidentified'' refers to
that output, and is not a literature-wide priority claim. The count is
of parameter tuples; pairwise distinctness and algebraic independence
of the resulting values are not asserted.

The earlier computational status was $26$ completed and $18$ outstanding.
The latter were rows $6,7,29,30,\ldots,44$. The present organization uses
the strongest common arguments, so it supersedes some earlier isolated
elliptic and real-multiplication proofs. The method codes are: L,
logarithms; I, incomplete second-kind periods; D, the uniform diagonal
theorem; E, an explicit elliptic quotient; B, a beta quotient; and S,
a simple genus-four Jacobian. All rows have final status O, meaning a
proof is supplied, rather than a numerical identification.

{\small
\begin{longtable}{@{}rcccccc@{}}
\caption{The exact forty-four tuples, in the source order.}\label{dbkz-tab:44}\\
\toprule Row&$a$&$b$&$c$&$d$&Method&Status\\\midrule\endfirsthead
\toprule Row&$a$&$b$&$c$&$d$&Method&Status\\\midrule\endhead
1&$\tfrac{1}{2}$&$\tfrac{1}{2}$&$\tfrac{1}{2}$&$\tfrac{1}{2}$&D&O\\
2&$0$&$\tfrac{1}{3}$&$1$&$\tfrac{2}{3}$&I&O\\
3&$0$&$\tfrac{1}{3}$&$\tfrac{1}{3}$&$\tfrac{2}{3}$&I&O\\
4&$0$&$\tfrac{2}{3}$&$1$&$\tfrac{1}{3}$&I&O\\
5&$0$&$\tfrac{2}{3}$&$\tfrac{1}{3}$&$\tfrac{1}{3}$&I&O\\
6&$\tfrac{1}{3}$&$\tfrac{1}{3}$&$\tfrac{1}{3}$&$\tfrac{1}{3}$&D&O\\
7&$\tfrac{2}{3}$&$\tfrac{2}{3}$&$\tfrac{1}{3}$&$\tfrac{2}{3}$&D&O\\
8&$0$&$\tfrac{1}{2}$&$1$&$\tfrac{1}{4}$&I&O\\
9&$0$&$\tfrac{1}{2}$&$1$&$\tfrac{3}{4}$&I&O\\
10&$0$&$\tfrac{1}{2}$&$\tfrac{1}{2}$&$\tfrac{1}{4}$&I&O\\
11&$0$&$\tfrac{1}{2}$&$\tfrac{1}{2}$&$\tfrac{3}{4}$&I&O\\
12&$0$&$\tfrac{1}{2}$&$\tfrac{1}{4}$&$0$&L&O\\
13&$0$&$\tfrac{1}{2}$&$\tfrac{1}{4}$&$\tfrac{1}{4}$&I&O\\
14&$0$&$\tfrac{1}{2}$&$\tfrac{1}{4}$&$\tfrac{3}{4}$&I&O\\
15&$0$&$\tfrac{1}{4}$&$\tfrac{1}{4}$&$\tfrac{1}{2}$&I&O\\
16&$0$&$\tfrac{3}{4}$&$\tfrac{1}{4}$&$\tfrac{1}{2}$&I&O\\
17&$\tfrac{1}{2}$&$\tfrac{1}{2}$&$\tfrac{1}{4}$&$\tfrac{1}{2}$&D&O\\
18&$\tfrac{1}{2}$&$\tfrac{1}{2}$&$\tfrac{1}{4}$&$\tfrac{1}{4}$&E&O\\
19&$\tfrac{1}{4}$&$\tfrac{1}{4}$&$\tfrac{1}{4}$&$\tfrac{1}{4}$&D&O\\
20&$\tfrac{3}{4}$&$\tfrac{3}{4}$&$\tfrac{1}{4}$&$\tfrac{3}{4}$&D&O\\
21&$0$&$\tfrac{1}{5}$&$\tfrac{1}{5}$&$0$&L&O\\
22&$0$&$\tfrac{1}{5}$&$\tfrac{1}{5}$&$\tfrac{3}{5}$&I&O\\
23&$0$&$\tfrac{2}{5}$&$\tfrac{1}{5}$&$0$&L&O\\
24&$0$&$\tfrac{2}{5}$&$\tfrac{1}{5}$&$\tfrac{1}{5}$&I&O\\
25&$0$&$\tfrac{3}{5}$&$\tfrac{1}{5}$&$0$&L&O\\
26&$0$&$\tfrac{3}{5}$&$\tfrac{1}{5}$&$\tfrac{4}{5}$&I&O\\
27&$0$&$\tfrac{4}{5}$&$\tfrac{1}{5}$&$0$&L&O\\
28&$0$&$\tfrac{4}{5}$&$\tfrac{1}{5}$&$\tfrac{2}{5}$&I&O\\
29&$\tfrac{1}{5}$&$\tfrac{1}{5}$&$1$&$\tfrac{3}{5}$&B&O\\
30&$\tfrac{1}{5}$&$\tfrac{1}{5}$&$\tfrac{1}{5}$&$\tfrac{1}{5}$&D&O\\
31&$\tfrac{1}{5}$&$\tfrac{1}{5}$&$\tfrac{1}{5}$&$\tfrac{3}{5}$&S&O\\
32&$\tfrac{1}{5}$&$\tfrac{1}{5}$&$\tfrac{2}{5}$&$\tfrac{3}{5}$&S&O\\
33&$\tfrac{1}{5}$&$\tfrac{2}{5}$&$1$&$\tfrac{4}{5}$&B&O\\
34&$\tfrac{1}{5}$&$\tfrac{2}{5}$&$\tfrac{1}{5}$&$\tfrac{4}{5}$&S&O\\
35&$\tfrac{1}{5}$&$\tfrac{2}{5}$&$\tfrac{2}{5}$&$\tfrac{4}{5}$&S&O\\
36&$\tfrac{1}{5}$&$\tfrac{3}{5}$&$1$&$\tfrac{2}{5}$&B&O\\
37&$\tfrac{1}{5}$&$\tfrac{3}{5}$&$\tfrac{1}{5}$&$\tfrac{2}{5}$&S&O\\
38&$\tfrac{1}{5}$&$\tfrac{3}{5}$&$\tfrac{2}{5}$&$\tfrac{2}{5}$&S&O\\
39&$\tfrac{2}{5}$&$\tfrac{2}{5}$&$\tfrac{1}{5}$&$\tfrac{2}{5}$&D&O\\
40&$\tfrac{2}{5}$&$\tfrac{4}{5}$&$1$&$\tfrac{3}{5}$&B&O\\
41&$\tfrac{2}{5}$&$\tfrac{4}{5}$&$\tfrac{1}{5}$&$\tfrac{3}{5}$&S&O\\
42&$\tfrac{2}{5}$&$\tfrac{4}{5}$&$\tfrac{2}{5}$&$\tfrac{3}{5}$&S&O\\
43&$\tfrac{3}{5}$&$\tfrac{3}{5}$&$\tfrac{1}{5}$&$\tfrac{3}{5}$&D&O\\
44&$\tfrac{4}{5}$&$\tfrac{4}{5}$&$\tfrac{1}{5}$&$\tfrac{4}{5}$&D&O\\
\bottomrule
\end{longtable}}

\subsection{The fixed-cycle period input}
For an abelian variety $A/\overline{\mathbb Q}$ and a rational cycle
$\gamma$, let $B_\gamma$ be the smallest abelian subvariety with
$\gamma\in H_1(B_\gamma,\mathbb Q)$, in the rational isogeny category.
The subgroup theorem for $1$-periods implies
\begin{equation}\label{dbkz-eq:periodkernel}
 \int_\gamma\alpha=0\quad\Longleftrightarrow\quad
 \alpha|_{B_\gamma}=0,
 \qquad \alpha\in H^1_{\rm dR}(A/\overline{\mathbb Q}).
\end{equation}
This is the pure-abelian specialization of
\cite[Theorem 9.7]{dbkz-hw}. In particular, if $A$ is simple and
$\gamma\ne0$, two independent algebraic de Rham classes have
$\overline{\mathbb Q}$-linearly independent periods on $\gamma$, so their
ratio is transcendental. The corresponding fixed-cycle statement also
appears in \cite[Satz 1]{dbkz-ww}.

A useful refinement avoids assuming simplicity. Let $\rho$ be a
finite-order automorphism such that every nonzero eigenspace of
$H^0(A,\Omega^1)$ is one-dimensional. Suppose a rational cycle $\gamma$
has nonzero period against every holomorphic eigenform. Then evaluation
on $\gamma$ is injective on each eigenspace of $H^1_{\rm dR}(A)$.
Indeed, $B=\sum_j\rho^jB_\gamma$ is $\rho$-stable. The kernel of
restriction of holomorphic forms to $B$ is stable and cannot contain
an eigenform, by the nonzero-period hypothesis. Thus $B=A$. An
eigenclass with zero $\gamma$-period vanishes on each
$\rho^jB_\gamma$ by \eqref{dbkz-eq:periodkernel}; the surjective addition
map from their product to $A$ then forces that class to vanish.
This criterion concerns the specified cycle, not arbitrary periods of
a reducible Jacobian.

\subsection{All twenty-one entries with first parameter zero}
\label{dbkz-sec:a0}
Suppose first that $0<b,d<1$ and $c>0$ is algebraic. Set
\[
 m=\frac{c}{1+c},\quad T=B_m(b,1-d),\quad
 H=m^b(1-m)^{1-d}.
\]
Two integrations of derivatives give
\[
 I:=\int_0^m t^b(1-t)^{-d}\,dt=\frac{bT-H}{b+1-d},
 \quad
 J:=\int_0^m t^b(1-t)^{-d-1}\,dt=\frac{(1+c)H-bT}{d}.
\]
For example, differentiate $t^b(1-t)^{1-d}$ for the first identity and
$t^{b+1}(1-t)^{-d}$ for the second. Substituting
$t=c(1-x)/(1+c)$ in the original integrals yields
\begin{equation}\label{dbkz-eq:a0mobius}
 C(0,b,c,d)=\frac{b+1-d}{d(1+c)}
       \frac{(1+c)H-bT}{bT-H}.
\end{equation}
The denominator is positive up to a positive factor, and the inner
fractional-linear determinant is $-bcH\ne0$. Consequently the target
is algebraic if and only if $T$ is algebraic.

We use the following precise incomplete-period consequence of
\cite[Theorem 13.9, Summaries 14.11 and 14.16, Corollary 14.22]{dbkz-hw}:
on a curve with simple Jacobian, the integral of a nonexact algebraic
differential of the second kind along one path between distinct algebraic
points is transcendental. The second-kind hypothesis matters: it excludes
the third-kind exception. One way to obtain the statement is that an
algebraic integral would, after subtracting an exact differential, give
a nonexact second-kind differential with zero path integral. The cited
vanishing criterion forces the Abel--Jacobi logarithm of the path to be
zero. The Abel--Jacobi map is injective on points of a positive-genus
curve, contradicting the distinct endpoints.

Write $b=p/q$, $1-d=r/q$ and normalize
\[
 w^q=t^p(1-t)^r,\qquad \omega=\frac{w\,dt}{t(1-t)}.
\]
With $g_0=\gcd(q,p)$, $g_1=\gcd(q,r)$ and
$g_\infty=\gcd(q,p+r)$, the genus and local orders are
\[
 2g-2=q-g_0-g_1-g_\infty,\qquad
 \operatorname{ord}_{0,1,\infty}\omega
 =\left(\frac p{g_0}-1,\frac r{g_1}-1,
           \frac{q-p-r}{g_\infty}-1\right).
\]
The ten beta types occurring here are listed below; repeated values of
$c$ account for the sixteen $d>0$ entries. In this table
$g_\infty=1$. Deck symmetry kills all residues at infinity. If
$k=\operatorname{lcm}(g_0,g_1)$, the positive lift $\delta$ of $[0,1]$
satisfies
$\int_{\delta-\tau^k\delta}\omega=(1-\zeta_q^k)B(b,1-d)\ne0$.
Hence $\omega$ is nonexact. For $q=4$ one must use $k=2$, not $k=1$,
in the four displayed types.
\[
\begin{array}{c|c|c|c|c|c}
 b&d&(q;p,r)&g&\operatorname{ord}_\infty\omega&k\\\hline
1/3&2/3&(3;1,1)&1&0&1\\
2/3&1/3&(3;2,2)&1&-2&1\\
1/2&1/4&(4;2,3)&1&-2&2\\
1/2&3/4&(4;2,1)&1&0&2\\
1/4&1/2&(4;1,2)&1&0&2\\
3/4&1/2&(4;3,2)&1&-2&2\\
1/5&3/5&(5;1,2)&2&1&1\\
2/5&1/5&(5;2,4)&2&-2&1\\
3/5&4/5&(5;3,1)&2&0&1\\
4/5&2/5&(5;4,3)&2&-3&1
\end{array}
\]
The genus-one Jacobians are simple. The genus-two Jacobians are simple
by the three-point prime-cover criterion
\cite[Corollary 19.16]{dbkz-hw}, since $5\not\equiv1\pmod3$.
The lifted path $[0,m]$ has distinct algebraic endpoints. The
incomplete-period theorem and \eqref{dbkz-eq:a0mobius} prove all sixteen
cases.

For the other five entries $d=0$, write $b=p/q\in(0,1)$ and
$A=(1+c)/c>1$. Substitution $1-x=u^q$ and polynomial division give
\[
 C(0,p/q,c,0)=\frac{p+q}{c}\left(AL-\frac1p\right),\quad
 L=\int_0^1\frac{u^{p-1}}{A-u^q}\,du>0.
\]
Partial fractions at the roots $\rho^q=A$ yield
\[
 L=-\frac1q\sum_{\rho^q=A}\rho^{p-q}
                 \operatorname{Log}(1-1/\rho).
\]
All coefficients and logarithm arguments are algebraic, and the
continued branches are unambiguous along the integration path.
Baker's theorem makes this nonzero logarithmic linear form
transcendental; its affine transformation is transcendental as well.
For instance,
$C(0,1/2,1/4,0)=12\sqrt5\,\operatorname{atanh}(1/\sqrt5)-12$.

The same proof yields useful infinite extensions. The $a=0$ result holds
for every positive algebraic $c$ when $d=0$, $0<b<1$ is rational; when
$0<b=d<1$ is rational; when $b,d\in\{1/4,1/2,3/4\}$; and when
$b=r/p,d=v/p$ with $1\le r,v<p$ and an odd prime
$p\not\equiv1\pmod3$. For $b=d=r/q$, substitute
$t=u^q/(1+u^q)$ to reduce $T$ to
$q\int_0^{c^{1/q}}u^{r-1}/(1+u^q)\,du$, again a positive logarithmic
linear form. Off-diagonal prime cases use the same simplicity criterion;
the remaining quarter cases are genus one by the displayed genus formula.

\subsection{The complete rational diagonal}
\label{dbkz-sec:diagonal}
\begin{theorem}\label{dbkz-thm:diagonal}
For every rational $h>-1$ and real algebraic $c>-1$,
$C(h,h,0,h)=1$, while $C(h,h,c,h)$ is transcendental whenever $c\ne0$.
\end{theorem}
\begin{proof}
First let $c>0$ and $h=k/n$ in lowest terms, with $n\ge2$. Put
$A=2+4/c$ and consider the genus-$n-1$ curve
\[
 X_A:\ v^2=t^{2n}-At^n+1,\qquad
 t^n=\frac{cx(1-x)}{1+cx},\quad v=2x+t^n-1.
\]
The deck map $\rho(t,v)=(\zeta_nt,v)$ has distinct eigenvalues on the
holomorphic basis $\omega_j=t^{j-1}dt/v$, $1\le j<n$.
The positive lift $\delta$ of $0\le x\le1$ runs from $(0,-1)$ to
$(0,1)$, through the first positive branch point. The identity
$dx/[x(1-x)]=-n\,dt/(tv)$ gives
\[
 -nc^{-j/n}\int_\delta\omega_j
 =\int_0^1[x(1-x)]^{j/n-1}(1+cx)^{-j/n}\,dx>0.
\]
Thus the closed cycle $\gamma=\delta-\rho\delta$ has nonzero period
against every $\omega_j$. The character-support criterion following
\eqref{dbkz-eq:periodkernel} applies even to reducible or CM fibers.

Set $g=c^{-h}t^k$, $\alpha=g\,dx$, and
$\beta=g\,dx/(1+cx)$. Both have the nontrivial deck character
$\zeta_n^k$. Their possible poles lie at the four deck-fixed branch
points $0,1,-1/c,\infty$, so invariance of residues makes them
second-kind differentials. They are regular at $0,1$, since the local
order of $g\,dx$ there is $n(h+1)-1\ge0$.
If $a\alpha+b\beta$ is exact, character averaging gives a primitive
$gR(x)$ with $R\in\overline{\mathbb Q}(x)$, satisfying
\begin{equation}\label{dbkz-eq:diagexact}
 R'+h\left(\frac1x-\frac1{1-x}-\frac c{1+cx}\right)R
       =a+\frac b{1+cx}.
\end{equation}
No finite pole of $R$ is possible: at a branch point the leading
coefficient involves a nonintegral residue, and at an ordinary point
differentiation raises the pole order. Consequently $R$ is polynomial,
and regularity on the right at $0,1$ forces $R(0)=R(1)=0$.
A nonzero $R$ of degree $M\ge2$ gives a leading term of degree $M-1$
with coefficient $(M+h)\operatorname{lc}(R)\ne0$, whereas the right
side is bounded at infinity. Thus $R=0$ and $a=b=0$.
The two de Rham classes are independent in one character. Their
$\gamma$-periods are $(1-\zeta_n^k)$ times the positive-path integrals.
Characterwise injectivity proves that their ratio is transcendental.

For integral $h=m\ge0$, $C_0(c)=\log(1+c)/c$. If $m\ge1$, put
$z=(c+2)/c>1$. Rodrigues' formula and integration by parts for the
Legendre functions give
\[
 C_m(c)=\frac{2(m+1)}{mc}
          \frac{Q_m(z)}{Q_{m-1}(z)-zQ_m(z)}.
\]
Here $Q_0(z)=\tfrac12\log(1+c)$,
$Q_m=P_mQ_0-W_{m-1}$ with $P_m,W_{m-1}\in\mathbb Q[z]$, and the
common three-term recurrence gives
$P_mQ_{m-1}-P_{m-1}Q_m=1/m$. Hence the displayed ratio is a
nonconstant fractional-linear function of the transcendental
$\log(1+c)$ over $\overline{\mathbb Q}$. Positivity of the original
denominator excludes a pole. This proves the integer case.
Finally, for $-1<c<0$, reflection $x\mapsto1-x$ gives
\[
 C_h(c)=\frac1{1+c}C_h\!\left(-\frac c{1+c}\right),
\]
reducing to positive algebraic parameters.
\end{proof}

There is also an inverse-value consequence. For every real $h>-1$,
$C_h(c)$ decreases strictly from $1$ to $0$ on $(0,\infty)$.
Indeed, with expectation for density proportional to
$[x(1-x)/(1+cx)]^h$ and $Y=(1+cx)^{-1}$, the involution
$x\mapsto(1-x)/(1+cx)$ gives $\mathbb E(Y^2)=1/(1+c)$.
Differentiation then yields
\[
 cC_h'=\frac{h+1}{1+c}-C_h-hC_h^2,\qquad C_h'(0)=-\frac12.
\]
At a first zero of $C_h'$ this equation would imply
$C_h''=-(h+1)/[c(1+c)^2]<0$, a contradiction. Moreover
$C_h^2\le\mathbb E(Y^2)\to0$.
Consequently the unique positive solution of $C_h(c)=r$, for rational
$h>-1$ and algebraic $0<r<1$, is transcendental. Equivalently, the
diagonal theorem classifies the contiguous quotient
\[
 \frac{{}_2F_1(h+1,h+1;2h+2;-c)}
      {{}_2F_1(h,h+1;2h+2;-c)}.
\]

\subsection{The additional genus-two entry has an elliptic quotient}
\label{dbkz-sec:row18}
For every positive algebraic $c$, $Q(c)=C(1/2,1/2,c,1/4)$ is
transcendental; row 18 is $c=1/4$. Here is an explicit reduction which
avoids any simplicity assertion for the genus-two Jacobian. Put
\[
 s=(1+c)^{1/4},\quad A=s^2+s^{-2},\quad B=A+2>4,\quad g=\sqrt B,
 \quad W(t)=\sqrt{t(At^2-t^4-1)}.
\]
Substitutions $U=\sqrt{1+cx}$ and $t=U/s$ transform the denominator
and numerator into
\[
 D=\frac{2s^{7/2}}{c^2}\int_{1/s}^sW(t)\,dt,
 \qquad N=\frac{2s^{3/2}}{c^2}\int_{1/s}^sW(t)t^{-2}\,dt.
\]
Use $W(1/t)=W(t)/t^3$, then $u=t+t^{-1}$ and
$V=W(t)(t-1)/t^2$ on $1<t<s$. The elliptic equation and transformed
integrals are
\[
 V^2=(u-2)(B-u^2),\quad
 D_t=\int_2^g(B-u^2)(u^2-u-1)\frac{du}{V},\quad
 N_t=\int_2^g(B-u^2)\frac{du}{V}.
\]
Set $\Omega=\int_2^gdu/V>0$, $H=\int_2^gu\,du/V$ and $R=H/\Omega$.
For $f=(u-2)(B-u^2)$, direct polynomial identities give
\begin{align*}
 (B-u^2)(u^2-u-1)
 &=\tfrac2{21}\{B(B-5)+(B-2)u\}+2r_D'f+r_Df',\\
 B-u^2&=\tfrac23(B-2u)+2r_N'f+r_Nf',
\end{align*}
where $r_D=(1-2B)/21+u/7+u^2/7$ and $r_N=1/3$.
The exact terms are $2d(r_DV)$ and $2d(r_NV)$ and vanish on integration
because $V$ vanishes at both endpoints. Hence
\begin{equation}\label{dbkz-eq:row18}
 Q(c)=\frac7{s^2}\frac{B-2R}{B(B-5)+(B-2)R}.
\end{equation}
The elliptic de Rham basis $du/V,u\,du/V$ evaluates on the nonzero
cycle obtained by returning on the other sheet to $2\Omega,2H$.
The fixed-cycle theorem makes $R$ transcendental. The determinant of
\eqref{dbkz-eq:row18}, before the nonzero scalar, is $3B(4-B)\ne0$;
its denominator is a nonzero multiple of $D_t>0$. This proves the claim.
If desired, the conventional parameter-form elliptic integrals give
$R=(g+2)E(\mu)/K(\mu)-g$, $\mu=-(g-2)/(g+2)$.

\subsection{Four complete-beta quotients}
\label{dbkz-sec:betafour}
For rows 29, 33, 36 and 40 the following exact identities hold:
\[
\begin{array}{c|c|c|c}
\text{row}&P&Q&C\\\hline
29&B(3/5,1/5)&B(11/10,1/5)&(2P-2Q)/(-2P+3Q)\\
33&B(7/10,1/5)&B(6/5,1/5)&(2P-2Q)/(-3P+4Q)\\
36&B(3/5,3/5)&B(11/10,3/5)&(-4P+7Q)/(2P-2Q)\\
40&B(9/10,7/5)&B(7/5,7/5)&(-104P+192Q)/(65P-96Q)
\end{array}
\]
Their fractional-linear determinants are $2,2,-6,-2496$. Positivity of
the original denominator excludes poles.
For verification, define
\[
 J_k(u,v)=\int_0^1t^v(1-t^2)^u(1+t)^{-k}\,dt,
 \quad P_0=\tfrac12B((v+1)/2,u),\quad
 Q_0=\tfrac12B((v+2)/2,u).
\]
Rationalization and integration of
$[t^{v+1}(1-t)(1-t^2)^{u-1}]'$ give
\[
 J_1=P_0-Q_0,\qquad
 J_2=\frac{u+v}{u-1}P_0-\frac{2u+v}{u-1}Q_0.
\]
For rows 29 and 33 use $x=(1-t)/(1+t)$, obtaining
$D=2^{b-d+1}J_2(a,b)$, $N=2^{b-d}J_1(a,b)$.
Row 36 directly gives $D=J_1(b,a)$, $N=J_2(b,a)$.
For row 40 the substitution gives
$D=2^{b-d+1}J_3(a,b)$, $N=2^{b-d}J_2(a,b)$.
A certificate avoiding any analytic-continuation issue in the last case
uses $h=t^{4/5}(1-t^2)^{2/5}$ and
\[
 R_2=\frac{5t(1-t)(4t+5)}{6(1+t)},\qquad
 R_3=-\frac{5t(1-t)(16t^2+33t+11)}{96(1+t)^2}.
\]
Direct differentiation yields
$h/(1+t)^2=h(-13/2+12t)+(hR_2)'$ and
$h/(1+t)^3=h(65/32-3t)+(hR_3)'$.
The primitives vanish at both endpoints and establish the listed beta
formulas.

For rational arguments, put
$\epsilon_k(r,s)=\{kr\}+\{ks\}-\{k(r+s)\}$.
The beta-quotient theorem of Wolfart--W\"ustholz
\cite[Satz 4]{dbkz-ww} says that unequal profiles, with $k$ coprime
to a common denominator and nonintegral arguments and sums, force the
beta quotient to be transcendental. Integer shifts only multiply a
beta value by a nonzero algebraic number. The following single
profile discrepancies therefore suffice:
\[
\begin{array}{c|c|c|c|c}
\text{row}&(r,s)_P&(r,s)_Q&k&(\epsilon_k(P),\epsilon_k(Q))\\\hline
29&(3/5,1/5)&(1/10,1/5)&3&(1,0)\\
33&(7/10,1/5)&(1/5,1/5)&3&(0,1)\\
36&(3/5,3/5)&(1/10,3/5)&1&(1,0)\\
40&(9/10,2/5)&(2/5,2/5)&1&(1,0)
\end{array}
\]
This proves all four assertions through the nondegenerate
fractional-linear transformations.

\subsection{Eight simple genus-four Jacobians: exact certificates}
\label{dbkz-sec:simple}
The eight remaining rows arise from the normalizations
\[
\begin{gathered}
 X_{r,s,t,c}:\quad y^5=x^r(1-x)^s(1+cx)^t,\\
 (r,s,t)\in\{(1,1,2),(1,3,3),(2,4,2),(1,2,1)\},
 \qquad c\in\{1/5,2/5\}.
\end{gathered}
\]
All four branch exponents, including infinity, are nonzero modulo five;
the curves have genus four. Good reductions at $31$ for $c=1/5$ and
at $11$ for $c=2/5$ have the following exact point counts:
\[
\begin{array}{c|c|r|rrrr}
c&(r,s,t)&p&N_1&N_2&N_3&N_4\\\hline
1/5&\text{all four}&31&19&949&29434&923449\\
2/5&(1,1,2),(1,3,3)&11&9&69&1164&14369\\
2/5&(2,4,2),(1,2,1)&11&9&129&1374&14369
\end{array}
\]
These integers can be obtained without period approximations. In
$\mathbb F_{p^k}$, there is one point above each of the three finite
branch points and one at infinity. Every other fiber contains five
points precisely when the nonzero radicand is a fifth power. Thus
\begin{equation}\label{dbkz-eq:counts}
 N_k=4+5\#\{x\in\mathbb F_{p^k}:x\notin\{0,1,-1/c\},\
 [x^r(1-x)^s(1+cx)^t]^{(p^k-1)/5}=1\}.
\end{equation}
Suitable irreducible presentations, for degrees $2,3,4$ respectively,
are
\[
\begin{array}{c|ccc}
p&f_2(T)&f_3(T)&f_4(T)\\\hline
11&T^2+1&T^3+4T^2+1&T^4+4T^3+1\\
31&T^2+1&T^3+3T^2+1&T^4+T^3+1
\end{array}
\]
Enumeration using these fields and \eqref{dbkz-eq:counts} completely
specifies the finite count certificate.

Newton's identities and reciprocity give, in table order,
\begin{align*}
P_1(X)={}&X^8-13X^7+78X^6-401X^5+2195X^4\\
 &{}-12431X^3+74958X^2-387283X+923521,\\
P_2(X)={}&X^8-3X^7-22X^6+19X^5+335X^4\\
 &{}+209X^3-2662X^2-3993X+14641,\\
P_3(X)={}&X^8-3X^7+8X^6-X^5-85X^4\\
 &{}-11X^3+968X^2-3993X+14641.
\end{align*}
Writing $P_i(X)=X^4F_i(X+p/X)$ gives
\begin{align*}
F_1(U)&=U^4-13U^3-46U^2+808U-719,\\
F_2(U)&=U^4-3U^3-66U^2+118U+1061,\\
F_3(U)&=U^4-3U^3-36U^2+98U-19.
\end{align*}
Each reduces modulo two to the irreducible polynomial $U^4+U^3+1$.
The Weil bounds make its four roots real. None equals $\pm2\sqrt p$,
as it has degree four. For a root $u$, the corresponding Frobenius
root $\pi$ is nonreal and solves $X^2-uX+p=0$; hence
$[\mathbb Q(\pi):\mathbb Q]=8$. Each $P_i$ is irreducible.

To certify absolute simplicity, let
\[
 R_i(Z)=\operatorname{Res}_X(P_i(X),P_i(ZX)),\qquad
 Q_i(Z)=R_i(Z)/(Z-1)^8.
\]
The roots of $R_i$ are ratios of Frobenius roots. There is no
nontrivial root-of-unity ratio: $Q_i$ has degree $56$ and is divisible
by none of $\Phi_n$ with $2\le n\le6272$, $\varphi(n)\le56$.
This finite check is exhaustive, since
$\varphi(n)\ge\sqrt{n/2}$. The latter bound follows by multiplying
$p^{2-a}/(p-1)^2$ over $p^a\Vert n$: each odd-prime factor is at most
one and the factor at two is at most two.

For a short reproducible certificate, all the relevant gcds are already
one after reduction modulo $10007$. The following complete SageMath
fragment performs the certificate directly from the three displayed
polynomials; it needs no external table or data file.
\begin{verbatim}
ZX.<X> = PolynomialRing(ZZ)
polys = [
 X^8-13*X^7+78*X^6-401*X^5+2195*X^4
   -12431*X^3+74958*X^2-387283*X+923521,
 X^8-3*X^7-22*X^6+19*X^5+335*X^4
   +209*X^3-2662*X^2-3993*X+14641,
 X^8-3*X^7+8*X^6-X^5-85*X^4
   -11*X^3+968*X^2-3993*X+14641]
FZ.<z> = PolynomialRing(GF(10007))
RX.<x> = PolynomialRing(FZ)
orders = [n for n in range(2,6273) if euler_phi(n)<=56]
for P in polys:
    p = RX(P.list())
    r = FZ(p.resultant(p(z*x)))
    q, rem = r.quo_rem((z-1)^8)
    assert rem == 0 and q.degree() == 56 and q(1) != 0
    for n in orders:
        assert q.gcd(FZ(cyclotomic_polynomial(n).list())) == 1
\end{verbatim}
The same assertions were checked with exact integer resultants and
rational cyclotomic division. This modular version is sufficient:
a monic integral cyclotomic divisor over $\mathbb Q$ would remain a
divisor modulo $10007$.

Thus every positive power of $\pi$ retains eight distinct conjugates.
The absolute-simplicity criterion
\cite[Proposition 3 and Lemma 4]{dbkz-hz} applies. A nontrivial
geometric idempotent in the characteristic-zero endomorphism algebra
would specialize to one in the geometric special fiber, so these
characteristic-zero Jacobians are absolutely simple as well.

For the eight targets, write
$g=x^a(1-x)^b(1+cx)^{-d}$ and consider $g\,dx$ and
$g\,dx/(1+cx)$ on the corresponding cover. Their branch signatures
are
\[
\begin{array}{c|c|c}
(a,b,d)&(r,s,t)&\text{rows}\quad(c=1/5,2/5)\\\hline
(1/5,1/5,3/5)&(1,1,2)&31,32\\
(1/5,2/5,4/5)&(1,2,1)&34,35\\
(1/5,3/5,2/5)&(1,3,3)&37,38\\
(2/5,4/5,3/5)&(2,4,2)&41,42
\end{array}
\]
Deck symmetry kills residues. Exactness of a linear combination would
require a rational $S$ solving
\[
 S'+\left(\frac ax-\frac b{1-x}-\frac{dc}{1+cx}\right)S
       =A+\frac B{1+cx}.
\]
As before, nonintegral residues force $S$ to be polynomial and divisible
by $x(1-x)$. A nonzero polynomial of degree $k\ge2$ has a nonzero
leading term of degree $k-1$, since $k+a+b-d>0$, contrary to the
bounded right-hand side. Thus the classes are independent. The closed
cycle obtained by subtracting a deck translate of the positive interval
has nonzero denominator period. Simplicity and the fixed-cycle period
theorem prove all eight ratios transcendental. Together with the
previous subsections this completes every row of Table~\ref{dbkz-tab:44}.

\subsection{A congruence family and a rational-weight classifier}
\label{dbkz-sec:classifier}
The finite certificates yield an infinite family and an exact decision
procedure for normalized rational weights. Let
$a,b,d\in\frac15\mathbb Z$, $a,b>-1$, and $c>0$ algebraic. Define
\[
 g=x^a(1-x)^b(1+cx)^{-d},\quad D=\int_0^1g\,dx,\quad
 M=\frac{\int_0^1xg\,dx}{D},\quad
 \mathcal LS=S'+\left(\frac ax-\frac b{1-x}-\frac{dc}{1+cx}\right)S.
\]
Assume that the ordered residues
\begin{equation}\label{dbkz-eq:signature}
 (5a,5b,-5d,5(d-a-b))\pmod5
\end{equation}
are a permutation of $(h,h,h,2h)$ for some $h\in\{1,2,3,4\}$.
Assume also that a number field containing $c$ has a finite place with
one of the following residue conditions:
\begin{equation}\label{dbkz-eq:congruence}
 \operatorname{char}k(v)=31,\ \bar c=25\in\mathbb F_{31};
 \qquad\hbox{or}\qquad
 \operatorname{char}k(v)=11,\ \bar c=7\in\mathbb F_{11}.
\end{equation}
For rational $c=m/n$ in lowest terms these are respectively
$31\nmid n$, $m\equiv25n\pmod{31}$, or
$11\nmid n$, $m\equiv7n\pmod{11}$.

\begin{theorem}\label{dbkz-thm:classifier}
Under \eqref{dbkz-eq:signature}--\eqref{dbkz-eq:congruence}, both
$C(a,b,c,d)$ and $M$ are transcendental. If
$R\in\overline{\mathbb Q}(x)$ has poles only in
$\{0,1,-1/c,\infty\}$ and $gR$ is integrable on $(0,1)$, there are
unique $u,v\in\overline{\mathbb Q}$ and
$S\in\overline{\mathbb Q}(x)$ such that
\begin{equation}\label{dbkz-eq:classifier}
 R=u+vx+\mathcal LS.
\end{equation}
Moreover
\[
 \frac{\int_0^1gR\,dx}{D}=u+vM,
\]
which is algebraic if and only if $v=0$. All coefficients and the
rational primitive can be found by a finite linear system over an
explicit coefficient number field.
For $H(c)={}_2F_1(d,a+1;a+b+2;-c)$, the three numbers
$1,H(c),H'(c)$ are linearly independent over
$\overline{\mathbb Q}$.
\end{theorem}
\begin{proof}
The curve with function field $\overline{\mathbb Q}(x,g)$ is a cyclic
fivefold cover with four fully ramified branch points. Integer shifts
of exponents and common multiplication of the residues by a unit
modulo five do not change its geometric function field, up to a twist.
Twelve of the sixteen signatures in \eqref{dbkz-eq:signature} reduce
to the four triples of Section~\ref{dbkz-sec:simple}. For the four with
the exceptional exponent at infinity, $x\mapsto-1/(cx)$ swaps
$0\leftrightarrow\infty$ and $1\leftrightarrow-1/c$, moving that
exponent to zero; they then reduce to $(1,3,3,3)$ by a scalar modulo
five. Condition \eqref{dbkz-eq:congruence} makes the geometric special
fiber one of the certified absolutely simple curves. The four branch
sections remain distinct and the characteristic is not five, so
specialization proves that the characteristic-zero Jacobian $A$ is
simple of dimension four.

Every $gR\,dx$ with the stated pole support is of the second kind by
deck symmetry. Its nontrivial deck eigenspace in de Rham cohomology has
dimension two: rational cohomology has no invariant part, since the
quotient has genus zero, and its characteristic polynomial is
$\Phi_5^2$. The forms $g\,dx$ and $xg\,dx$ are independent. Indeed,
a rational solution of $\mathcal LS=A_0+A_1x$ must be polynomial and
vanish at all three finite branch points. If nonzero it has degree
$k\ge3$, while its leading image has degree $k-1\ge2$ and coefficient
proportional to $k+a+b-d\ne0$. Nonvanishing follows from the
nonintegral residue at infinity. This contradicts the right-hand side.
The same argument with right-hand side $A_0+A_1/(1+cx)$ proves
independence of $g\,dx$ and $g\,dx/(1+cx)$.

The first independent pair forms a basis of the eigenspace, proving
existence and uniqueness of $u,v$ in \eqref{dbkz-eq:classifier}.
An exact difference has a rational primitive $gS$ by character
averaging. The primitive $S$ is unique, since
$\mathcal LS=0$ would make $gS$ constant, contradicting the
nontrivial fifth-degree function-field extension unless $S=0$.

The closed cycle $\gamma=\sigma-\tau\sigma$, where $\sigma$ is the
positive lift of $[0,1]$, has nonzero denominator period
$(1-\zeta_5)D$. The simple fixed-cycle theorem makes $C$ and $M$
transcendental. Integrating \eqref{dbkz-eq:classifier} around this
cycle proves the integral formula, with no separate assumption on
endpoint values of the rational primitive.

For the last assertion put $B=B(a+1,b+1)>0$, so $D=BH(c)$.
The signature condition implies $a+b\notin\mathbb Z$.
The value $B$, up to a nonzero algebraic deck factor, is a second-kind
period of the genus-two three-point cover with function field
$\overline{\mathbb Q}(x,x^a(1-x)^b)$. Its Jacobian has no simple
factor isogenous to the simple fourfold $A$. On their product the
subgroup theorem separates the two isogeny components. Consequently
$B,D,N$ are linearly independent over $\overline{\mathbb Q}$, where
$N=\int_0^1g/(1+cx)\,dx$. Finally
\[
 N=D+\frac cd D',\qquad C=1+\frac cd\frac{H'}H
\]
transfers that independence to $1,H,H'$. The nonintegrality of $d$
ensures $d\ne0$. This conclusion is linear, not algebraic,
independence.
\end{proof}

To implement the finite algorithm, let $m_j$ denote the pole order of
$R$ at each finite branch point, and put
$M_\infty=\max(1,\deg_\infty R)$.
Every possible $S$ in \eqref{dbkz-eq:classifier} has pole order at most
$\max(m_j-1,0)$ there and polynomial part of degree at most
$M_\infty+1$. These bounds follow from the nonintegral local residues
and the nonzero leading coefficient at infinity. Substitute the
corresponding partial-fraction ansatz and solve for its coefficients
and $u,v$. Existence follows from the cohomology argument, rather than
from an experimental fit.

The theorem applies, for example, to $(a,b,d)=(1/5,1/5,3/5)$ and
\[
 c=1/8,\quad c=7+11n,\quad c=25+31n\quad(n\ge0),\quad c=\sqrt5.
\]
For the last value use the place $(11,\sqrt5-7)$.
The classifier also detects genuine algebraic exceptions among
nonconstant weights:
\[
 \mathcal L[x(1-x)(1+cx)]
 =(1+a)+\{(a+2)(c-1)-b-dc\}x-c(a+b+3-d)x^2.
\]
Its normalized integral is zero. Thus the theorem does not claim
transcendence for every nonconstant rational weight.

\clearpage

\part{Third-kind periods and uniform families}
\section{Third-kind periods and twenty-two values in the Log1 output}
\label{log1-section}

This section concerns the slice with fourth parameter zero,
\begin{equation}\label{log1-definition}
 F_h(c)=C(h,h,c,0)
 =\frac{\displaystyle\int_0^1[x(1-x)]^h(1+cx)^{-1}\,dx}
 {B(h+1,h+1)}
 ={}_2F_1(1,h+1;2h+2;-c),\qquad h,c>-1.
\end{equation}
All real fractional powers are positive. This slice differs from
$C(h,h,c,h)$ considered elsewhere in the paper. The results below are
applications of existing period and torsion theorems, supplemented by
explicit algebraic and local calculations. They are not assertions of
historical priority for the corresponding special-function values.

\subsection{An exact catalogue of the public entries}

Table~\ref{log1-table} lists the two initial third-kind targets and the
twenty subsequent targets in the public DBKZ \texttt{Log1} output
\cite{log1-dbkz}. Each row is specified exactly by $(h,c)$; decimal
approximations are included only to identify the source entry. Source
lines refer to the text snapshot examined in this work. All twenty-two
rows are proved transcendental below. Only rows 3 and 4 are asserted here
to be algebraically independent of one another. None is a duplicate of a
fourth-parameter-nonzero entry from the earlier \texttt{Log2} catalogue.
The absence of an identification in a computer-generated list does not
establish that the transcendence of an entry was previously unknown.

\begin{table}[htbp]
\centering\small
\caption{Twenty-two catalogued values $F_h(c)=C(h,h,c,0)$ in Log1.}
\label{log1-table}
\begin{tabular}{rccrl}
\toprule
No. & $h$ & $c$ & Approximate value & Source line\\
\midrule
1 & $1/3$ & $1/3$ & $0.86196583755412212\ldots$ &124\\
2 & $2/3$ & $1/3$ & $0.86121914488406966\ldots$ &194\\
3 & $1/3$ & $1$ & $0.68815073461701577\ldots$ &133\\
4 & $2/3$ & $1$ & $0.68472526272182017\ldots$ &203\\
5 & $1/4$ & $1/4$ & $0.89204558868041487\ldots$ &232\\
6 & $1/4$ & $1/2$ & $0.80934781900111164\ldots$ &238\\
7 & $1/4$ & $1$ & $0.68921689098230447\ldots$ &244\\
8 & $3/4$ & $1/4$ & $0.89134153873904941\ldots$ &338\\
9 & $3/4$ & $1/2$ & $0.80724556697156396\ldots$ &344\\
10 & $3/4$ & $1$ & $0.68403133820765647\ldots$ &350\\
11 & $1/5$ & $1/5$ & $0.91131086128452211\ldots$ &416\\
12 & $2/5$ & $1/5$ & $0.91107655008701731\ldots$ &469\\
13 & $3/5$ & $1/5$ & $0.91088694358067085\ldots$ &575\\
14 & $4/5$ & $1/5$ & $0.91073036736539899\ldots$ &628\\
15 & $1/7$ & $1/7$ & $0.93459900758455854\ldots$ &955\\
16 & $2/7$ & $1/7$ & $0.93449760261342181\ldots$ &979\\
17 & $3/7$ & $1/7$ & $0.93441123426490072\ldots$ &1003\\
18 & $4/7$ & $1/7$ & $0.93433678973658805\ldots$ &1045\\
19 & $5/7$ & $1/7$ & $0.93427195982569432\ldots$ &1069\\
20 & $6/7$ & $1/7$ & $0.93421499535486150\ldots$ &1093\\
21 & $1/4$ & $1/8$ & $0.94210870681369684\ldots$ &1176\\
22 & $3/4$ & $1/8$ & $0.94190133226732776\ldots$ &1336\\
\bottomrule
\end{tabular}
\end{table}

\subsection{The period-theoretic input}

We use the following consequence of the Huber--W\"ustholz theorem on
linear relations among 1-periods \cite[Theorem~9.7]{dbkz-hw}.
If $G$ is a semi-abelian variety over $\overline{\mathbb Q}$ and
$v\in H_1(G(\mathbb C),\mathbb Q)$, there is a connected semi-abelian
subgroup $B\subseteq G$, understood up to isogeny, such that $v$ comes
from $H_1(B,\mathbb Q)$ and the annihilator of $v$ in
$H^1_{\mathrm{dR}}(G)$ is the kernel of restriction to $B$.
This is a statement about the full algebraic de Rham realization,
including second-kind classes; it is not restricted to invariant
differentials on the semi-abelian variety.

\begin{lemma}\label{log1-eigencriterion}
Let $0\to T\to G\to A\to0$ have a finite-order automorphism $s$.
Assume $A$ has no proper nonzero $s$-stable abelian subvariety,
$X^*(T)\otimes\mathbb Q$ is an irreducible $s$-representation, and the
extension does not split up to isogeny. If $v$ has nonzero projection to
$H_1(A,\mathbb Q)$, evaluation at $v$ is injective on every eigenspace of
$s^*$ in $H^1_{\mathrm{dR}}(G)$.
\end{lemma}
\begin{proof}
If an eigenclass $\omega$ annihilates $v$, the subgroup $B$ just
described annihilates its restriction. The connected sum
$H=\sum_j s^jB$ also annihilates $\omega$. Its nonzero $s$-stable image
in $A$ is all of $A$. Irreducibility of the torus representation gives
$(H\cap T)^0=0$ or $T$. The former case makes $H$ abelian and $H\to A$
an isogeny, splitting the pullback extension, contrary to the hypothesis.
The latter makes $H=G$, whence $\omega=0$. In this last step restriction
to the factors of the surjection $\prod_j s^j B\to G$ detects first
de Rham cohomology.
\end{proof}

For a curve $X$ punctured at a reduced set $S$, the generalized Jacobian
has torus character group $\mathbb Z[S]^0$, and its extension class is
the map
\begin{equation}\label{log1-classmap}
 \mathbb Z[S]^0\longrightarrow\operatorname{Jac}(X)(\overline{\mathbb Q}),
 \qquad D\longmapsto[D]
\end{equation}
\cite[Lemma~4.24]{dbkz-hw}. A non-torsion value of this map obstructs
splitting up to isogeny. When the map is injective after tensoring with
$\mathbb Q$, $G$ has no nonzero algebraic character, up to isogeny.
If its abelian quotient is elliptic, no proper connected subgroup can
then surject onto that quotient: its quotient in $G$ would otherwise
be a positive-dimensional torus and provide such a character.

\subsection{Prime denominators, including the exceptional fifth-order case}

\begin{theorem}\label{log1-prime}
Let $h=m/p>-1$ be in lowest terms, where $p\ge5$ is prime. Then
$F_h(c)$ is transcendental for every real algebraic $c>-1$, $c\ne0$.
At $c=0$ it equals $1$.
\end{theorem}
\begin{proof}
First take $c>0$. On the smooth projective curve
\[
 X:\ w^p=x(1-x),\qquad \rho(x,w)=(x,\zeta_p w),
\]
the points above $0,1,\infty$ are totally ramified, so
$g(X)=(p-1)/2$. The quotient by $\rho$ is the projective line.
Consequently $H_1(J,\mathbb Q)$, for $J=\operatorname{Jac}(X)$, is the
irreducible $(p-1)$-dimensional $\Phi_p$ representation. In particular,
$J$ has no proper nonzero $\rho$-stable abelian subvariety, and
$1-\rho$ is an isogeny. This reasoning does not require absolute
simplicity of $J$.

Let $S$ be the $p$ points above $x=-1/c$, and set
\[
 \alpha=w^m\,dx,\qquad \beta=\frac{w^m\,dx}{1+cx},\qquad
 G=J(X\setminus S).
\]
At $0,1$, both differentials have order $m+p-1\ge0$. At infinity,
their residues vanish because that point is fixed by $\rho$ and their
eigenvalue is $\zeta_p^m\ne1$. Thus $\alpha$ is of the second kind,
whereas $\beta$ has nonzero residues $w(P)^m/c$ at every $P\in S$.
Both represent classes in $H^1_{\mathrm{dR}}(G)$; the meromorphic
description of this cohomology permits second-kind poles away from $S$.
The rational character group of its torus is $\mathbb Q[S]^0$, again
the irreducible $\Phi_p$ representation.

Let $\sigma$ be the positive-real-branch path from $x=0$ to $1$ and
$v=\sigma-\rho\sigma$. Its endpoints are fixed by $\rho$, so it is a
closed cycle, and
\[
 \int_v\alpha=(1-\zeta_p^m)B(h+1,h+1),\qquad
 F_h(c)=\frac{\int_v\beta}{\int_v\alpha}.
\]
Its image in $H_1(J,\mathbb Q)$ is nonzero: the holomorphic form
$dx/w^{p-1}$ has period $(1-\zeta_p)$ times a positive integral on
$\sigma$. If $F_h(c)=a\in\overline{\mathbb Q}$, the nonzero eigenclass
$\beta-a\alpha$ annihilates $v$. The subgroup argument in
Lemma~\ref{log1-eigencriterion} forces either a contradiction or an
abelian subgroup $H\subset G$ isogenous onto $J$, containing $v$ in
rational homology.

Use infinity as the Abel--Jacobi base point and write
$\iota(P)=[P-\infty]$. For $P\in S$,
$[P-\rho P]=(1-\rho)\iota(P)$. Hence non-torsion of $\iota(P)$ rules
out the abelian-subgroup alternative via \eqref{log1-classmap}.
For $p\ge7$, Grant--Shaulis \cite[Theorem~1.1]{log1-grant-shaulis}
classify the entire geometric torsion packet on this curve: its points
have $x=0,1,1/2,\infty$. None has $x=-1/c<0$.
For $p=5$, Boxall--Grant \cite[Theorem~0.2]{log1-boxall-grant} add the
points with $x=(1\pm\sqrt5)/2$. The sole remaining positive parameter
is $c=\varphi=(1+\sqrt5)/2$.

For this exception put $t=-w$, $s=(1-\sqrt5)/2$, and
\[
 q(t)=-\sqrt5\,t^2+\sqrt5\,t+s,
 \qquad q(t)^2-q(t)-t^5=-(t-1)^5.
\]
On $t^5=x^2-x$, define, with indices modulo $5$,
\[
 f_j=x-q(\zeta_5^{-j}t),\qquad P_j=(\zeta_5^j,s).
\]
The identity
$(x-q(t))(x+q(t)-1)=(t-1)^5$ shows that
$\operatorname{div}(f_j)=5(P_j-\infty)$: the second factor is
$-\sqrt5$ at $P_0$, and $x$ has a pole of order $5$ at infinity.
Thus $\chi=f_2/f_1$, after multiplication by a nonzero constant,
defines an algebraic character of $G$.

We compute its period without numerical winding tests. On $\sigma$,
write $u=[x(1-x)]^{1/5}$, so
\[
 f_j=x-s+\sqrt5(u^2\zeta_5^{-2j}+u\zeta_5^{-j}).
\]
Every $f_j$ starts and ends on the positive real axis. Let $n_j$ be
its continuous argument change divided by $2\pi$. Since $f_0>0$ and
$\operatorname{Im} f_1<0$ in the interior, $n_0=n_1=0$. Moreover,
\[
 \operatorname{Im}f_2
 =\sqrt5\,u\sin(\pi/5)(\varphi u-1).
\]
There are two interior real-axis crossings, at
$x_-<1/2<x_+$ with $x(1-x)=\varphi^{-5}$. At these crossings
$\operatorname{Re}f_2=x-k$, where $k=\sqrt5-2$ and
$k(1-k)=2\varphi^{-5}$. Therefore $x_-<k<x_+$: the first crossing
passes from below to above the negative real axis; the second crosses
the positive axis. It follows that $n_2=-1$. Since
$\rho^*f_j=f_{j-1}$, the endpoint logarithms cancel and
\begin{equation}\label{log1-winding}
 \int_v d\log\chi=2\pi i(n_2-2n_1+n_0)=-2\pi i\ne0.
\end{equation}
Every character restricts trivially to an abelian variety. Thus an
abelian $H$ containing $v$ would force the left side of
\eqref{log1-winding} to vanish. This excludes the last alternative.

Finally substitution $x\mapsto1-x$ gives, for all $h,c>-1$,
\begin{equation}\label{log1-symmetry}
 F_h(c)=\frac1{1+c}F_h\!\left(-\frac c{1+c}\right).
\end{equation}
For $-1<c<0$ the new parameter is positive algebraic, which completes
the proof.
\end{proof}

An inverse consequence also follows. For an exponent in
Theorem~\ref{log1-prime} and algebraic $r\in(0,1)$, the unique positive
solution of $F_h(c)=r$ is transcendental. Indeed, $F_h$ decreases
strictly from $1$ to $0$ on $[0,\infty)$, since
\[
 F'_h(c)=-\frac1{B(h+1,h+1)}
 \int_0^1\frac{x^{h+1}(1-x)^h}{(1+cx)^2}\,dx<0,
\]
and dominated convergence gives its limit at infinity. An algebraic
solution would contradict the theorem.

\subsection{Quarter exponents and all rational parameters}

\begin{theorem}\label{log1-quarter}
If $h>-1$ and $4h$ is odd, then $F_h(c)$ is transcendental for every
rational $c>-1$, $c\ne0$.
\end{theorem}
\begin{proof}
For $c>0$ use the elliptic curve $w^4=x(1-x)$ with
$\rho(x,w)=(x,iw)$. The coordinates
$U=x/w^2$, $V=x/w^3$ give $V^2=U^3+U$ and
$\rho(U,V)=(-U,iV)$. The four points above $x=-1/c$ are
$P,iP,-P,-iP$, where
\[
 U^2=-\frac1{c+1},\qquad V^2=\frac c{c+1}U,
 \qquad V^4=-\frac{c^2}{(c+1)^3}.
\]
We first prove uniformly that $P$ is non-torsion. Exact duplication gives
\begin{equation}\label{log1-quarter-double}
 2P=(AU,BV),\qquad A=-\frac{(c+2)^2}{4c},\qquad
 B=\frac{(c+2)(c^2-4c-4)}{8c^2}.
\end{equation}
Write $c=m/n$ in lowest positive terms. An odd prime $\ell$ dividing
$(m+2n)(m^2-4mn-4n^2)$ divides none of $m,n,m+n$, as reduction modulo
each shows. Over $\mathbb Q_\ell$, the polynomial
\[
 Z^4+\frac{m^2n}{(m+n)^3}
\]
has unit discriminant. Hence every completion of $\mathbb Q(V)$ above
$\ell$ is unramified, and $U,V$ are units. The elliptic model has good
reduction, since its discriminant is $-64$. If $r=v_\ell(B)>0$, then
$2P$ reduces to a nonzero point of order two, while duplication yields
\[
 v_\ell(U(4P))=-2r,\qquad v_\ell(V(4P))=-3r.
\]
Thus $4P\ne O$ lies in the formal subgroup. This subgroup is torsion-free
over an unramified extension of $\mathbb Q_\ell$ for odd $\ell$:
the formal logarithm has expansion
$z+\sum_{j\ge2}a_jz^j/j$, with integral $a_j$, and for $v(z)\ge1$
every higher term has valuation greater than $v(z)$. Its kernel is zero
and its target is an additive characteristic-zero field. This proves
non-torsion whenever such an odd prime exists.

If both numerator factors have no odd prime divisor, $m$ must be even;
write $m=2r$ with $n$ odd. The absence of odd divisors in $r+n$ forces
$r$ odd and $r+n=2^s$. Also
$v_2(r^2-2rn-n^2)=1$, so $r^2-2rn-n^2=\pm2$. Substitution yields
$r^2=2^{2s-1}\pm1$. For $s=1$ this gives $r=n=1$. For $s\ge2$,
one sign is impossible modulo $8$; factoring the other as
$(r-1)(r+1)=2^{2s-1}$ gives only $r=3$, $s=2$, $n=1$.
The exceptional parameters are therefore $c=2,6$. They admit the
following exact triplication certificates:
\[
\begin{array}{c|c|c|c}
c&U(3P)/U&V(3P)/V&\ell\\\hline
2&1/9&-11/27&11\\
6&169/121&-1443/1331&13
\end{array}
\]
At the indicated prime, $3P$ reduces to a nonzero point of order two
and $6P\ne O$ lies in the torsion-free formal subgroup. Thus $P$ is
non-torsion for every positive rational $c$.

There is a small representation-theoretic point in applying the period
theorem. The torus of $G=J(E\setminus\{P,iP,-P,-iP\})$ has rank three,
with rational character types $\Phi_2\oplus\Phi_4$. The connected image
$T_2$ of $1+\rho^2$ on $G$ lies in the torus, because $\rho^2=[-1]$
on $E$; it is one-dimensional of type $\Phi_2$. The quotient
$\overline G=G/T_2$ has a rank-two torus of type $\Phi_4$. Its
classifying map sends the anti-pair divisors
$[P]-[-P]$ and $[iP]-[-iP]$ to $2P$ and $2iP$.
This map is injective over $\mathbb Q$: any nonzero
$a+b\rho\in\mathbb Q(i)$ is an isogeny, since its product with
$a-b\rho$ is $[a^2+b^2]$, and cannot send $P$ to torsion.
Hence $\overline G$ has no nonzero character, and no proper connected
subgroup of it surjects onto $E$.

Put $k=4h$ and use $\alpha=w^kdx$, $\beta=w^kdx/(1+cx)$.
At $0,1$ their order is $k+3\ge0$. The two points at infinity are
fixed by $\rho^2$, whereas the eigenvalue of these forms under
$\rho^2$ is $-1$, so both residues vanish there. Their cohomology
classes have eigenvalue $i^k\in\{i,-i\}$ and descend uniquely to
$H^1_{\mathrm{dR}}(\overline G)$, because the removed torus has only
the eigenvalue $-1$. For $v=\sigma-\rho\sigma$ on the positive branch,
\[
 \int_v\alpha=(1-i^k)B(h+1,h+1)\ne0.
\]
The class of $\alpha$ comes from $E$, so the projected cycle is
nonzero. The nonzero residues of $\beta$ at the punctures prove that
the two classes are independent. The subgroup supplied by the
Huber--W\"ustholz theorem must surject onto $E$, hence must be all of
$\overline G$. Its period-evaluation map is therefore injective.
An algebraic ratio $F_h(c)$ would contradict this injectivity.
Equation~\eqref{log1-symmetry} handles the negative rational parameters.
\end{proof}

For the eight public quarter-denominator rows, the duplication factors
in \eqref{log1-quarter-double} give particularly short certificates:
\[
\begin{array}{c|c|c|c}
c&A&B&\ell\\\hline
1/8&-289/32&-4879/64&41\\
1/4&-81/16&-711/32&79\\
1/2&-25/8&-115/16&23\\
1&-9/4&-21/8&7
\end{array}
\]
The indicated valuation of $B$ is one. For each fixed positive rational
$c$, the four unnormalized integrals $D_{1/4},N_{1/4}(c),D_{3/4},N_{3/4}(c)$,
where $D_h=B(h+1,h+1)$ and $N_h(c)=D_hF_h(c)$, are linearly independent
over $\overline{\mathbb Q}$. They form a de Rham basis, with two
independent classes in each of the $i$ and $-i$ eigenspaces, and the
whole period-evaluation map just used is injective. This conclusion
does not assert algebraic independence of the two ratios.

\subsection{The two initial third-kind values at $c=1/3$}

\begin{theorem}\label{log1-initial-thirds}
The numbers $F_{1/3}(1/3)$ and $F_{2/3}(1/3)$ are transcendental.
The corresponding four unnormalized numerator and denominator
integrals are linearly independent over $\overline{\mathbb Q}$.
\end{theorem}
\begin{proof}
The curve $w^3=x(1-x)$ becomes $V^2=U^3+16$ under
$U=-4w$, $V=8x-4$. Write $\tau(w)=\zeta_3w$. The three punctures
above $x=-3$ are $P,\tau P,\tau^2P$, with
\[
 P=(U_0,-28),\quad U_0^3=768,\qquad
 2P=\left(\frac{10}{49}U_0,-\frac{1628}{343}\right).
\]
At every place above $7$, the field $\mathbb Q(U_0)$ is unramified,
because $X^3-768$ has discriminant prime to $7$; the elliptic model
also has good reduction. The coordinates of $2P$ have valuations
$-2,-3$. The same formal-logarithm argument as above proves that $P$
is non-torsion.

The generalized Jacobian has rank-two torus with character group
$\mathbb Q[\{P,\tau P,\tau^2P\}]^0$. Its classifying map is injective:
for rational $a,b$, the divisor
$a(\tau P-P)+b(\tau^2P-P)$ maps to
\[
 \bigl((-a-2b)+(a-b)\tau\bigr)P.
\]
Since $1+\tau+\tau^2=0$, this is a nonzero rational CM endomorphism
unless $a=b=0$, and cannot send the non-torsion $P$ to torsion.
Thus the generalized Jacobian has no nonzero character. The cycle
$v=\sigma-\tau\sigma$ has nonzero elliptic projection, as its integral
against $w^kdx$ is $(1-\zeta_3^k)B(1+k/3,1+k/3)$ for $k=1,2$.
Its period-evaluation map is injective by the subgroup theorem.
For each $k$, the classes $w^kdx$ and $w^kdx/(1+x/3)$ are independent:
the first has zero residues, the second has residues $3w(Q)^k$ at the
punctures. Infinity has zero residue by its nontrivial $\tau$-character.
The two pairs lie in distinct eigenspaces and form a basis of the
four-dimensional de Rham cohomology. The two ratio statements and the
four-period linear independence follow.
\end{proof}

\subsection{Explicit evaluations and independence at $c=1$}

\begin{theorem}\label{log1-explicit-thirds}
Put $r=\sqrt[3]{2}$, $\Omega=B(1/3,1/3)$, and
$\Psi=B(2/3,2/3)$. Then
\begin{align}
 L_1:=F_{1/3}(1)&=5-\frac{10}{3r^2}\Psi,\label{log1-L1}\\
 L_2:=F_{2/3}(1)&=\frac{21}{2}-\frac7{3r}\Omega.\label{log1-L2}
\end{align}
The pair $L_1,L_2$ is algebraically independent over
$\overline{\mathbb Q}$, and
\begin{equation}\label{log1-field}
 \overline{\mathbb Q}(L_1,L_2)
 =\overline{\mathbb Q}(\pi,\Gamma(1/3)^3).
\end{equation}
\end{theorem}
\begin{proof}
For $w=[x(1-x)]^{1/3}$ substitute $y=(1-x)/w$. This decreases from
$\infty$ to $0$ as $x$ increases from $0$ to $1$, and direct
differentiation gives
\[
 y^3+4=\frac{(1+x)^2}{x},\qquad
 \frac{dx}{(1+x)w}=-\frac{3\,dy}{y^3+4},\qquad
 \left(\frac1{1+x}-\frac12\right)\frac{dx}{w^2}
 =-\frac32\frac{y\,dy}{y^3+4}.
\]
Using $\int_0^\infty dt/(1+t^3)=\int_0^\infty t\,dt/(1+t^3)
=2\pi/(3\sqrt3)$ and scaling $y=r^2t$ yields
\[
 T_2:=\int_0^1\frac{dx}{(1+x)w}=\frac{\pi\sqrt3}{3r},\qquad
 T_1:=\int_0^1\frac{dx}{(1+x)w^2}
 =\frac\Omega2+\frac{\pi\sqrt3}{3r^2}.
\]
The beta recurrence gives $D_{1/3}=\Omega/10$ and $D_{2/3}=\Psi/7$.
Polynomial division
$x(1-x)/(1+x)=-x+2-2/(1+x)$ and symmetry about $x=1/2$ give
\[
 N_{1/3}(1)=\tfrac32\Omega-2T_1,\qquad
 N_{2/3}(1)=\tfrac32\Psi-2T_2.
\]
Finally $\Omega\Psi=2\pi\sqrt3$ proves
\eqref{log1-L1}--\eqref{log1-L2}. Writing $g=\Gamma(1/3)$, reflection
also gives
\[
 \Omega=\frac{\sqrt3\,g^3}{2\pi},\qquad
 \Psi=\frac{4\pi^2}{g^3}.
\]
The affine transformations are invertible over $\overline{\mathbb Q}$,
so they imply \eqref{log1-field}. Chudnovsky's established theorem
states that $\pi$ and $\Gamma(1/3)$ are algebraically independent
\cite[p.~344, equation~(9)]{log1-chudnovsky}. Passing from $g$ to $g^3$
does not change transcendence degree. This proves the assertion.
\end{proof}

For example, the exact reconstruction formula
\[
 \pi=\frac9{70\sqrt3}(5-L_1)\left(\frac{21}{2}-L_2\right)
\]
is consistent with independence, since it is not a relation with
algebraic coefficients between $L_1,L_2$ alone. Every nonconstant
polynomial in this pair with algebraic coefficients is transcendental.
The explicit evaluations can also be approached through classical
hypergeometric transformations; their derivation here is not a claim
of a historically new special-function identity.

These evaluations extend to all third-integer exponents:
\begin{equation}\label{log1-all-thirds}
 F_h(1)\ \hbox{and}\ F_h(-1/2)\ \hbox{are transcendental}
 \qquad(h\in\tfrac13\mathbb Z,\ h>-1).
\end{equation}
Indeed write $h=h_0+j$ with $j\ge0$ and
$h_0\in\{-2/3,-1/3,0\}$. Put $P_j(x)=[x(1-x)]^j$ and
$Q_j=(P_j-(-2)^j)/(1+x)\in\mathbb Q[x]$. Under the beta probability
measure with exponent $h_0$, all polynomial moments are rational:
$\mathbb E_{h_0}(x^k)=(h_0+1)_k/(2h_0+2)_k$. Therefore
\[
 F_h(1)=\frac{q_j+(-2)^jF_{h_0}(1)}{d_j},\qquad
 q_j=\mathbb E_{h_0}(Q_j)\in\mathbb Q,\quad
 d_j=\frac{(h_0+1)_j^2}{(2h_0+2)_{2j}}>0.
\]
The three base values are
\[
 F_{-2/3}(1)=\frac12+\frac\Psi{6r^2},\qquad
 F_{-1/3}(1)=\frac\Omega{6r},\qquad F_0(1)=\log2.
\]
The first two follow from the preceding substitutions; the last is
transcendental by Hermite--Lindemann. The transfer coefficient is
nonzero, and \eqref{log1-symmetry} gives $F_h(-1/2)=2F_h(1)$.
Moreover, any value in the class $h\equiv1/3\pmod1$ and any value in
$h\equiv2/3\pmod1$, at either of these two parameters, are
algebraically independent. This makes no independence assertion
involving the integer-exponent class and $\log2$.

\subsection{An algebraic complementary family}

Not all unidentified entries of this type are transcendental. For
rational $0<a<1$ and algebraic $c>0$, one has
\begin{equation}\label{log1-algebraic-complement}
 C(a,1-a,c,0)
 =\frac{2\{1+(1-a)c-(1+c)^{1-a}\}}{a(1-a)c^2}
 \in\overline{\mathbb Q}.
\end{equation}
The left side is ${}_2F_1(1,a+1;3;-c)$. Expanding the binomial on the
right at zero gives its hypergeometric series, and analytic
continuation on $c>0$ proves the identity. For example,
$C(1/3,2/3,1,0)=15-9\,2^{2/3}$. This example explains why selection
by proved geometric hypotheses, rather than a missing numerical
identification, is essential.

The proofs in this section use exact divisor identities, local
valuations, and established transcendence theorems. Numerical agreement
with the source entries serves only identification and arithmetic
checking. A general algorithmic treatment of linear relations among
1-periods is itself already part of the literature
\cite{log1-sertoz}; the contribution recorded here is the explicit
application and its stated parameter ranges, not the invention of that
general theory.

\clearpage

\part{Multiple-integral evaluations}
\section{Six double-integral and two denominator-six triple-integral values}
\label{extra-eight}
The following reductions preserve the boundary terms of the hypergeometric
summations and Barnes contours. They are applications of classical summation
and beta-transcendence theorems, with explicit identities supplied here.

\subsection{Definitions and scope}
We use
\[
B(r,s)=\frac{\Gamma(r)\Gamma(s)}{\Gamma(r+s)},\qquad
{}_3F_2\left(\begin{matrix}a,b,c\\d,e\end{matrix};1\right)
=\sum_{n\ge0}\frac{(a)_n(b)_n(c)_n}{(d)_n(e)_n n!}.
\]
The six double-integral constants are exactly the six entries under
``Denominator 7 with not yet identified constants'' on p.~14 of
\cite{extra8-dbkz}. Their common definition is
\begin{align*}
C_2(a_1,a_2,b_1,b_2)
&=\frac{\displaystyle\int_0^1\!\int_0^1
 \frac{x^{-a_1}(1-x)^{-a_2}y^{-b_1}(1-y)^{-b_2}}{1-xy}\,dx\,dy}
 {B(1-a_1,1-a_2)B(1-b_1,1-b_2)}\\
&={}_3F_2\left(\begin{matrix}1,1-a_1,1-b_1\\
2-a_1-a_2,2-b_1-b_2\end{matrix};1\right).
\end{align*}
For the triple-integral constants define
\begin{align*}
J(a,b,c,d,e)&=\iiint_{[0,1]^3}
\frac{x^b(1-x)^cy^e(1-y)^az^a(1-z)^c}{(1-z+xyz)^d}\,dx\,dy\,dz,\\
K(a,b,c,d,e)&=\frac{J(a,b,c,d+1,e)}{J(a,b,c,d,e)}.
\end{align*}
The two denominator-six entries treated here were left in apparently
unreduced hypergeometric form in the 2023 revision of \cite{extra8-zudilin}.
Their evaluation is established below rather than inferred from that
description. A statement about these dated sources does not establish
that no later or differently phrased proof exists.

\subsection{The arithmetic input}
Schneider's beta theorem \cite{extra8-schneider}, also stated in
\cite[Theorem~12]{extra8-waldschmidt} and in the introduction of the research
article \cite{extra8-dhillon-saha}, says that for positive rational numbers
\(r,s\), if \(r,s,r+s\notin\mathbb Z\), then \(B(r,s)\) is
transcendental. The condition does not require \(r+s<1\).
Thus every beta value used below is transcendental.
If \(\alpha,\beta\) are algebraic with \(\alpha\ne0\), then
\(\alpha B(r,s)+\beta\) is transcendental. More generally, an invertible
fractional-linear transformation over \(\Qbar\) preserves transcendence.
These elementary transfer steps do not assume algebraic independence of
arbitrary gamma values.

\subsection{A common residue lemma}
\label{extra8-sec:residue}
For later use set \(S=A+B+C+E-2\), and let
\[
I(A,B,C,E)=\frac1{2\pi i}\int_{\Re s=-1/12}
\frac{\Gamma(A+s)\Gamma(B+s)\Gamma(C+s)
      \Gamma(-2-s)\Gamma(E-s)}{\Gamma(S+s)}\,ds.
\]
On \(\Re A,\Re B,\Re C>1/12\) and \(\Re E>-1/12\), this integral
is holomorphic, with exponential decay on the vertical line locally
uniform in the parameters. Barnes' second lemma \cite[Eq.~(5.6)]{extra8-barnes}
and two residues give the meromorphic identity
\begin{align*}
I(A,B,C,E)&=G+R_{-2}+R_{-1},\\
G&=\frac{\Gamma(A-2)\Gamma(B-2)\Gamma(C-2)
\Gamma(A+E)\Gamma(B+E)\Gamma(C+E)}
{\Gamma(A+B+E-2)\Gamma(A+C+E-2)\Gamma(B+C+E-2)},\\
R_{-2}&=-\frac{\Gamma(A-2)\Gamma(B-2)\Gamma(C-2)\Gamma(E+2)}
{\Gamma(S-2)},\\
R_{-1}&=\frac{\Gamma(A-1)\Gamma(B-1)\Gamma(C-1)\Gamma(E+1)}
{\Gamma(S-1)}.
\end{align*}
For a rigorous starting domain take \(\Re A,\Re B,\Re C>9/4\),
\(\Re E>0\). The line \(\Re s=-9/4\) separates the pole families,
so Barnes' lemma applies there. Moving this line right to \(-1/12\)
crosses precisely the poles \(-2,-1\) of \(\Gamma(-2-s)\).
With upward orientation on both lines the right-line integral equals
the left-line integral plus these two residues. Horizontal segments
vanish by Stirling's estimate. Meromorphic continuation gives the
displayed identity throughout the connected domain stated above.
At the specializations below no displayed gamma factor has a pole.
In particular this argument does not apply a separating-contour formula
directly to a nonseparating vertical line.

\subsection{Reduction and transcendence of the first denominator-six triple-integral target}
\label{extra8-sec:triple6a}

Define
\[
 I_d=\iiint_{[0,1]^3}
 \frac{x^{1/2}(1-x)^{1/2}y^{1/6}(1-z)^{1/2}}
 {(1-z+xyz)^d}\,dx\,dy\,dz,
 \qquad d\in\left\{\frac13,\frac43\right\},
\]
and
\[
 C_6=\frac{I_{4/3}}{I_{1/3}}
 =K\left(0,\frac12,\frac12,\frac13,\frac16\right).
\]
This is the parameter convention of \cite{extra8-zudilin}:
\[
 K(a,b,c,d,e)=\frac{J_3(a,b,c,d+1,e;0)}{J_3(a,b,c,d,e;0)}.
\]
The integrals are finite and positive. Indeed,
$1-z+xyz\geq1-z$, and even for $d=4/3$ the remaining power
$(1-z)^{-5/6}$ is integrable.

\paragraph{Evaluation and result.}
Set
\[
 V=B\left(\frac13,\frac12\right)
 =\frac{\sqrt\pi\,\Gamma(1/3)}{\Gamma(5/6)}.
\]
We will prove
\begin{equation}
 \boxed{C_6=\frac{35(17-4V)}{2(32V-133)}.}
 \label{extra8-eq:triple6a-main}
\end{equation}
In particular $C_6$ is transcendental. For orientation only,
\[
 V=4.206546315976362783525057237\ldots,
 \qquad C_6=1.889898533450806850143673974\ldots.
\]

\paragraph{The original Mellin--Barnes formula.}
Specializing the formula of \cite{extra8-zudilin} gives
\[
 I_d=P_d\frac1{2\pi i}\int_{\Re s=-1/12}
 \frac{\Gamma(d+s)\Gamma(1+s)\Gamma(3/2+s)
 \Gamma(7/6+s)\Gamma(-s)\Gamma(3/2-d-s)}
 {\Gamma(13/6+s)\Gamma(3+s)}\,ds,
\]
where
\[
 P_d=\frac{\Gamma(3/2)}{\Gamma(d)\Gamma(5/2-d)}.
\]
The line is admissible for both values of $d$: it lies to the right
of all the left pole families and to the left of both right pole
families in this original representation. The identity
\[
 \frac{\Gamma(1+s)\Gamma(-s)}{\Gamma(3+s)}=\Gamma(-2-s)
\]
reduces the integrand to
\[
 F_d(s)=\frac{\Gamma(d+s)\Gamma(3/2+s)\Gamma(7/6+s)
 \Gamma(-2-s)\Gamma(3/2-d-s)}{\Gamma(13/6+s)}.
\]
The five numerator parameters now sum to the denominator parameter:
\[
 d+\frac32+\frac76-2+\left(\frac32-d\right)=\frac{13}6.
\]
This is precisely the parameter relation in Barnes' second lemma.
The two newly exposed poles at $s=-2,-1$ require residue corrections.

\paragraph{Corrected Barnes evaluation.}
Apply the residue identity of Section~\ref{extra8-sec:residue} with
\[
(A,B,C,E)=(d,3/2,7/6,3/2-d).
\]
Both $d=1/3$ and $d=4/3$ satisfy its hypotheses. The two residues are retained;
no separating-contour formula is used on a nonseparating line.

\paragraph{Exact gamma reduction.}
At the required parameters the three terms are
\[
 G_d=\frac{\Gamma(d-2)\Gamma(-1/2)\Gamma(-5/6)
 \Gamma(3/2)\Gamma(3-d)\Gamma(8/3-d)}
 {\Gamma(1)\Gamma(2/3)\Gamma(13/6-d)},
\]
\[
 R_{-2,d}=-\frac{\Gamma(d-2)\Gamma(-1/2)\Gamma(-5/6)
 \Gamma(7/2-d)}{\Gamma(1/6)},
\]
\[
 R_{-1,d}=\frac{\Gamma(d-1)\Gamma(1/2)\Gamma(1/6)
 \Gamma(5/2-d)}{\Gamma(7/6)}.
\]
For example, write
\[
 Q_d=\frac{\sqrt\pi\,\Gamma(1/6)\Gamma(3-d)\Gamma(8/3-d)}
 {\Gamma(5/2-d)\Gamma(2/3)\Gamma(13/6-d)}.
\]
Gamma shifts alone show $Q_{1/3}=(64/21)V$ and
$Q_{4/3}=(4/3)V$, and
\[
 \frac{P_dG_d}{\pi}
 =-\frac{Q_d}{2(d-1)(d-2)(-5/6)},
\]
\[
 \frac{P_dR_{-2,d}}{\pi}
 =\frac{5/2-d}{(d-1)(d-2)(-5/6)},\qquad
 \frac{P_dR_{-1,d}}{\pi}=\frac3{d-1}.
\]
Thus the complete exact coefficient table is
\[
\begin{array}{c|cc}
 &d=1/3&d=4/3\\ \hline
 P_dG_d/\pi&\frac{288}{175}V&-\frac{18}{5}V\\[2pt]
 P_dR_{-2,d}/\pi&-\frac{117}{50}&\frac{63}{10}\\[2pt]
 P_dR_{-1,d}/\pi&-\frac92&9
\end{array}
\]
Consequently
\[
 \frac{I_{1/3}}\pi=\frac{9(32V-133)}{175},\qquad
 \frac{I_{4/3}}\pi=\frac{9(17-4V)}{10},
\]
which proves \eqref{extra8-eq:triple6a-main}. Positivity also proves
$133/32<V<17/4$, so none of the displayed denominators vanish.

\paragraph{Transcendence and scope.}
Schneider's beta-value theorem \cite{extra8-schneider}, explicitly
recalled in the primary research article \cite{extra8-dhillon-saha}, states
that $B(a,b)$ is transcendental for positive rational nonintegral
$a,b$ with $a+b$ nonintegral. It applies to
$a=1/3$, $b=1/2$, $a+b=5/6$, proving that $V$ is transcendental.
Solving \eqref{extra8-eq:triple6a-main} for $V$ gives
\[
 V=\frac{266C_6+595}{64C_6+140}.
\]
The denominator is positive. If $C_6$ were algebraic then $V$ would
be algebraic, a contradiction. Hence $C_6$ is transcendental.

The 2023 version of \cite{extra8-zudilin}, page~3, lists this precise target
among expressions involving ${}_3F_2$ values that seem not reducible
to a simpler form. The present proof supplies such a reduction and
transfers a classical transcendence theorem to the integral ratio.
Historical priority of this specific evaluation has not been
established. This does not prove any proposed irrationality measure
or divisibility assertion associated with the original construction.

% Bibliography metadata for the enclosing document:
% schneider1941: T. Schneider, Zur Theorie der Abelschen Funktionen und
% Integrale, J. reine angew. Math. 183 (1941), 110--128,
% DOI 10.1515/crll.1941.183.110.
% dhillon-saha: S. Dhillon and E. Saha, On arithmetic nature of special
% values of the incomplete beta function, Acta Arith. 218 (2025),
% 273--284, DOI 10.4064/aa240610-6-10.

\subsection{The second denominator-six triple-integral quotient}
\label{extra8-sec:triple6b}
We continue to use the normalization in Zudilin's note \cite{extra8-zudilin},
so $J_3(a,b,c,d,e;0)=J(a,b,c,d,e)$ as defined above.
All powers in the real integrals use their positive real values.
In this subsection only, $C_2$ denotes the scalar below, rather than the
four-argument double-integral function $C_2(a_1,a_2,b_1,b_2)$.
\begin{theorem}
Let
\[
 C_2=2^{-1/3}\frac{\Gamma(1/3)^3}{\pi^2}.
\]
Then
\[
 \boxed{\displaystyle
 K\left(0,\frac12,\frac12,\frac16,\frac13\right)
 =\frac{1001C_2-1540}{275-175C_2}.}
\]
In particular this constant is transcendental. Its numerical value is
\[1.73295989536011838\ldots.\]
\end{theorem}

\subsubsection*{Barnes representation and exact cancellation}
Put $J_d=J_3(0,1/2,1/2,d,1/3;0)$, where
$d\in\{1/6,7/6\}$.  Zudilin's Barnes representation gives
\[
 J_d=P_d I_d,\qquad
 P_d=\frac{\Gamma(3/2)}{\Gamma(d)\Gamma(5/2-d)},
\]
\[
 I_d=\frac{1}{2\pi i}\int_{\Re s=-1/12}
 \frac{\Gamma(d+s)\Gamma(3/2+s)\Gamma(4/3+s)
       \Gamma(-2-s)\Gamma(3/2-d-s)}{\Gamma(7/3+s)}\,ds.
\]
Indeed, the original expression has factors
$\Gamma(1+s)\Gamma(-s)/\Gamma(3+s)$, and the gamma recurrence gives the exact identity
\[
 \frac{\Gamma(1+s)\Gamma(-s)}{\Gamma(3+s)}=\Gamma(-2-s).
\]
The original contour condition is
\[
 -\min\{1,3/2,d,4/3\}<-1/12<3/2-d,
\]
which holds for both values of $d$.  In particular the common vertical
contour is legitimate before making any analytic-continuation argument.

\subsubsection*{Barnes lemma and the contour correction}
Apply the common residue lemma in Section~\ref{extra8-sec:residue} with
\[
 A=d,\quad B=3/2,\quad C=4/3,\quad E=3/2-d,
 \qquad S=A+B+C+E-2=7/3.
\]
For both values of $d$, the hypotheses $A,B,C>1/12$ and $E>-1/12$
hold.  The analytic-continuation argument of that section begins with
a genuine separating line and keeps the two residues at $s=-2,-1$.
Thus no separating-contour formula is being applied without correction
to the nonseparating line $\Re s=-1/12$.  None of the gamma factors below
has a pole at either specialization.  Explicitly,
\[
 I_d=C_d+R_{-2,d}+R_{-1,d},
\]
where
\begin{align*}
 C_d&=
 \frac{\Gamma(d-2)\Gamma(-1/2)\Gamma(-2/3)\Gamma(3/2)
       \Gamma(3-d)\Gamma(17/6-d)}
 {\Gamma(5/6)\Gamma(7/3-d)},\\
 R_{-2,d}&=-\frac{\Gamma(d-2)\Gamma(-1/2)\Gamma(-2/3)
                         \Gamma(7/2-d)}{\Gamma(1/3)},\\
 R_{-1,d}&=\frac{\Gamma(d-1)\Gamma(1/2)\Gamma(1/3)
                        \Gamma(5/2-d)}{\Gamma(4/3)}.
\end{align*}
The signs follow from
$\operatorname{Res}_{s=-2}\Gamma(-2-s)=-1$ and
$\operatorname{Res}_{s=-1}\Gamma(-2-s)=1$.

\subsubsection*{Gamma recurrence simplification}
Write
\[
 H=\sqrt\pi\,\frac{\Gamma(2/3)}{\Gamma(1/6)}.
\]
Gamma recurrence alone gives
\[
 P_d(R_{-2,d}+R_{-1,d})
 =\frac{3\pi(4d-9)}{4(d-1)(d-2)}.
\]
For clarity, the individual exact terms are
\[
\begin{array}{c|ccc}
 d&P_dC_d&P_dR_{-2,d}&P_dR_{-1,d}\\ \hline
 1/6&(135/14)\pi H&-(126/55)\pi&-(9/5)\pi\\
 7/6&-54\pi H&(72/5)\pi&9\pi
\end{array}
\]
and hence
\[
 J_{1/6}=\pi\left(\frac{135}{14}H-\frac{45}{11}\right),
 \qquad
 J_{7/6}=\pi\left(-54H+\frac{117}{5}\right).
\]
Taking the quotient yields
\[
 K=\frac{2002-4620H}{825H-350}.
\]
Reflection and duplication give
\[
 \Gamma(1/3)\Gamma(2/3)=\frac{2\pi}{\sqrt3},\qquad
 \Gamma(1/6)\Gamma(2/3)=2^{2/3}\sqrt\pi\,\Gamma(1/3),
\]
so
\[
 H=\frac{B(2/3,2/3)}{3\,2^{2/3}}
   =\frac{2}{3C_2}.
\]
Substitution proves the asserted exact evaluation.

\subsubsection*{Transcendence and verification limits}
The numbers $2/3,2/3,4/3$ are rational and nonintegral.  Schneider's
beta-value theorem \cite{extra8-schneider} therefore makes $B(2/3,2/3)$
transcendental, even though its two arguments sum to more than one.
The identity $H=B(2/3,2/3)/(3\,2^{2/3})$ makes $H$ transcendental.
The fractional-linear map in $H$ has determinant
\[
 (-4620)(-350)-2002\cdot825=-34650\ne0,
\]
and inverse $H=(350K+2002)/(825K+4620)$.  Since $K>0$, the inverse
has nonzero denominator; hence $K$ is transcendental.

As an optional alternative, the transcendence of $C_2$ is explicitly
recalled in \cite[p.~2]{extra8-zudilin}.  The rational fractional-linear map
in $C_2$ has determinant
\[
 1001\cdot275-1540\cdot175=5775\ne0,
\]
and its inverse is
\[
 C_2=\frac{275K+1540}{175K+1001}.
\]
Since the original integrals are positive, $K>0$, and this inverse has
nonzero denominator.  Algebraicity of $K$ would contradict transcendence
of $C_2$.  The beta-value argument above does not require an
algebraic-independence theorem for gamma values.

\subsection{The first two double-integral candidates}
\label{extra8-sec:double12}

Put $s_j=\sin(j\pi/7)$ and $g_j=\Gamma(j/7)$.  For the first two
denominator-seven double-integral candidates, the hypergeometric forms
and evaluations are
\begin{align}
 F_1&={}_3F_2\left(\begin{matrix}1,13/7,11/7\\26/7,23/7\end{matrix};1\right)
 =\frac{418}{7}\frac{s_1}{s_2}B(2/7,4/7)-\frac{437}{3},
 \label{extra8-eq:double1-final}\\
 F_2&={}_3F_2\left(\begin{matrix}1,13/7,10/7\\25/7,22/7\end{matrix};1\right)
 =\frac{198}{7}\frac{s_2}{s_3}B(1/7,4/7)-180.
 \label{extra8-eq:double2-final}
\end{align}
These correspond respectively to
$C_2(-6/7,-6/7,-4/7,-5/7)$ and
$C_2(-6/7,-5/7,-3/7,-5/7)$ in the DBKZ list, with the
hypergeometric normalization used there.  Both are transcendental.
The proof below gives explicit finite telescoping certificates; it
does not rely on numerical recognition and makes no priority claim.

\subsubsection{Classical identities used}

Watson's summation, in the normalization of DLMF 16.4.6, is
\begin{equation}\label{extra8-eq:double12-watson}
 {}_3F_2\left(\begin{matrix}a,b,c\\(a+b+1)/2,2c\end{matrix};1\right)
 =\frac{\Gamma(1/2)\Gamma(c+1/2)\Gamma((a+b+1)/2)
              \Gamma(c+(1-a-b)/2)}
             {\Gamma((a+1)/2)\Gamma((b+1)/2)
              \Gamma(c+(1-a)/2)\Gamma(c+(1-b)/2)}.
\end{equation}
All applications below satisfy its convergence condition
$\operatorname{Re}(2c-a-b)>-1$.
We also use Euler's transformation of ${}_2F_1$ inside its beta integral.
For the positive real parameters appearing here this gives directly
\begin{equation}\label{extra8-eq:double12-euler}
 {}_3F_2\left(\begin{matrix}a,b,1\\e,f\end{matrix};1\right)
 =\frac{f-1}{e+f-a-b-1}
 {}_3F_2\left(\begin{matrix}e-a,e-b,1\\e,e+f-a-b\end{matrix};1\right).
\end{equation}
To see the prefactor explicitly, write the left side as
$(f-1)\int_0^1(1-t)^{f-2}{}_2F_1(a,b;e;t)\,dt$ and use
${}_2F_1(a,b;e;t)=(1-t)^{e-a-b}{}_2F_1(e-a,e-b;e;t)$.
This is a special case of a Thomae transformation.

\subsubsection{First candidate}

Set
\[
 W_1={}_3F_2\left(\begin{matrix}1,4/7,13/7\\9/7,26/7\end{matrix};1\right),
\]
\[
 h_n=\frac{(4/7)_n(13/7)_n}{(23/7)_n(26/7)_n},\qquad
 P_1(n)=\frac{(7n+16)(7n+19)(21n+23)}{48}.
\]
The summands of $F_1$ and $W_1$, respectively, are
\[
 h_n\frac{n+4/7}{4/7},\qquad
 h_n\frac{(n+9/7)(n+16/7)}{(9/7)(16/7)}.
\]
Direct multiplication, using
$h_{n+1}/h_n=(n+4/7)(n+13/7)/((n+23/7)(n+26/7))$, proves
\begin{equation}\label{extra8-eq:double1-certificate}
 P_1(n+1)h_{n+1}-P_1(n)h_n
 =h_n\left\{\frac{n+4/7}{4/7}
   -99\frac{(n+9/7)(n+16/7)}{(9/7)(16/7)}\right\}.
\end{equation}
Here $h_n=O(n^{-32/7})$, so $P_1(n)h_n\to0$, and
$P_1(0)=437/3$.  Summing proves
\begin{equation}\label{extra8-eq:double1-contiguous}
 F_1=99W_1-437/3.
\end{equation}
Watson's formula with $(a,b,c)=(1,4/7,13/7)$ gives
\[
 W_1=
 \frac{\sqrt\pi\,\Gamma(33/14)\Gamma(9/7)\Gamma(11/7)}
 {\Gamma(11/14)\Gamma(13/7)\Gamma(29/14)}
 =\frac{38}{63}\frac{g_2^2g_4g_5}{g_1g_6^2}
 =\frac{38}{63}\frac{s_1}{s_2}B(2/7,4/7).
\]
The middle equality uses only the gamma shift and duplication formulas;
the last uses $g_jg_{7-j}=\pi/s_j$.  Combining this with
\eqref{extra8-eq:double1-contiguous} proves \eqref{extra8-eq:double1-final}.

\subsubsection{Second candidate}

Taking $(a,b,e,f)=(13/7,10/7,22/7,25/7)$ in
\eqref{extra8-eq:double12-euler} gives
\[
 F_2=\frac{18}{17}V_2,\qquad
 V_2={}_3F_2\left(\begin{matrix}1,9/7,12/7\\22/7,24/7\end{matrix};1\right).
\]
Now set
\[
 W_2={}_3F_2\left(\begin{matrix}1,2/7,12/7\\8/7,24/7\end{matrix};1\right),
\]
\[
 h_n=\frac{(2/7)_n(12/7)_n}{(22/7)_n(24/7)_n},\qquad
 P_2(n)=\frac{(7n+15)(7n+17)(35n+32)}{48}.
\]
The exact difference certificate is
\begin{equation}\label{extra8-eq:double2-certificate}
 P_2(n+1)h_{n+1}-P_2(n)h_n
 =h_n\left\{\frac{n+2/7}{2/7}
  -\frac{275}{2}\frac{(n+8/7)(n+15/7)}{(8/7)(15/7)}\right\}.
\end{equation}
Again $h_n=O(n^{-32/7})$, so the upper boundary vanishes.
Since $P_2(0)=170$, this proves
\[
 V_2=\frac{275}{2}W_2-170,\qquad
 F_2=\frac{2475}{17}W_2-180.
\]
Watson's formula with $(a,b,c)=(1,2/7,12/7)$ evaluates the remaining
series:
\[
 W_2=
 \frac{\sqrt\pi\,\Gamma(31/14)\Gamma(8/7)\Gamma(11/7)}
 {\Gamma(9/14)\Gamma(12/7)\Gamma(29/14)}
 =\frac{34}{175}\frac{g_1g_3g_4^2}{g_2g_5^2}
 =\frac{34}{175}\frac{s_2}{s_3}B(1/7,4/7).
\]
This proves \eqref{extra8-eq:double2-final}.

\subsubsection{Transcendence conclusion and scope of verification}

Schneider's beta-value transcendence theorem for rational arguments
\cite{extra8-schneider,extra8-waldschmidt} applies to
$B(2/7,4/7)$ and $B(1/7,4/7)$: their arguments and their sums are
nonintegral.  Both sine ratios in
\eqref{extra8-eq:double1-final}--\eqref{extra8-eq:double2-final} are nonzero algebraic
numbers.  A nonconstant affine transformation with algebraic
coefficients preserves transcendence.  Therefore both existing DBKZ
candidates $F_1$ and $F_2$ are transcendental.

\subsection{The third and fourth double-integral candidates}
\label{extra8-sec:double34}

The third and fourth denominator-seven entries in the list of
unidentified double-integral constants in~\cite[p.~14]{extra8-dbkz} are
\[
 F_3={}_3F_2\left(\begin{matrix}1,13/7,9/7\\25/7,17/7\end{matrix};1\right),
 \qquad
 F_4={}_3F_2\left(\begin{matrix}1,13/7,8/7\\24/7,16/7\end{matrix};1\right).
\]
The source has $17/7$, not $15/7$, as the last lower parameter of $F_3$.
The parameter tuple gives the lower parameter
\[2-(-2/7)-(-1/7)=17/7.\]
The excesses of these two series are $13/7$
and $12/7$, respectively.  Thus both series converge.

Put
\[
 s_j=\sin(j\pi/7),\qquad
 T=\frac{\pi\Gamma(3/7)}{\Gamma(1/7)\Gamma(2/7)}
   =\frac{s_2}{7}B\left(\frac37,\frac57\right).
\]
The beta identity follows from
$\Gamma(2/7)\Gamma(5/7)=\pi/s_2$ and
$\Gamma(8/7)=\Gamma(1/7)/7$.

\paragraph{Exact evaluations and transcendence.}
We have
\begin{align}
 F_3&=40-\frac{66s_2}{s_1s_3}\,T
     =40-\frac{66s_2^2}{7s_1s_3}
       B\left(\frac37,\frac57\right),\label{extra8-eq:double3}\\
 F_4&=\frac{323}{24}-\frac{425s_1}{8s_2^2}\,T
     =\frac{323}{24}-\frac{425s_1}{56s_2}
       B\left(\frac37,\frac57\right).\label{extra8-eq:double4}
\end{align}
Schneider's beta-value theorem~\cite{extra8-schneider}, in the form stated
in~\cite[Theorem~12]{extra8-waldschmidt}, implies that
$B(3/7,5/7)$ is transcendental: $3/7$, $5/7$, and their sum $8/7$
are all nonintegers.  The coefficients of this beta value in
\eqref{extra8-eq:double3} and~\eqref{extra8-eq:double4} are nonzero algebraic numbers.
Consequently both $F_3$ and $F_4$ are transcendental.
This conclusion uses the transcendence of a single beta value;
no conjecture about gamma values or their algebraic independence is needed.

\paragraph{Watson sums.}
Define the two convergent Watson series
\[
 W_3={}_3F_2\left(\begin{matrix}1,6/7,9/7\\10/7,18/7\end{matrix};1\right),
 \qquad
 W_4={}_3F_2\left(\begin{matrix}1,6/7,8/7\\10/7,16/7\end{matrix};1\right).
\]
Watson's formula~\cite[Eq.~16.4.6]{extra8-dlmf} yields
\begin{align}
 W_3&=\frac{\sqrt\pi\,\Gamma(25/14)\Gamma(10/7)\Gamma(6/7)}
 {\Gamma(13/14)\Gamma(9/7)\Gamma(19/14)}
 =\frac{33s_2}{10s_1s_3}\,T,\label{extra8-eq:double-w3}\\
 W_4&=\frac{\sqrt\pi\,\Gamma(23/14)\Gamma(10/7)\Gamma(5/7)}
 {\Gamma(13/14)\Gamma(8/7)\Gamma(17/14)}
 =\frac{9s_1}{s_2^2}\,T.\label{extra8-eq:double-w4}
\end{align}
The last equalities use only the gamma recurrence, reflection, and
duplication formulas.  The convergence condition of Watson's formula,
$\operatorname{Re}(2c-a-b)>-1$, gives $5/7>-1$ and $3/7>-1$ here.

\paragraph{Exact telescoping certificates.}
The remaining relations are
\begin{equation}
 F_3+20W_3=40,\qquad
 F_4+\frac{425}{72}W_4=\frac{323}{24}.
 \label{extra8-eq:double34-contiguous}
\end{equation}
They can be checked term by term without a generalized Watson table.
Write $\Delta g_n=g_{n+1}-g_n$.  First put
\[
 h_n=\frac{(1)_n(6/7)_n(9/7)_n}{(17/7)_n(25/7)_n n!},\qquad
 P_3(n)=-\frac{(7n+10)(7n+12)(7n+18)}{54}.
\]
If $f_{3,n},w_{3,n}$ are the defining summands of $F_3,W_3$, then
\[
 \frac{f_{3,n}}{h_n}=\frac{7n+6}{6},\qquad
 \frac{w_{3,n}}{h_n}=\frac{(7n+10)(7n+18)}{180},
\]
and direct rational-function cancellation gives
\[
 \Delta\bigl(h_nP_3(n)\bigr)=f_{3,n}+20w_{3,n}.
\]
The initial boundary is $h_0P_3(0)=-40$; since
$h_n=O(n^{-27/7})$ and $P_3(n)=O(n^3)$, the upper boundary is zero.
This proves the first identity in~\eqref{extra8-eq:double34-contiguous}.

For the second identity put
\[
 k_n=\frac{(1)_n(6/7)_n(8/7)_n}{(24/7)_n(16/7)_n n!},\qquad
 P_4(n)=-\frac{(7n+9)(7n+17)(21n+38)}{432}.
\]
Writing $f_{4,n},w_{4,n}$ for the summands of $F_4,W_4$ gives
\[
 \frac{f_{4,n}}{k_n}=\frac{7n+6}{6},\qquad
 \frac{w_{4,n}}{k_n}=\frac{(7n+10)(7n+17)}{170},
\]
and
\[
 \Delta\bigl(k_nP_4(n)\bigr)=f_{4,n}+\frac{425}{72}w_{4,n}.
\]
Here $k_0P_4(0)=-323/24$, and the upper boundary is zero because
$k_n=O(n^{-26/7})$ and $P_4(n)=O(n^3)$.
Combining the contiguous relations with the Watson sums proves the
exact evaluations.

The argument proves the qualitative transcendence assertions.  It does
not prove the quantitative irrationality bounds conjectured for the
original rational-approximation constructions.

\subsection{The fifth and sixth double-integral targets}

The last two entries in the denominator-seven list of unidentified
double-integral constants in \cite[p.~14]{extra8-dbkz} are exactly
\begin{align*}
 F_5&=C_2(-6/7,-3/7,-5/7,-3/7)
 ={}_3F_2\!\left(\begin{matrix}1,13/7,12/7\\23/7,22/7\end{matrix};1\right),\\
 F_6&=C_2(-5/7,-3/7,-4/7,-2/7)
 ={}_3F_2\!\left(\begin{matrix}1,12/7,11/7\\22/7,20/7\end{matrix};1\right).
\end{align*}
Let $s=\sin(3\pi/7)/\sin(\pi/7)$. We prove
\begin{equation}\label{extra8-eq:double56-evaluation}
 \boxed{\begin{aligned}
 F_5&=-152+\frac{48}{7}s\,B(1/7,2/7),\\
 F_6&=-\frac{195}{2}+\frac{195}{7}s\,B(5/7,6/7).
 \end{aligned}}
\end{equation}
In particular, both specified pre-existing constants are transcendental.

\subsubsection{An elementary transformation and exact certificates}

Euler's transformation of ${}_2F_1$, followed by beta integration, yields
\begin{equation}\label{extra8-eq:double56-thomae}
 {}_3F_2\!\left(\begin{matrix}1,a,b\\e,f\end{matrix};1\right)
 =\frac{f-1}{e+f-a-b-1}
 {}_3F_2\!\left(\begin{matrix}1,e-a,e-b\\e,e+f-a-b\end{matrix};1\right).
\end{equation}
For the present positive parameters both beta integrals converge; this
identity also follows from \cite[Eq.~16.4.11]{extra8-dlmf}. Explicitly, the integral
representation on the left is
\[
 (f-1)\int_0^1(1-t)^{f-2}{}_2F_1(a,b;e;t)\,dt,
\]
and Euler's factor $(1-t)^{e-a-b}$ gives the right side of
\eqref{extra8-eq:double56-thomae}. Thus
\[
 F_5=\frac{16}{13}G_5,\qquad
 G_5={}_3F_2\!\left(\begin{matrix}1,9/7,10/7\\22/7,20/7\end{matrix};1\right).
\]

For $F_5$, set
\[
 h_n=\frac{(2/7)_n(10/7)_n}{(22/7)_n(20/7)_n},\quad
 U(n)=\frac{7n+2}{2},\quad
 V(n)=\frac{(7n+8)(7n+15)}{120},
\]
and
\[
 Q(n)=\frac{(7n+13)(7n+15)(77n+76)}{120}.
\]
The following is an identity of rational functions of $n$:
\begin{equation}\label{extra8-eq:double5-certificate}
 Q(n+1)\frac{(n+2/7)(n+10/7)}{(n+22/7)(n+20/7)}-Q(n)
   =U(n)-99V(n).
\end{equation}
Multiplication by $h_n$ and summation telescope. Since
$Q(n)h_n=O(n^{-9/7})\longrightarrow0$ and $Q(0)=247/2$, this proves
\[
 G_5=99W_5-\frac{247}{2},\qquad
 W_5={}_3F_2\!\left(\begin{matrix}1,2/7,10/7\\8/7,20/7\end{matrix};1\right).
\]
Indeed $U(n)h_n$ and $V(n)h_n$ are respectively the summands of $G_5$
and $W_5$.

For $F_6$, instead put
\[
 h_n=\frac{(12/7)_n(11/7)_n}{(20/7)_n(22/7)_n},\qquad
 V(n)=\frac{7n+13}{13},\qquad
 Q(n)=\frac{(7n+13)(7n+15)}{2}.
\]
The exact certificate is
\begin{equation}\label{extra8-eq:double6-certificate}
 Q(n+1)\frac{(n+12/7)(n+11/7)}{(n+20/7)(n+22/7)}-Q(n)
    =1-\frac{65}{2}V(n).
\end{equation}
Here $Q(n)h_n=O(n^{-5/7})\longrightarrow0$ and $Q(0)=195/2$; hence
\[
 F_6=\frac{65}{2}W_6-\frac{195}{2},\qquad
 W_6={}_3F_2\!\left(\begin{matrix}1,12/7,11/7\\13/7,22/7\end{matrix};1\right).
\]
All series used in these two telescoping arguments converge absolutely.
Both displayed certificates can be verified simply by clearing their
quadratic denominators and expanding over $\mathbb Q$.

\subsubsection{Watson summation and beta reduction}

Watson's sum \cite[Eq.~16.4.6]{extra8-dlmf} applies to $W_5$ and $W_6$ with,
respectively,
\[
 (a,b,c)=(1,2/7,10/7),\qquad (a,b,c)=(1,12/7,11/7).
\]
The convergence quantities $c+(1-a-b)/2$ are $9/7$ and $5/7$, both
strictly positive. The resulting exact values are
\begin{align*}
 W_5&=\frac{\sqrt\pi\,\Gamma(8/7)\Gamma(27/14)\Gamma(9/7)}
                 {\Gamma(9/14)\Gamma(10/7)\Gamma(25/14)},\\
 W_6&=\frac{\sqrt\pi\,\Gamma(13/7)\Gamma(29/14)\Gamma(5/7)}
                 {\Gamma(19/14)\Gamma(11/7)\Gamma(17/14)}.
\end{align*}
Writing $g_j=\Gamma(j/7)$, gamma translation and duplication give
\[
 W_5=\frac{13}{231}\frac{g_1^2g_2g_6}{g_3^2g_4},\qquad
 W_6=\frac32\frac{g_1g_5g_6^2}{g_4^2g_3}.
\]
Reflection gives
\[
 \frac{g_1g_6}{g_3g_4}
  =\frac{\sin(3\pi/7)}{\sin(\pi/7)}=s.
\]
Using $B(1/7,2/7)=g_1g_2/g_3$ and
$B(5/7,6/7)=7g_5g_6/(4g_4)$, we therefore obtain
\[
 W_5=\frac{13}{231}s\,B(1/7,2/7),\qquad
 W_6=\frac67s\,B(5/7,6/7).
\]
Substituting these into the two proved contiguous relations establishes
\eqref{extra8-eq:double56-evaluation}.

\subsubsection{Transcendence}

Schneider's theorem \cite[p.~114]{extra8-schneider} states that $B(r,t)$ is
transcendental when $r,t\in\mathbb Q$ and none of $r,t,r+t$ is an integer.
Its hypotheses hold for $(r,t)=(1/7,2/7)$ and $(5/7,6/7)$.
The coefficient $s$ is a nonzero algebraic number, because its sine values
are formed from roots of unity. Both expressions in
\eqref{extra8-eq:double56-evaluation} are therefore nonzero algebraic multiples
of transcendental numbers plus rational numbers, and both are
transcendental.

The displayed finite telescoping identities and subsequent gamma algebra
are exact. The numerical values
\[
 F_5=1.6723484102511823299843852299\ldots,\qquad
 F_6=1.6644935592570095489190072527\ldots
\]
are included only for identifying the constants, not as evidence for the
identities. This argument uses classical transformations and transcendence
theorems with explicit certificates; it does not establish a historical
priority claim for these evaluations or their transcendence consequences.



\section{The five denominator-seven triple-integral quotients}
\label{extra-seven}
These five entries are separate from the six double integrals above.
The first evaluation is already explicit in Zudilin's note; the remaining
formulas are retained as proof records without a historical-priority claim.

\subsection{Definitions and the five conclusions}
Set
\begin{align*}
J(a,b,c,d,e)&=\iiint_{[0,1]^3}
\frac{x^b(1-x)^cy^e(1-y)^az^a(1-z)^c}{(1-z+xyz)^d}
\,dx\,dy\,dz,\\
K(a,b,c,d,e)&=\frac{J(a,b,c,d+1,e)}{J(a,b,c,d,e)}.
\end{align*}
All integrals used below converge and are positive. Write
\(g_j=\Gamma(j/7)\), \(t=2\cos(2\pi/7)\), and
\[
U=\frac{2g_2^2g_5}{g_1g_4^2},\quad
H=\frac{g_1^2g_6}{g_2^2g_4},\quad
L=\frac{t^2}{1+t}\frac{g_1g_4}{g_2g_3},\quad
W=\frac{g_4^2g_3}{g_2g_1^2},\quad
T=\frac{g_3^2g_4}{g_5g_6^2}.
\]
The exact evaluations are
\begin{align*}
K_1&=K(1/7,0,2/7,3/7,4/7)=\frac{945U-1512}{832-525U},\\
K_2&=K(1/7,0,2/7,5/7,3/7)=\frac{15H-48}{70-23H},\\
K_3&=K(1/7,0,3/7,4/7,5/7)=\frac{525L-567}{462-425L},\\
K_4&=K(1/7,0,4/7,2/7,5/7)=\frac{9702W-363}{101-2744W},\\
K_5&=K(2/7,0,3/7,4/7,5/7)=\frac{520-120T}{88T-373}.
\end{align*}
Their decimal values, for identification only, are
\begin{align*}
K_1&=2.56248224568\ldots,& K_2&=3.59939534133\ldots,\\
K_3&=2.52360910522\ldots,& K_4&=1.99773527607\ldots,\\
K_5&=2.72543856463\ldots.
\end{align*}
These are the five entries under ``Denominator 7'' in \cite{extra8-dbkz}.
The first formula is printed in \cite{extra8-zudilin}. The same source states that
similar gamma evaluations exist for the next three, without displaying them,
and lists the fifth among the expressions not apparently reducible beyond
\(\F32\). Accordingly, the explicit formulas for the middle three should
not be described as the first discovery of their gamma evaluability.

\subsection{Arithmetic step: a finite certificate for each value}
For positive rational arguments with none of \(a,b,a+b,c,d,c+d\) integral,
define
\[
h_{a,b}(u)=\{ua\}+\{ub\}-\{u(a+b)\}.
\]
The beta-value theorem of Wolfart--W\"ustholz, Satz 4, together with the
classical-relation criterion on pp.~6--8 of \cite{dbkz-ww}, implies
\[
h_{a,b}(u)\ne h_{c,d}(u)\text{ for some }(u,N)=1
\quad\Longrightarrow\quad
\frac{B(a,b)}{B(c,d)}\notin\Qbar,
\]
where \(N\) is a common denominator. The theorem includes both first-kind
and second-kind beta periods; the sums of the arguments need not both be
less than one. This is a theorem about ratios of \emph{two beta values}, not
an unconditional classification of arbitrary products of gamma values.

Reflection, duplication and gamma shifts give
\begin{align*}
U&=\frac{2\sin(3\pi/7)}{\sin(2\pi/7)}
 \frac{B(2/7,3/7)}{B(1/7,4/7)},\\
H&=\frac{\sin(3\pi/7)}{\sin(\pi/7)}
 \frac{B(1/7,3/7)}{B(2/7,2/7)},\\
L&=\frac{t^2}{1+t}\frac{B(1/7,4/7)}{B(2/7,3/7)},\\
W&=\frac17\frac{B(4/7,4/7)}{B(1/7,2/7)},\qquad
T=\frac74\frac{B(3/7,3/7)}{B(5/7,6/7)}.
\end{align*}
All prefactors are nonzero algebraic numbers. The following table supplies
exact witnesses; its two middle columns are sums of fractional parts.
\begin{center}
\begin{tabular}{ccccc}
\toprule
Parameter & witness \(u\) & numerator beta & denominator beta & carries\\
\midrule
\(U\)&2&\(10/7\)&\(3/7\)&\(1,0\)\\
\(H\)&3&\(5/7\)&\(12/7\)&\(0,1\)\\
\(L\)&2&\(3/7\)&\(10/7\)&\(0,1\)\\
\(W\)&1&\(8/7\)&\(3/7\)&\(1,0\)\\
\(T\)&1&\(6/7\)&\(11/7\)&\(0,1\)\\
\bottomrule
\end{tabular}
\end{center}
Thus all five parameters are transcendental. In the order \(K_1,\ldots,K_5\),
the fractional-linear determinants are
\[
-7560,\quad-54,\quad1575,\quad-16170,\quad-1000.
\]
They are nonzero. An algebraic value of any \(K_i\) would therefore imply
an algebraic value of its parameter, a contradiction.

\begin{theorem}
All five specified denominator-seven DBKZ integral quotients are
transcendental.
\end{theorem}

\subsection{Common analytic reduction}
We use the Barnes representation in \cite{extra8-zudilin}:
\begin{align*}
J(a,b,c,d,e)
&=\frac{\Gamma(1+a)\Gamma(1+c)}{\Gamma(d)\Gamma(2+a+c-d)}
\frac1{2\pi i}\int_{\Re s=\sigma}
\frac{\Gamma(d+s)\Gamma(1+a+s)\Gamma(1+b+s)\Gamma(1+e+s)}
 {\Gamma(2+a+e+s)\Gamma(2+b+c+s)}\\[-2mm]
&\hspace{45mm}\cdot\Gamma(-s)\Gamma(1+c-d-s)\,ds,
\end{align*}
where \(-\min(a+1,b+1,d,e+1)<\sigma<c-d+1\).
The usual Mellin--Barnes separation also requires \(\sigma<0\); the
line \(\sigma=-1/2\) used for the fifth case satisfies all these bounds.
For the other cases the two-residue-family expression is
\[
J=\frac{\pi}{\sin\pi(d-c)}(A_dF_d-B_dG_d),
\]
with
\begin{align*}
A_d&=\frac{\Gamma(a+1)^2\Gamma(b+1)\Gamma(c+1)\Gamma(e+1)}
 {\Gamma(d-c)\Gamma(a+e+2)\Gamma(b+c+2)\Gamma(a+c-d+2)},\\
B_d&=\frac{\Gamma(a+1)\Gamma(c+1)^2
 \Gamma(b+c-d+2)\Gamma(e+c-d+2)}
 {\Gamma(c-d+2)\Gamma(d)\Gamma(a+e+c-d+3)\Gamma(b+2c-d+3)},\\
F_d&=\F43\left(\begin{matrix}a+1,b+1,d,e+1\\d-c,a+e+2,b+c+2\end{matrix};1\right),\\
G_d&=\F43\left(\begin{matrix}a+c-d+2,b+c-d+2,c+1,e+c-d+2\\
 c-d+2,a+e+c-d+3,b+2c-d+3\end{matrix};1\right).
\end{align*}
In the first case the printed formula in \cite{extra8-zudilin} uses
\(2^{2/7}\sqrt\pi\Gamma(9/14)/[\cos(3\pi/14)g_4^2]\).
Duplication and reflection identify this with \(U\), proving the first row.
The following sections give exact certificates for the four remaining rows.

\subsection{The second denominator-seven candidate}
\label{extra7-sec:target2}

Write $g_k=\Gamma(k/7)$ for $1\leq k\leq6$.  The evaluation proved
in this section is
\begin{equation}\label{extra7-eq:target2-main}
 K\left(\frac17,0,\frac27,\frac57,\frac37\right)
 =\frac{15H-48}{70-23H},\qquad
 H=\frac{g_1^2g_6}{g_2^2g_4}.
\end{equation}
Together with the beta-period criterion of Wolfart and W\"ustholz
\cite{dbkz-ww}, this proves that this $K$ is transcendental.  The argument
uses the hypergeometric representation in \cite{extra8-zudilin} and explicit
telescoping identities.  No claim of historical priority is made.

For clarity, the conventions for the integrals are
\begin{align*}
 J(a,b,c,d,e)&=\int_0^1\int_0^1\int_0^1
 \frac{x^b(1-x)^c y^e(1-y)^a z^a(1-z)^c}
 {(1-z+xyz)^d}\,dx\,dy\,dz,\\
 K(a,b,c,d,e)&=\frac{J(a,b,c,d+1,e)}{J(a,b,c,d,e)}.
\end{align*}
The particular integrals used below converge and are strictly positive.

\subsubsection{Two elementary balanced summations}

Let
\[
 t_n=\frac{(A)_n(B)_n}{(C)_n(D)_n},\qquad
 \delta=1-C,\quad p=A+\delta,\quad q=B+\delta,
 \quad R=\frac{\Gamma(C)\Gamma(D)}{\Gamma(A)\Gamma(B)}.
\]
The following identities are used only at positive rational parameters
with nonzero displayed denominators; they also hold for generic complex
parameters for which the series and expressions are defined.

First suppose $C+D=A+B+2$.  Then
\begin{equation}\label{extra7-eq:target2-sum1}
 {}_3F_2\left(\begin{matrix}1,A,B\\ C,D\end{matrix};1\right)
 =\frac{\delta(D-1)+R}{pq}.
\end{equation}
Indeed, with $Q(n)=(n+C-1)(n+D-1)$, direct multiplication gives
\[
 Q(n+1)t_{n+1}-Q(n)t_n=pq\,t_n.
\]
The gamma-ratio asymptotic gives $n^2t_n\longrightarrow R$.
Summing the difference from $n=0$ to $N$ and then taking the limit yields
$pq\sum_{n\geq0}t_n=R-(C-1)(D-1)$, proving
\eqref{extra7-eq:target2-sum1}.  In particular, the boundary at infinity must
not be discarded.

Next suppose $C+D=A+B+3$.  Put
\[
 E=p(p+1)q(q+1),\qquad
 k=\frac{pq+1+\delta}{E},\qquad
 \ell=\frac{(1+\delta)(D-2)}{E},
\]
\[
 P(n)=(n+C-1)(n+D-1)(kn+\ell).
\]
Another polynomial identity is
\[
 P(n+1)t_{n+1}-P(n)t_n=(n+1)t_n.
\]
Here $n^3t_n\longrightarrow R$, so the upper boundary is $kR$.
Since $(2)_n/n!=n+1$, telescoping proves
\begin{equation}\label{extra7-eq:target2-sum2}
 {}_3F_2\left(\begin{matrix}2,A,B\\ C,D\end{matrix};1\right)
 =\frac{\delta(1+\delta)(D-1)(D-2)
       +R(pq+1+\delta)}{p(p+1)q(q+1)}.
\end{equation}
Thus both summations have finite algebraic difference certificates.

\subsubsection{Reduction of the four hypergeometric series}

Specialize $a=1/7$, $b=0$, $c=2/7$, and $e=3/7$, and abbreviate
$J_d=J(a,b,c,d,e)$.  The Barnes-residue representation from
\cite{extra8-zudilin}, specialized to these parameters, is
\begin{equation}\label{extra7-eq:target2-barnes}
 J_d=\frac{\pi}{\sin\pi(d-2/7)}
       \bigl(A_dF_d-B_dG_d\bigr),
 \qquad d\in\{5/7,12/7\},
\end{equation}
where the prefactors are
\[
\begin{array}{c|cc}
 d&A_d&B_d\\ \hline
 5/7&\displaystyle\frac{7g_1^2}{660g_4g_5}
     &\displaystyle\frac{49g_2^2}{2340g_5g_6}\\[2mm]
 12/7&\displaystyle\frac{7g_1^2}{396g_4g_5}
      &\displaystyle\frac{7g_2^2}{60g_5g_6}
\end{array}
\]
and, after cancelling equal upper and lower parameters where possible,
\begin{align*}
 F_{5/7}&={}_4F_3\left(
  \begin{matrix}8/7,1,5/7,10/7\\3/7,18/7,16/7\end{matrix};1\right),\\
 G_{5/7}&={}_3F_2\left(
  \begin{matrix}2,9/7,12/7\\20/7,22/7\end{matrix};1\right),\\
 F_{12/7}&={}_3F_2\left(
  \begin{matrix}1,8/7,12/7\\16/7,18/7\end{matrix};1\right),\\
 G_{12/7}&={}_3F_2\left(
  \begin{matrix}1,5/7,9/7\\13/7,15/7\end{matrix};1\right).
\end{align*}
All four series converge absolutely.  Formulas
\eqref{extra7-eq:target2-sum1}--\eqref{extra7-eq:target2-sum2} immediately evaluate
the last three series.

To reduce $F_{5/7}$ as well, set
\[
 U={}_3F_2\left(
    \begin{matrix}1,5/7,8/7\\16/7,18/7\end{matrix};1\right),
 \qquad V=F_{12/7}.
\]
For the summands $t_n$ of $U$, the corresponding summands of $V$ are
$t_n(1+7n/5)$.  Also
$(10/7)_n/(3/7)_n=1+7n/3$, whence
\[
 F_{5/7}=-\frac23U+\frac53V.
\]
To eliminate $U$, use $A=5/7$, $B=8/7$, $C=16/7$, $D=18/7$ and
$Q(n)=(n+C-1)(n+D-1)$.  Here $C+D=A+B+3$, and
\[
 Q(n+1)t_{n+1}-Q(n)t_n
 =\left(AB-(C-1)(D-1)-n\right)t_n.
\]
This time $Q(n)t_n\longrightarrow0$.  Consequently
\[
 \sum_{n\geq0}nt_n=-\frac{59}{49}U+\frac{99}{49},\qquad
 V=-\frac{24}{35}U+\frac{99}{35}.
\]
It follows that
\[
 U=\frac{99-35V}{24},\qquad
 F_{5/7}=-\frac{11}{4}+\frac{95}{36}V.
\]
Using only $\Gamma(x+1)=x\Gamma(x)$ in the balanced summations gives
\begin{align*}
 F_{5/7}&=\frac{11(23g_1g_5-38g_2g_4)}{3g_1g_5},\\
 G_{5/7}&=-\frac{26(19g_1g_6-35g_2g_5)}{7g_2g_5},\\
 F_{12/7}&=\frac{33(5g_1g_5-8g_2g_4)}{5g_1g_5},\\
 G_{12/7}&=-\frac{8(g_1g_6-2g_2g_5)}{g_2g_5}.
\end{align*}

Substituting in \eqref{extra7-eq:target2-barnes}, and noting that
$\sin(10\pi/7)=-\sin(3\pi/7)$, proves the two exact integral values
\begin{align}
 J_{5/7}&=\frac{\pi}{\sin(3\pi/7)}
   \frac{7(23g_1^2g_6-70g_2^2g_4)}{180g_4g_5g_6},
   \label{extra7-eq:target2-j0}\\
 J_{12/7}&=\frac{\pi}{\sin(3\pi/7)}
   \frac{-7(5g_1^2g_6-16g_2^2g_4)}{60g_4g_5g_6}.
   \label{extra7-eq:target2-j1}
\end{align}
Their ratio proves \eqref{extra7-eq:target2-main}.  In particular the
denominator in that equation is nonzero, since $J_{5/7}>0$.

\subsubsection{The beta-period obstruction}

Euler's reflection formula yields
\[
 H=\frac{\sin(3\pi/7)}{\sin(\pi/7)}
   \frac{g_1g_3}{g_2^2}
  =\frac{\sin(3\pi/7)}{\sin(\pi/7)}
   \frac{B(1/7,3/7)}{B(2/7,2/7)}.
\]
The sine ratio is a nonzero algebraic number.  For rational beta
arguments the criterion of \cite{dbkz-ww} tests algebraic proportionality
using the profiles
\[
 h_{r,s}(u)=\{ur\}+\{us\}-\{u(r+s)\},
 \qquad u\in(\mathbb Z/7\mathbb Z)^\times.
\]
The two profiles differ already at $u=3$:
\[
 h_{1/7,3/7}(3)=\frac37+\frac27-\frac57=0,
 \qquad
 h_{2/7,2/7}(3)=\frac67+\frac67-\frac57=1.
\]
Thus the beta ratio, and hence $H$, is transcendental.  Finally,
the rational fractional linear transformation in
\eqref{extra7-eq:target2-main} has determinant
$15\cdot70-(-48)(-23)=-54\ne0$.  Such a transformation preserves
transcendence at every point where its denominator is nonzero.
This completes the proof of the transcendence of the second candidate.

\subsection{The third denominator-seven candidate}
\label{extra7-sec:target3}

Put
\begin{align*}
 J(d)&=J_3\left(\frac17,0,\frac37,d,\frac57;0\right)\\
 &=\int_{[0,1]^3}
 \frac{(1-x)^{3/7}y^{5/7}(1-y)^{1/7}z^{1/7}(1-z)^{3/7}}
 {(1-z+xyz)^d}\,dx\,dy\,dz,\\
 K_3&=\frac{J(11/7)}{J(4/7)}.
\end{align*}
These are the parameters of the third denominator-seven candidate; the
notation and the hypergeometric reduction of $J_3$ are those of
Zudilin~\cite{extra8-zudilin}.  Define
\[
 t=2\cos\frac{2\pi}{7},\qquad u=2\cos\frac{\pi}{7}=t^2+t-1,
 \qquad
 R=\frac{\Gamma(1/7)\Gamma(4/7)}{\Gamma(2/7)\Gamma(3/7)},
\]
\[
 L=\frac{t^2}{1+t}R=(1+t-t^2)R,
 \qquad D=\frac{u\Gamma(1/7)\Gamma(3/7)^2}{180}.
\]
The equalities involving $t$ follow from
$t^3+t^2-2t-1=0$.

\paragraph{Exact evaluation.}
The following identities hold:
\begin{equation}
 \label{extra7-eq:t3-evaluation}
 J(4/7)=D(462-425L),\qquad
 J(11/7)=D(525L-567),\qquad
 K_3=\frac{525L-567}{462-425L}.
\end{equation}
We give certificates for the evaluations, rather than relying on a
numerical identification.

All hypergeometric series below have argument $1$.  Set
\[
 C=\frac{\Gamma(20/7)\Gamma(17/7)}{\Gamma(4/7)\Gamma(12/7)},
 \qquad
 C'=\frac{\Gamma(20/7)\Gamma(17/7)}{\Gamma(11/7)\Gamma(12/7)}.
\]
The four required series are
\begin{align}
 F_A&={}_4F_3\left(\begin{matrix}8/7,1,4/7,12/7\\1/7,20/7,17/7\end{matrix};1\right)
       =\frac{8575}{36}C-\frac{1105}{3}, \label{extra7-eq:t3-fa}\\
 F_B&={}_3F_2\left(\begin{matrix}1,11/7,12/7\\20/7,17/7\end{matrix};1\right)
       =\frac{49}{2}C'-65, \label{extra7-eq:t3-fb}\\
 F_C&={}_3F_2\left(\begin{matrix}2,10/7,18/7\\26/7,23/7\end{matrix};1\right)
       =\frac{38000W_2-42408}{77}, \label{extra7-eq:t3-fc}\\
 F_E&={}_3F_2\left(\begin{matrix}1,10/7,11/7\\19/7,16/7\end{matrix};1\right)
       =9(W_1-1), \label{extra7-eq:t3-fe}
\end{align}
where Whipple's summation formula~\cite[Eq.~16.4.7]{extra7-dlmf}, followed by
the gamma recurrence, reflection, and duplication formulas, gives
\begin{align}
 W_2&={}_3F_2\left(\begin{matrix}2,3/7,4/7\\19/7,16/7\end{matrix};1\right)
      =\frac{\sin(3\pi/7)}{2\sin(\pi/7)}, \label{extra7-eq:t3-w2}\\
 W_1&={}_3F_2\left(\begin{matrix}1,3/7,4/7\\12/7,9/7\end{matrix};1\right)
      =\frac{5\sin^2(\pi/7)}{\sin(2\pi/7)\sin(3\pi/7)}.
      \label{extra7-eq:t3-w1}
\end{align}
In particular the use of Whipple's formula is inside its convergent
range: its parameter $c$ equals $2$ and $1$, respectively.

\paragraph{Four telescoping certificates.}
Write $\Delta g_n=g_{n+1}-g_n$, and use $(a)_n$ for the rising factorial.
For~\eqref{extra7-eq:t3-fa}, let
\[
 a_n=\frac{(8/7)_n(1)_n(4/7)_n(12/7)_n}
 {(1/7)_n(20/7)_n(17/7)_n n!},\qquad
 r_A(n)=\frac{(1225n/36+119/6)(n^2+23n/7+130/49)}{n+1/7}.
\]
Direct cancellation of rising factorials gives
$\Delta(a_n r_A(n))=a_n$.  The boundary values are
\[
 a_0r_A(0)=\frac{1105}{3},\qquad
 \lim_{n\to\infty}a_nr_A(n)=\frac{8575}{36}C.
\]
Thus the boundary at infinity must be retained, and summation proves
\eqref{extra7-eq:t3-fa}.

For~\eqref{extra7-eq:t3-fb}, let
\[
 b_n=\frac{(1)_n(11/7)_n(12/7)_n}{(20/7)_n(17/7)_n n!},
 \qquad r_B(n)=\frac{49}{2}n^2+\frac{161}{2}n+65.
\]
Then $\Delta(b_nr_B(n))=b_n$, with initial value $65$ and limit
$49C'/2$ at infinity.  This proves~\eqref{extra7-eq:t3-fb}.

For~\eqref{extra7-eq:t3-fc}, use the common term
\[
 h_n=\frac{(2)_n(3/7)_n(4/7)_n}{(26/7)_n(23/7)_n n!}
\]
and the rational certificate
\[
 r_C(n)=\frac{(7n+16)(7n+19)}{49(n+1)}
 \left(\frac{1953}{22}+\frac{30779}{264}n
       +\frac{3087}{88}n^2-\frac{343}{132}n^3\right).
\]
Let $c_n$ and $w_{2,n}$ be the usual summands of $F_C$ and $W_2$.
Then the exact rational identity
\[
 c_n-\frac{38000}{77}w_{2,n}=\Delta(h_nr_C(n))
\]
holds.  Its initial boundary is $h_0r_C(0)=42408/77$ and its boundary
at infinity is zero.  Summing proves~\eqref{extra7-eq:t3-fc}.

For~\eqref{extra7-eq:t3-fe}, take
\[
 \widetilde h_n=\frac{(1)_n(3/7)_n(4/7)_n}
 {(19/7)_n(16/7)_n n!},\qquad
 r_E(n)=\frac{(7n+9)(7n+12)}{12}.
\]
If $e_n,w_{1,n}$ are the summands of $F_E,W_1$, respectively, then
\[
 e_n-9w_{1,n}=\Delta(\widetilde h_nr_E(n)).
\]
The initial boundary is $9$ and the limit at infinity is zero.
This proves~\eqref{extra7-eq:t3-fe}.  All the displayed boundary limits follow
from $\Gamma(n+a)/\Gamma(n+b)\sim n^{a-b}$; all four series converge.

\paragraph{Substitution into the integral.}
For transparency, put $g_j=\Gamma(j/7)$ and $s_j=\sin(j\pi/7)$.
Zudilin's reduction~\cite{extra8-zudilin} specializes to
\[
 J(4/7)=\frac{\pi}{s_1}(A F_A-B F_C),\qquad
 J(11/7)=-\frac{\pi}{s_1}(A' F_B-B' F_E),
\]
where
\[
 A=\frac{g_8^2g_{10}g_{12}}{g_1g_{20}g_{17}},\quad
 B=\frac{g_8g_{10}^2g_{18}}{g_4g_{26}g_{23}},\quad
 A'=\frac{g_8g_{10}g_{12}}{g_{20}g_{17}},\quad
 B'=\frac{g_8g_{10}^2}{g_{19}g_{16}}.
\]
Substitute~\eqref{extra7-eq:t3-fa}--\eqref{extra7-eq:t3-w1}, reduce gamma arguments
by recurrence, reflection, and duplication, and use
\[
 s_2=u s_1,\qquad s_3=(1+t)s_1,\qquad
 s_1^2=\frac{2-t}{4},\qquad t^3+t^2-2t-1=0.
\]
These substitutions give exactly~\eqref{extra7-eq:t3-evaluation}; every
step is an identity in rational functions, gamma recurrences, or the
explicit cubic field defined above.

\paragraph{Transcendence.}
The common-sum beta representation is
\[
 R=\frac{B(1/7,4/7)}{B(2/7,3/7)}.
\]
For rational $a,b>0$ with $a+b<1$, write
$T_{a,b}=\{v\in(\mathbb Z/7\mathbb Z)^*:
\{va\}+\{vb\}<1\}$.  Here
\[
 T_{1/7,4/7}=\{1,2,4\},\qquad
 T_{2/7,3/7}=\{1,4,5\}.
\]
Already at $v=2$ their fractional-part sums are $3/7$ and $10/7$.
Wolfart--W\"ustholz's beta-value theorem and classical-relation
criterion~\cite[Satz~4 and Proposition~2]{dbkz-ww} therefore show that
$R$ is transcendental.  The factor $t^2/(1+t)$ is nonzero algebraic,
so $L$ is transcendental.  Finally the determinant of the fractional
linear transformation in~\eqref{extra7-eq:t3-evaluation} is
\[
 525\cdot462-(-567)(-425)=1575\ne0.
\]
Consequently $K_3$ is transcendental.  Positivity and finiteness of
$J(4/7)$ also ensure that the displayed denominator is nonzero.
This argument establishes transcendence; it does not establish the
quantitative irrationality-measure conjecture accompanying the candidate.

\subsection{The fourth denominator-seven target}

We consider the constant already listed in the DBKZ collection,
\[
 K_4=K\left(\frac17,0,\frac47,\frac27,\frac57\right)
       =\frac{J_1}{J_0},\qquad
 J_j=J_3\left(\frac17,0,\frac47,\frac27+j,\frac57;0\right)
 \quad(j=0,1).
\]
Write $g_r=\Gamma(r/7)$ for $1\leq r\leq6$, and put
\[
 W=\frac{g_4^2g_3}{g_2g_1^2},\qquad
 P=\frac{7\pi g_1^2}{2970\sin(2\pi/7)g_3g_6}>0.
\]
We will prove the exact evaluations
\begin{equation}\label{extra7-eq:target4-J}
 J_0=-P(101-2744W),\qquad
 J_1=33P(11-294W),
\end{equation}
and hence
\begin{equation}\label{extra7-eq:target4-K}
 \boxed{K_4=\frac{9702W-363}{101-2744W}.}
\end{equation}
The identification of this particular target with a gamma expression was
already announced, without its formula, in \cite[p.~3]{extra8-zudilin}. The purpose
of the calculation below is to give the explicit expression and an exact,
independently checkable certificate, rather than to claim that its gamma
evaluation was previously unknown.

\subsubsection{Four exact telescoping certificates}

For rational $A,B,C,D$ let
\[
 T_n=\frac{(A)_n(B)_n}{(C)_n(D)_n}.
\]
If polynomials $p,q$ satisfy
\begin{equation}\label{extra7-eq:target4-certificate}
 q(n+1)(n+A)(n+B)
   =(q(n)+p(n))(n+C)(n+D),
\end{equation}
then $p(n)T_n=q(n+1)T_{n+1}-q(n)T_n$. In each of the four cases below,
$m=C+D-A-B=\deg q$ is a positive integer, and the gamma-ratio asymptotic
gives
\[
 \lim_{N\longrightarrow\infty}q(N)T_N
   =\operatorname{lc}(q)\frac{\Gamma(C)\Gamma(D)}{\Gamma(A)\Gamma(B)}.
\]
Consequently the convergent sum is exactly
\begin{equation}\label{extra7-eq:target4-sum}
 \sum_{n=0}^{\infty}p(n)T_n
 =\operatorname{lc}(q)\frac{\Gamma(C)\Gamma(D)}{\Gamma(A)\Gamma(B)}-q(0).
\end{equation}
Here the four parameter tuples and polynomials are
\[
\begin{array}{c|rrrr|l}
 r&A&B&C&D&p_r(n)\\ \hline
 1&2/7&8/7&18/7&20/7&-(7n-2)(7n+5)/10\\
 2&11/7&17/7&27/7&29/7&(n+1)(n+2)/2\\
 3&8/7&9/7&18/7&20/7&(7n+5)/5\\
 4&10/7&11/7&20/7&22/7&n+1
\end{array}
\]
and
\begin{align*}
 q_1(n)&=\frac{(7n+11)(7n+13)(2107n^2+2765n+606)}{990},\\
 q_2(n)&=-\frac{(7n+20)(7n+22)(43n^2+167n+156)}{198},\\
 q_3(n)&=-\frac{(7n+11)(7n+13)(21n+22)}{50},\\
 q_4(n)&=-\frac{(3n+4)(7n+13)(7n+15)}{10}.
\end{align*}
For every row, identity \eqref{extra7-eq:target4-certificate} follows by polynomial
expansion over $\mathbb Q$. The leading coefficients and initial values are
\[
\begin{array}{c|r|r|r}
 r&m&\operatorname{lc}(q_r)&q_r(0)\\ \hline
 1&4&103243/990&1313/15\\
 2&4&-2107/198&-1040/3\\
 3&3&-1029/50&-1573/25\\
 4&3&-147/10&-78
\end{array}
\]
All four original summands are $O(n^{-2})$; thus there is no regularized or
formally assigned sum in this calculation.

Define the four ordinary hypergeometric values
\begin{align*}
 F_1&={}_4F_3\!\left(\begin{matrix}1,2/7,8/7,12/7\\-2/7,18/7,20/7\end{matrix};1\right),\\
 F_2&={}_3F_2\!\left(\begin{matrix}3,11/7,17/7\\27/7,29/7\end{matrix};1\right),\\
 F_3&={}_4F_3\!\left(\begin{matrix}1,8/7,9/7,12/7\\5/7,18/7,20/7\end{matrix};1\right),\\
 F_4&={}_3F_2\!\left(\begin{matrix}2,10/7,11/7\\20/7,22/7\end{matrix};1\right).
\end{align*}
Their summands are precisely $p_r(n)T_n$ in the respective rows. Formula
\eqref{extra7-eq:target4-sum} proves all four evaluations:
\begin{align*}
 F_1&=\frac{103243}{990}
       \frac{\Gamma(18/7)\Gamma(20/7)}{\Gamma(2/7)\Gamma(8/7)}
       -\frac{1313}{15},\\
 F_2&=\frac{1040}{3}-\frac{2107}{198}
       \frac{\Gamma(27/7)\Gamma(29/7)}{\Gamma(11/7)\Gamma(17/7)},\\
 F_3&=\frac{1573}{25}-\frac{1029}{50}
       \frac{\Gamma(18/7)\Gamma(20/7)}{\Gamma(8/7)\Gamma(9/7)},\\
 F_4&=78-\frac{147}{10}
       \frac{\Gamma(20/7)\Gamma(22/7)}{\Gamma(10/7)\Gamma(11/7)}.
\end{align*}

\subsubsection{Substitution in the Barnes reduction}

The Barnes-to-hypergeometric formula in \cite[p.~3]{extra8-zudilin}, with the
canceling upper and lower parameters removed, gives
\begin{align*}
 J_0&=-\frac{\pi}{\sin(2\pi/7)}(A_0F_1-B_0F_2),\\
 J_1&=\frac{\pi}{\sin(2\pi/7)}(A_1F_3-B_1F_4),
\end{align*}
where the gamma recurrence $\Gamma(z+1)=z\Gamma(z)$ reduces its prefactors to
\[
 A_0=-\frac{7g_1^2}{2574g_3g_6},\qquad
 B_0=\frac{2401g_4^2}{128700g_2g_6},\qquad
 A_1=\frac{35g_1^2}{2574g_3g_6},\qquad
 B_1=\frac{343g_4^2}{1170g_2g_6}.
\]
The same recurrence reduces the four proven evaluations to
\begin{align*}
 F_1&=-\frac{13(101g_1g_2-1204g_4g_6)}{15g_1g_2},\\
 F_2&=-\frac{520(43g_1g_6-98g_3g_4)}{147g_3g_4},\\
 F_3&=\frac{143(11g_1g_2-126g_4g_6)}{25g_1g_2},\\
 F_4&=-\frac{78(3g_1g_6-7g_3g_4)}{7g_3g_4}.
\end{align*}
Direct rational arithmetic, with no further gamma identities, now gives
\begin{align*}
 A_0F_1-B_0F_2
   &=\frac{7(101g_1^2g_2-2744g_3g_4^2)}{2970g_2g_3g_6},\\
 A_1F_3-B_1F_4
   &=\frac{7(11g_1^2g_2-294g_3g_4^2)}{90g_2g_3g_6}.
\end{align*}
These identities establish \eqref{extra7-eq:target4-J} and \eqref{extra7-eq:target4-K}.
Both defining integrals are positive and finite; in particular
$101-2744W\ne0$ follows from $J_0>0$ and $P>0$.

\subsubsection{Transcendence transfer}

The gamma recurrence also gives the exact beta quotient
\[
 W=\frac17\frac{B(4/7,4/7)}{B(1/7,2/7)}.
\]
For completeness, write
\[
 \epsilon_u(r,s)=\{ur\}+\{us\}-\{u(r+s)\},
\]
where braces denote fractional parts. The beta-quotient algebraicity
criterion of Wolfart--W\"ustholz \cite{dbkz-ww} implies that algebraicity of the
quotient here would require the two corresponding functions
$\epsilon_u$ to agree for every $u$ relatively prime to $7$. Already $u=1$
contradicts this:
\[
 \epsilon_1(4/7,4/7)=1,\qquad
 \epsilon_1(1/7,2/7)=0.
\]
Thus $W$ is transcendental. Finally the rational fractional-linear map in
\eqref{extra7-eq:target4-K} has nonzero determinant
\[
 9702\cdot101-363\cdot2744=-16170.
\]
It cannot send a transcendental argument to an algebraic value: an
algebraic $K_4$ would give
\[
 W=\frac{101K_4+363}{2744K_4+9702}\in\overline{\mathbb Q}.
\]
The denominator in this inverse cannot vanish for a value of the
nonconstant map. Therefore $K_4$ is transcendental.

\subsection{A Barnes evaluation and transcendence proof for the fifth denominator-seven target}
\label{extra7-sec:target5}

For $d\in\{4/7,11/7\}$, define the convergent positive integral
\[
 J_d=\iiint_{[0,1]^3}
 \frac{(1-x)^{3/7}y^{5/7}(1-y)^{2/7}
 z^{2/7}(1-z)^{3/7}}
 {(1-z+xyz)^d}\,dx\,dy\,dz,
 \qquad C_5=\frac{J_{11/7}}{J_{4/7}}.
\]
Thus $C_5=K(2/7,0,3/7,4/7,5/7)$ in the notation of
\cite{extra8-zudilin}. Both integrals are finite: near $z=1$, the
$z$-integral for the larger exponent is bounded by a constant
times $(xy)^{-1/7}$, which leaves integrable powers of $x$ and $y$.

\paragraph{Exact evaluation.}
Put $g_j=\Gamma(j/7)$ for $1\leq j\leq6$ and
\[
 T=\frac{g_3^2g_4}{g_5g_6^2}.
\]
Then
\begin{equation}
 \boxed{C_5=\frac{40(13-3T)}{88T-373}.}
 \label{extra7-eq:target5-evaluation}
\end{equation}
The approximate values are
\[
 T=4.270216179039710293695950611\ldots,
 \qquad C_5=2.725438564625374875727251628\ldots.
\]
The proof below uses Barnes' second lemma and explicit residues;
the numerical values are checks and play no role in the proof.

\paragraph{Mellin--Barnes reduction.}
The Mellin--Barnes representation in \cite{extra8-zudilin}, specialized to
$a=2/7,b=0,c=3/7,e=5/7$, gives
\[
 J_d=P_dM_d,\qquad
 P_d=\frac{\Gamma(9/7)\Gamma(10/7)}
 {\Gamma(d)\Gamma(19/7-d)},
\]
\[
 M_d=\frac{1}{2\pi i}\int_{\Re t=-1/2}
 \frac{\Gamma(d+t)\Gamma(9/7+t)\Gamma(12/7+t)
 \Gamma(-2-t)\Gamma(10/7-d-t)}{\Gamma(17/7+t)}\,dt.
\]
Here the simplifying identity is
\[
 \frac{\Gamma(1+t)\Gamma(-t)}{\Gamma(3+t)}
 =\frac{\Gamma(-t)}{(t+1)(t+2)}=\Gamma(-2-t).
\]
The chosen vertical line is valid for both values of $d$.

\paragraph{A contour identity with all residue corrections.}
Let $S=A+B+C+E-2$ and set
\[
 F(t)=\frac{\Gamma(A+t)\Gamma(B+t)\Gamma(C+t)
 \Gamma(-2-t)\Gamma(E-t)}{\Gamma(S+t)}.
\]
For complex parameters with
\[
 \Re A,\Re B,\Re C>\frac12,\qquad \Re E>-\frac12,
\]
one has, with removable singularities interpreted by continuation,
\begin{equation}
 \frac{1}{2\pi i}\int_{\Re t=-1/2}F(t)\,dt
 =G+R_{-2}+R_{-1},
 \label{extra7-eq:target5-barnes-corrected}
\end{equation}
where
\[
 G=\frac{
 \Gamma(A-2)\Gamma(B-2)\Gamma(C-2)
 \Gamma(A+E)\Gamma(B+E)\Gamma(C+E)}
 {\Gamma(A+B+E-2)\Gamma(A+C+E-2)
 \Gamma(B+C+E-2)},
\]
\[
 R_{-2}=-\frac{\Gamma(A-2)\Gamma(B-2)\Gamma(C-2)
 \Gamma(E+2)}{\Gamma(S-2)},\qquad
 R_{-1}=\frac{\Gamma(A-1)\Gamma(B-1)\Gamma(C-1)
 \Gamma(E+1)}{\Gamma(S-1)}.
\]

For a rigorous justification, first assume
$\Re A,\Re B,\Re C>9/4$ and $\Re E>0$.
The line $\Re t=-9/4$ separates the left poles from the right poles,
so Barnes' second lemma \cite{extra8-barnes} evaluates its integral as $G$.
Shifting that line to $\Re t=-1/2$ crosses precisely the poles
$t=-2,-1$ of $\Gamma(-2-t)$. Their residues are respectively
$R_{-2}$ and $R_{-1}$, since the residues of this gamma factor are
$-1$ and $+1$. Stirling's estimate gives exponential decay on the
horizontal sides. Both vertical lines are oriented upwards, and
the resulting identity is
$I(-1/2)=I(-9/4)+R_{-2}+R_{-1}$, with the plus signs displayed in
\eqref{extra7-eq:target5-barnes-corrected}.

The integral on the fixed line $\Re t=-1/2$ is holomorphic throughout
the larger connected parameter domain specified above. No numerator
pole meets that line; reciprocal gamma is entire; and Stirling's
estimate gives locally uniform integrable decay. The right-hand side
is meromorphic there. The meromorphic identity theorem therefore
continues the identity to that domain, and any apparent poles of the
sum on the right are removable. This establishes the formula at the
desired parameters without incorrectly using a separating straight
contour at their final values.

\paragraph{Application and exact gamma reduction.}
Take
\[
 (A,B,C,E)=\left(d,\frac97,\frac{12}7,\frac{10}7-d\right),
 \qquad S=\frac{17}7.
\]
Both $d=4/7$ and $d=11/7$ satisfy the preceding domain conditions.
For example, the Barnes term becomes
\[
 G_d=\frac{\Gamma(d-2)\Gamma(-5/7)\Gamma(-2/7)
 \Gamma(10/7)\Gamma(19/7-d)\Gamma(22/7-d)}
 {\Gamma(5/7)\Gamma(8/7)\Gamma(17/7-d)}.
\]
For a compact reduction table, define
\[
 p=\frac{g_2g_3}{g_1g_4},\qquad
 u=\frac{g_2g_3g_4^2}{g_6},\qquad
 v=\frac{g_2g_4g_5g_6}{g_3}.
\]
Successive uses of $\Gamma(z+1)=z\Gamma(z)$ give
\[
\begin{array}{c|cc}
 &d=4/7&d=11/7\\ \hline
 P_d&\frac34p&\frac32p\\[2pt]
 G_d&\frac{308}{75}u&-\frac{14}{5}u\\[2pt]
 R_{-2,d}&-\frac{637}{50}v&\frac{49}{5}v\\[2pt]
 R_{-1,d}&-\frac{14}{3}v&\frac73v
\end{array}
\]
Consequently,
\[
 J_{4/7}=\frac{7g_2^2}{200g_1g_6}
 \bigl(88g_3^2g_4-373g_5g_6^2\bigr),
\]
\[
 J_{11/7}=-\frac{7g_2^2}{5g_1g_6}
 \bigl(3g_3^2g_4-13g_5g_6^2\bigr).
\]
Their quotient proves \eqref{extra7-eq:target5-evaluation}.
In particular $88T-373>0$, as $J_{4/7}>0$.

\paragraph{Transcendence.}
The beta-function definition and $\Gamma(11/7)=(4/7)\Gamma(4/7)$
give
\[
 T=\frac74\frac{B(3/7,3/7)}{B(5/7,6/7)}.
\]
We use the following consequence of Wolfart--W\"ustholz
\cite{dbkz-ww}, Satz~4 and Proposition~2. For positive rational arguments,
with the individual arguments and their pair sums nonintegral,
an algebraic ratio of two beta values requires the equality of their
fractional-part profiles
\[
 h_{a,b}(r)=\{ra\}+\{rb\}-\{r(a+b)\}
\]
for every $r$ coprime to a common denominator.
The criterion allows a pair sum larger than one; these are among
the first- and second-kind beta periods treated in that theorem.
For the present ratio and $r=1$,
\[
 h_{3/7,3/7}(1)=0,\qquad h_{5/7,6/7}(1)=1.
\]
Thus $T$ is transcendental. The rational fractional-linear map in
\eqref{extra7-eq:target5-evaluation} has determinant
\[
 (-120)(-373)-520\cdot88=-1000\ne0.
\]
Equivalently,
\[
 T=\frac{373C_5+520}{88C_5+120},
\]
whose denominator is positive because $C_5>0$.
If $C_5$ were algebraic, this inverse formula would make $T$
algebraic. Hence
\[
 \boxed{K(2/7,0,3/7,4/7,5/7)\text{ is transcendental}.}
\]

\paragraph{Scope and priority.}
The published discussion in \cite{extra8-zudilin} explicitly lists this
target among the expressions involving ${}_3F_2$ values that seem
not reducible to a simpler form. The argument above supplies a gamma
evaluation and a transcendence proof using established theorems.
Historical priority of this particular reduction has not been
established. No claim is made here to prove the associated proposed
irrationality measure or divisibility assertions.



\clearpage

\part{Gamma quotient fields and continued fractions}
\section{Balanced Gamma quotients: classification and exact fields}
\label{fields-sec-gammafields}
This section concerns balanced products with two Gamma factors in each
numerator and denominator. It determines a finite arithmetic
classification and the fields generated by the algebraic expressions.
The motivating field-dimension question is recorded in
\cite{fields-MSE}; the unconditional transcendence input is the
beta-ratio theorem of Wolfart--W\"ustholz \cite{dbkz-ww}.
All counts below refer to expressions, with a quotient and its reciprocal
identified, rather than to pairwise distinct real numbers.


\begin{theorem}
Let $N\geq3$, $1\leq a\leq b<N$, $1\leq c\leq d<N$, and
$a+b=c+d\ne N$.  Put
\[
 R_N(a,b;c,d)=\frac{\Gamma(a/N)\Gamma(b/N)}
 {\Gamma(c/N)\Gamma(d/N)}.
\]
For a unit $u\pmod N$, define
\[
 h_{a,b}(u)=\frac{(ua\bmod N)+(ub\bmod N)
 -(u(a+b)\bmod N)}{N}\in\{0,1\}.
\]
Then $R_N(a,b;c,d)$ is algebraic if and only if
$h_{a,b}(u)=h_{c,d}(u)$ for every unit $u\pmod N$.
If one equality fails, the quotient is transcendental.
\end{theorem}

\begin{proof}
The balance condition cancels the common gamma factor and gives
\[
 R_N(a,b;c,d)=\frac{B(a/N,b/N)}{B(c/N,d/N)}.
\]
All four arguments are positive nonintegers and their common sum is a
noninteger.  The beta-ratio criterion of Wolfart--W\"ustholz therefore
applies, including when the common sum exceeds one.  Its character test is
precisely the equality of the displayed fractional-part signatures.
\end{proof}

The enumeration convention identifies a quotient with its reciprocal and
removes the identity by requiring $(a,b)<(c,d)$ in lexicographic order.
The expressions counted need not be pairwise distinct numbers.  An exact
integer enumeration gives
\[
\begin{array}{c|rrr|rrr}
&\multicolumn{3}{c|}{\text{denominator divides }N}
&\multicolumn{3}{c}{\text{exact common denominator }N}\\
N&\text{all}&\text{algebraic}&\text{transcendental}
 &\text{all}&\text{algebraic}&\text{transcendental}\\\hline
21&570&54&516&560&54&506\\
33&2480&30&2450&2420&30&2390
\end{array}
\]
Here ``exact common denominator'' imposes
$\gcd(N,a,b,c,d)=1$.  This distinction matters because the published
question uses the phrase ``overall denominator''.

For $N=21$, the nontrivial equal-signature classes consist of $18$ classes
of size $2$ and $12$ classes of size $3$, giving
$18+12\binom32=54$ algebraic quotient expressions.  For $N=33$, there are
$30$ classes of size $2$, giving $30$ algebraic quotient expressions.
The embedded certificate in Section~\ref{fields-sec-certificate}
reconstructs these finite lists using integer arithmetic. For each
transcendental expression, the first unit with unequal signature entries
is an independently checkable witness.

For example,
\[
 \frac{\Gamma(1/21)\Gamma(3/21)}{\Gamma(2/21)^2}
 \quad\text{and}\quad
 \frac{\Gamma(1/33)\Gamma(3/33)}{\Gamma(2/33)^2}
\]
are transcendental.  The witnesses are respectively $u=8$ and $u=13$:
the numerator signature equals $0$ and the denominator signature equals
$1$ in each case.



\subsection{The complete denominator-21 field}

Let $\mathcal A_{21}$ consist of the $54$ algebraic balanced quotient
representatives
\[
 r(a,b;c,d)=\frac{\Gamma(a/21)\Gamma(b/21)}
 {\Gamma(c/21)\Gamma(d/21)},
\]
where $1\leq a\leq b<21$, $1\leq c\leq d<21$,
$a+b=c+d\ne21$, and the unordered pairs are distinct. We choose one
orientation from each quotient and its reciprocal. Let
\[
 H_{21}=\langle r^7:r\in\mathcal A_{21}\rangle\subset\mathbb R_{>0},
 \qquad K_{21}=\mathbb Q(H_{21}).
\]
Thus $H_{21}$ includes arbitrary products and integer powers, as in the
source question~\cite{fields-MSE}. This section determines the whole field and
whole linear span, not just a lower bound from one counterexample.
The omitted reflection-only cases $a+b=c+d=21$ do not change the
answer: their values are sine quotients $s_c/s_a\in F$, and trivial
quotients equal $1$.

\begin{theorem}
Put
\[
 y=\left(\frac{\Gamma(1/21)\Gamma(4/21)}
 {\Gamma(2/21)\Gamma(3/21)}\right)^7,
 \qquad F=\mathbb Q(\zeta_{21})^+.
\]
Then
\[
 \boxed{K_{21}=F(y)=\mathbb Q(y),\qquad
 [K_{21}:\mathbb Q]=\dim_{\mathbb Q}\operatorname{span}_{\mathbb Q}H_{21}=36.}
\]
More precisely, among the $54$ chosen seventh powers, $18$ lie in $F$,
$18$ lie in $yF$, and $18$ lie in $y^{-1}F$.
\end{theorem}

\subsubsection{The cyclotomic containment lemma}
For $1\leq k\leq10$, set
\[
 q_k=2\sqrt3\sin(k\pi/21).
\]
These numbers belong to $F$. Indeed $\eta=e^{\pi i/21}$ lies in
$\mathbb Q(\zeta_{21})$, as does $\sqrt{-3}$, and
\[
 q_k=-\sqrt{-3}\,(\eta^k-\eta^{-k})
\]
is real. Consequently, whenever $C,h,m_1,\ldots,m_{10}$ are integers
and $h-M$ is even, where $M=\sum_{k=1}^{10}m_k$, one has
\begin{equation}\label{fields-eq:n21-cyclotomic-containment}
 2^C3^{h/2}\prod_{k=1}^{10}\sin(k\pi/21)^{m_k}
 =2^{C-M}3^{(h-M)/2}\prod_{k=1}^{10}q_k^{m_k}\in F.
\end{equation}
All the sine factors are nonzero, so negative exponents are harmless.

\subsubsection{Exact certificates for every quotient}
Write $g_k=\log\Gamma(k/21)$ and $s_k=\sin(k\pi/21)$. The identities
used in the verification are exactly
\begin{align*}
 R_k&:g_k+g_{21-k}=\log\pi-\log s_k
       &&(1\leq k\leq10),\\
 T_k&:g_k+g_{k+7}+g_{k+14}-g_{3k}
       =\log(2\pi)+(1/2-k/7)\log3
       &&(1\leq k\leq6),\\
 P_k&:\sum_{j=0}^{6}g_{k+3j}-g_{7k}
       =3\log(2\pi)+(1/2-k/3)\log7
       &&(k=1,2).
\end{align*}
These are reflection, triplication, and septuplication, with positive
real arguments. Thus rational linear combinations of their logarithms
are valid without any branch ambiguity.

The exact certificate in Section~\ref{fields-sec-certificate} enumerates the balanced pairs and applies
the previously established Wolfart--W\"ustholz test
\[
 \left\{\frac{ua}{21}\right\}+\left\{\frac{ub}{21}\right\}
 =\left\{\frac{uc}{21}\right\}+\left\{\frac{ud}{21}\right\}
 \quad\text{for every }u\in(\mathbb Z/21\mathbb Z)^\times.
\]
Exactly $54$ of the $570$ nontrivial representatives satisfy it.
For each of these $54$ targets $z=r(a,b;c,d)^7$, the certificate gives
an explicit rational linear combination of the displayed identities
and verifies an equality
\begin{equation}\label{fields-eq:n21-all-targets}
 \boxed{z=y^{\varepsilon}
  2^C3^{h/2}\prod_{k=1}^{10}s_k^{m_k},}
\end{equation}
where
\[
 \varepsilon\in\{-1,0,1\},\quad C,h,m_k\in\mathbb Z,
 \quad h-\sum_k m_k\in2\mathbb Z.
\]
The certificate constructs these identities by rational row reduction of
the displayed reflection and multiplication equations. It checks every
Gamma exponent vector, the vanishing of the remaining $\log\pi$ and
$\log7$ coefficients, and the stated parity condition. Its complete code
is embedded below; no external coefficient table is required.

For example, if $z=r(1,10;2,9)^7$, the table gives
\[
 z=y\,2^7\sqrt3\,s_4^7\in yF.
\]
For $z=r(2,8;4,6)^7$, it gives
\[
 z=y^{-1}\,2^7 3^{5/2}s_6^7\in y^{-1}F.
\]
The exact counts for $\varepsilon=-1,0,1$ are $18,18,18$ respectively.
Equation~\eqref{fields-eq:n21-cyclotomic-containment} therefore proves
$K_{21}\subseteq F(y)$.

\subsubsection{Equality and the exact dimension}
For completeness, the degree-$36$ argument is recalled here. Put
$t=2\cos(2\pi/21)$ and $R=(t+1)/t$. For an explicit seed certificate, the Gamma exponents cancel in
\[
 2\log R_{21}(1,4;2,3)
 =-R_2+R_4-R_6-R_8+T_1-T_2-T_3+P_1.
\]
Thus, with $s_k=\sin(k\pi/21)$,
\[
 R_{21}(1,4;2,3)=2\,3^{1/28}7^{1/12}
 \sqrt{s_2s_6s_8/s_4}
 =3^{1/28}7^{1/12}\sqrt{(t+1)/t}.
\]
Here $s_2/s_4=1/t$ and $s_6s_8=(t+1)/4$.
Consequently
\[
 y=3^{1/4}7^{7/12}R^{7/2},\qquad
 a=y^6=1029\sqrt{21}\,R^{21}\in F.
\]
The identity $\sqrt{21}\in F$ can be verified explicitly from
\[
 \sqrt{21}=2t^5+2t^4-10t^3-8t^2+12t+3,
 \quad t^6-t^5-6t^4+6t^3+8t^2-8t+1=0.
\]
If an automorphism $\sigma$ of $F$ fixes $a$, then
$\sigma(R)=\pm R$. The positive choice fixes $t$. The negative choice
would give $\sigma(t)=-t/(2t+1)\in(-1/2,0)$, while all six conjugates
$2\cos(2\pi u/21)$, for $u=1,2,4,5,8,10$, are positive or less
than $-1$. Hence $\mathbb Q(a)=F$.

The norm $\operatorname N_{F/\mathbb Q}(a)=-3^9 7^{42}<0$ shows
that $a$ is not a square in $F$. Also
\[
 a=(7\sqrt{21}R^7)^3/7
\]
shows it is not a cube: otherwise the abelian extension $F$ would
contain the nonnormal field $\mathbb Q(\sqrt[3]7)$.
Let $b=y^3=\sqrt a$. Then $[F(b):F]=2$. If $Z^3-b$ had a root
in $F(b)$, taking its norm to $F$ would make $-a$, hence $a$, a cube
in $F$. Thus $[F(y):F]=6$, and $[\mathbb Q(y):\mathbb Q]=36$.

Since $y\in H_{21}$, one has $\mathbb Q(y)\subseteq K_{21}$;
since $F=\mathbb Q(y^6)$, the upper bound already proved gives
$K_{21}=F(y)=\mathbb Q(y)$. Finally $1,y,\ldots,y^{35}$ belong to
$H_{21}$ and are linearly independent. Every group member lies in the
$36$-dimensional field $K_{21}$. Its linear span therefore has dimension
exactly $36$ as claimed.

The degree-$36$ seed is already a counterexample to a proposed
six-dimensional span of the multiplicative group of seventh powers:
$1,y,\ldots,y^{35}$ are linearly independent. This is a statement about
this explicitly defined group, not about all Gamma monomials.


\subsection{The denominator-33 field and individual degrees}
Write $g_j=\Gamma(j/33)$ and let $\mathcal A_{33}$ consist of the
nonidentity algebraic expressions
\[
 Q=\frac{g_ag_b}{g_cg_d},\qquad
 1\leq a\leq b<33,\quad 1\leq c\leq d<33,\quad
 a+b=c+d\ne33,
\]
with a quotient and its reciprocal identified.  The previously proved
Wolfart--W\"ustholz classification shows that there are exactly $30$ such
expressions.  This notation counts expressions, not necessarily distinct
real numbers.

\begin{theorem}
Let $\zeta=\exp(2\pi i/33)$ and
$F=\mathbb Q(\zeta+\zeta^{-1})$.  Then
\[
 \mathbb Q(Q^{11}:Q\in\mathcal A_{33})=F,
 \qquad [F:\mathbb Q]=10,
\]
and
\[
 \mathbb Q(Q:Q\in\mathcal A_{33})=F(3^{1/11}),
 \qquad [F(3^{1/11}):\mathbb Q]=110.
\]
In fact every individual $Q\in\mathcal A_{33}$ generates the respective
field: $\mathbb Q(Q^{11})=F$ and
$\mathbb Q(Q)=F(3^{1/11})$.
\end{theorem}

\begin{proof}
Put
\[
 t_j=2\sqrt3\sin(j\pi/33),\qquad 1\leq j\leq16.
\]
Since $\exp(\pi i/33)=-\zeta^{17}$ and
$\sqrt{-3}=1+2\zeta^{11}$, we have the exact cyclotomic expression
\[
 t_j=-(1+2\zeta^{11})
 \bigl((-\zeta^{17})^j-(-\zeta^{17})^{-j}\bigr).
\]
It is real, so $t_j\in F$.

The reflection and triplication identities take the form
\[
 g_jg_{33-j}=\frac{\pi}{\sin(j\pi/33)},\qquad
 \frac{g_jg_{j+11}g_{j+22}}{g_{3j}}
   =2\pi\,3^{1/2-j/11}\quad(1\leq j\leq10).
\]
For each of the $30$ expressions, the exact certificate in Section~\ref{fields-sec-certificate}
gives one triplication identity, with exponent $\epsilon\in\{1,-1\}$,
and a product of reflection identities.  The powers of $\pi$ cancel.
More explicitly, every expression has a certified identity
\[
 Q=3^{m/11}\prod_jt_j^{e_j},
 \qquad m\in\{-5,-4,-3,-2,-1,1,2,3,4,5\},
 \qquad e_j\in\{-1,0,1\}.
\]
Consequently all eleventh powers belong to $F$, and all original
expressions belong to $F(3^{1/11})$.  These $30$ identities are verified
by comparing their exponent vectors in the free abelian group on
$g_1,\ldots,g_{32}$, together with their explicit scalar factors.

For the reverse inclusion a single expression suffices:
\[
 Q_0=\frac{g_1g_{12}}{g_3g_{10}}
     =\frac{t_{10}}{3^{1/11}},\qquad
 y=Q_0^{11}=\frac{t_{10}^{11}}3.
\]
Its square is
\[
 y^2=3^9\bigl(2-2\cos(20\pi/33)\bigr)^{11}.
\]
The $10$ conjugates of $2-2\cos(20\pi/33)$ in $F$ are pairwise distinct
positive real numbers, because $\gcd(10,33)=1$. Raising them to the
eleventh power preserves distinctness. Thus
$[\mathbb Q(y^2):\mathbb Q]=10$, and
$\mathbb Q(y)=F$.  Moreover $3^{1/11}=t_{10}/Q_0$ belongs to
$\mathbb Q(Q_0)$, so
$\mathbb Q(Q_0)=F(3^{1/11})$.  The degrees of $F$ and
$\mathbb Q(3^{1/11})$ are respectively $10$ and $11$; the latter follows
from Eisenstein's criterion.  Their compositum therefore has degree
$110$.

For the stronger assertion about every individual expression, the
embedded certificate performs exact arithmetic in
$\mathbb Q[z]/(\Phi_{33}(z))$, applies every substitution $z\mapsto z^u$
for $u\in(\mathbb Z/33\mathbb Z)^\times$, and verifies exactly $10$
different conjugates of each $Q^{11}$.  This proves
$\mathbb Q(Q^{11})=F$.  Since $\gcd(m,11)=1$, the displayed formula for
$Q$ then recovers $3^{1/11}$ by B\'ezout's identity and proves
$\mathbb Q(Q)=F(3^{1/11})$ for every expression.
\end{proof}

\paragraph{Exact-computation specification.}
The modulus used by the exact polynomial calculation is
\[
 \Phi_{33}(z)=z^{20}-z^{19}+z^{17}-z^{16}+z^{14}-z^{13}
 +z^{11}-z^{10}+z^9-z^7+z^6-z^4+z^3-z+1.
\]
The embedded certificate verifies this polynomial by the identity
\[
 \Phi_{33}(z)=\frac{(z^{33}-1)(z-1)}{(z^{11}-1)(z^3-1)}.
\]
The calculation uses rational coefficients, derives all $30$ exponent
certificates directly from reflection and triplication, reconstructs
the complete fractional-part list, and counts exact Galois images.  Neither floating-point approximations nor
integer-relation guesses are used.

\paragraph{Relation to the published question.}
This proves the conjectured dimension $10$ for the field generated by
eleventh powers of the balanced four-gamma expressions at denominator
$33$.  However, the accompanying assertion that $10$ of the original
expressions have degree $55$ is false as written: all $30$ have degree
$110$.  No assertion here extends the classification to arbitrary
products of gamma values with more than two factors in numerator and
denominator.


\section{Generalized Brouncker fractions and effective decisions}\label{fields-sec-brouncker}

For $s,r\in\mathbb Q$, $s>0$, $r>1/2$, let
\[
 Y(s,r)=s+\underset{n=1}{\overset{\infty}{\mathbf K}}
 \frac{(2n-1)^2r^2-(r-1)^2}{2s}.
\]
Kushel's Theorem 1.4 \cite{fields-Kushel} gives the convergent value
\begin{equation}\label{fields-eq:brouncker-gamma}
 Y(s,r)=4r\frac{\Gamma(a+1/2)\Gamma(b+1/2)}{\Gamma(a)\Gamma(b)},
 \qquad a=\frac{s+1}{4r},\quad b=\frac{s+2r-1}{4r}.
\end{equation}
Notice that $a,b>0$ and $|a-b|<1/2$.

\paragraph{Theorem.}
If $2a\in\mathbb Z$ or $2b\in\mathbb Z$, then $Y(s,r)$ is
transcendental. Otherwise choose any even positive integer $N$ for which
$Na,Nb\in\mathbb Z$. Then $Y(s,r)$ is algebraic if and only if, for
every integer $k$ coprime to $N$, the two fractional parts $\{ka\}$
and $\{kb\}$ lie in opposite open intervals $(0,1/2)$ and $(1/2,1)$.
Consequently one same-half witness $k$ is a finite transcendence certificate.

\paragraph{Proof.}
Write $R(x)=\Gamma(x+1/2)/\Gamma(x)$. For an integer $m\geq1$,
$R(m)\in\mathbb Q^{\times}\sqrt\pi$, and for a positive nonintegral
half-integer $h$, $R(h)\in\mathbb Q^{\times}/\sqrt\pi$.
If both $a,b$ are half-integers, the strict bound $|a-b|<1/2$
forces $a=b$, so their product is a nonzero rational multiple of
$\pi$ or $1/\pi$ and is transcendental. If precisely one, say $a$,
is a half-integer, there are two possibilities. For integral $a$,
\[
 \sqrt\pi R(b)=b B(b+1/2,1/2),
\]
and for nonintegral half-integral $a$,
\[
 R(b)/\sqrt\pi=1/B(b,1/2).
\]
The beta arguments and their sum are nonintegers. Schneider's beta-value
theorem therefore proves transcendence. This also explains why the strict
parameter restriction matters: an integer and a nonintegral half-integer
can cancel their square-root-of-$\pi$ factors in unrestricted gamma products.

Now assume neither $2a$ nor $2b$ is integral. The gamma quotient equals
\begin{equation}\label{fields-eq:brouncker-beta}
 Y(s,r)=4r(a+b)\frac{B(a+1/2,b+1/2)}{B(a,b)}.
\end{equation}
First suppose $a+b\notin\mathbb Z$. All arguments and sums needed for
the Wolfart--W\"ustholz beta-ratio criterion \cite{dbkz-ww} are nonintegral.
For rational $u,v$ put
$h_{u,v}(k)=\{ku\}+\{kv\}-\{k(u+v)\}$.
That criterion says that the quotient in \eqref{fields-eq:brouncker-beta}
is algebraic exactly when its two profiles agree on $(\mathbb Z/N\mathbb Z)^\times$.
Every such $k$ is odd. If $x=\{ka\}$ and $y=\{kb\}$, then neither is
$0$ or $1/2$, and direct calculation gives
\[
 h_{a+1/2,b+1/2}(k)-h_{a,b}(k)
 =1-\mathbf1_{x>1/2}-\mathbf1_{y>1/2}.
\]
This difference vanishes exactly when $x$ and $y$ lie in opposite halves.

If $a+b\in\mathbb Z$, gamma recurrence and reflection directly make
the quotient algebraic: writing $t=\{a\}\in(0,1)\setminus\{1/2\}$,
one reduces it, up to a nonzero rational factor, to
\[
 \frac{\Gamma(t+1/2)\Gamma(3/2-t)}
 {\Gamma(t)\Gamma(1-t)}=(1/2-t)\tan(\pi t)\in\overline{\mathbb Q}.
\]
For every unit $k$, the corresponding fractional parts sum to $1$;
thus they lie in opposite halves, as required. This covers all cases.


For example, the opposite-half test gives
\[
 Y(2,3)=\frac{3^{1/4}(\sqrt6+\sqrt2)}2\in\overline{\mathbb Q},
\]
whereas $Y(2,4)$ is transcendental: $(a,b)=(3/16,9/16)$ and $k=3$
gives the same-half witness $(9/16,11/16)$. These are applications to a
pre-existing family, without a claim of historical priority.


\subsection{A terminating minimality and positivity algorithm}

The preceding finite classifier can be upgraded from existence of algebraic
minimal initial data to a decision procedure for \emph{every specified pair}
of real algebraic initial values. No general period-equality conjecture is
needed for this particular recurrence family.

\paragraph{Effective classical gamma certificates.}
Put $G(x)=\Gamma(x)/\sqrt\pi$ for $x\in\mathbb Q_{>0}$.
The following identities have explicitly computable positive real algebraic
right-hand sides:
\begin{align*}
 G(x+1)/G(x)&=x,\\
 G(x)G(1-x)&=1/\sin(\pi x)\quad(0<x<1),\\
 \frac{\prod_{j=0}^{m-1}G(x+j/m)}{G(mx)}
 &=2^{(m-1)/2}m^{1/2-mx}\quad(m\geq2),\\
 G(1/2)&=1.
\end{align*}
For each identity record its finitely supported rational exponent vector
in the symbols $E_x$, together with its algebraic constant. Enumerate the
positive rational arguments $x$ and positive integers $m$, and at every
finite stage solve a finite rational linear system to ask whether
\[
 T=E_{a+1/2}+E_{b+1/2}-E_a-E_b
\]
lies in the span of the enumerated exponent vectors. If
$T=\sum_j q_jR_j$, with identity constants $c_j>0$, then
\[
 \frac{\Gamma(a+1/2)\Gamma(b+1/2)}{\Gamma(a)\Gamma(b)}
 =\prod_j c_j^{q_j}.
\]
All rational powers mean positive real roots, so no branch ambiguity arises.
Real algebraic arithmetic computes the right-hand side exactly, for example
by a polynomial together with an isolating rational interval.

This search terminates whenever the classifier returns algebraic.
Wolfart--W\"ustholz, Satz 4, says that an algebraic quotient of two positive
rational beta values is classically proportional; their printed page 6
defines classical relations as consequences of the functional, reflection,
and distribution relations. Proposition 1 on that page is the finite
fractional-part test for these relations. Thus a finite certificate exists.
Arguments outside $(0,1)$ are reduced by recurrence; negative intermediate
arguments, if used in a classical derivation, can likewise be shifted to
positive ones without changing its validity up to explicit nonzero rational
factors. The two cases excluded from a direct nonresonant beta argument
are already covered in the classifier: half-integer boundaries are
transcendental, and integral $a+b$ has an explicit reflection evaluation.
Consequently no unproved general assertion about gamma algebraic relations
is being invoked.

\paragraph{Theorem (specified initial data).}
There is an unconditional algorithm which, given rational $r>1/2$, $s>0$
and real algebraic $v_{-1},v_0$, decides whether the nonzero solution of
\[
 (2rn+1)v_n=2s v_{n-1}+(2rn-1)v_{n-2}\qquad(n\geq1)
\]
is minimal.

\paragraph{Algorithm and proof.}
To check the continued-fraction normalization, put $P_{-1}=P_0=1$,
$P_n=\prod_{j=1}^n(2rj+1)$ and $u_n=P_nv_n$. Then
\[
 u_n=2s u_{n-1}+(2rn-1)(2r(n-1)+1)u_{n-2}.
\]
The second coefficient equals
$(2n-1)^2r^2-(r-1)^2$, exactly the numerator of the fraction defining
$Y(s,r)-s$. Scaling preserves ratios of solutions and therefore
minimality. Pincherle's theorem gives the unique minimal initial slope
$v_0/v_{-1}=s-Y(s,r)$.

The zero initial pair is separately identified; the theorem concerns
nonzero solutions. If $v_{-1}=0$, answer nonminimal. Otherwise run the
finite classifier for $Y(s,r)$. If it returns transcendental, answer
nonminimal, since Pincherle's criterion requires
$v_0/v_{-1}=s-Y(s,r)$, whereas the left-hand side is algebraic.
If it returns algebraic, run the terminating classical-certificate search
to obtain an exact algebraic representation of $Y(s,r)$. Decide the equality
\[
 v_0/v_{-1}=s-Y(s,r)
\]
using effective real algebraic arithmetic, answering minimal exactly when
it holds. Each branch terminates and the criterion is necessary and sufficient.

This decides every specified real algebraic input pair in the stated
two-parameter family. It resolves this
restriction of the polynomial-continued-fraction equality/minimality problem
discussed by Kenison et al.~\cite{fields-Kenison}. It does
not settle general second-order holonomic minimality, and the historical
priority of this particular corollary has not been established.

\paragraph{A positivity consequence with an elementary proof.}
The same algorithm decides positivity (all terms, including the specified
initial terms, nonnegative) for
\[
 (2rn+1)w_n=-2s w_{n-1}+(2rn-1)w_{n-2}\qquad(n\geq1).
\]
The zero sequence is positive. Every other positive sequence must have
$w_{-1}>0$: a negative initial term is forbidden, and $w_{-1}=0<w_0$
would force $w_1<0$. Normalize $w_{-1}=1$ and put $t=w_0/w_{-1}$.
After the same positive scaling as before, let $P_n/Q_n$ be the $n$th
convergent of the positive continued fraction $Y(s,r)-s$, with
$P_0=0$, $Q_0=1$, $P_{-1}=1$, $Q_{-1}=0$. Its numerator and denominator
recurrences have $b_n=2s$ and
$a_n=(2rn-1)(2r(n-1)+1)>0$.
The scaled solution is
\[
 W_n=(-1)^n(Q_nt-P_n).
\]
Consequently positivity is equivalent to $t$ lying above every even
convergent and below every odd convergent. Both subsequences converge
to $Y(s,r)-s$, so the condition is exactly
\[
 w_{-1}>0,\qquad w_0/w_{-1}=Y(s,r)-s,
\]
apart from the separately accepted zero sequence. This is an equality
between the same effectively classified constant and an algebraic input,
and is decided by the preceding algorithm. The proof uses convergence
already supplied by the generalized Brouncker formula and the elementary
ordering of positive continued-fraction convergents.


\section{Barnes reductions to beta values and logarithms}
\label{fields-sec:barnes-family}

For $e\in\mathbb Q$, $0<e<\tfrac12$, put $d=\tfrac12-e$ and define
\[
 J_{p,e}=\int_{[0,1]^3}
 \frac{\sqrt{x(1-x)(1-z)}\,y^e}
 {(1-z+xyz)^p}\,dx\,dy\,dz,
 \qquad K_e=\frac{J_{d+1,e}}{J_{d,e}}.
\]
Both integrals are finite and strictly positive. The underlying
Mellin--Barnes representation is given in \cite{extra8-zudilin}.  Every $K_e$ in this
family is transcendental.  The particular values $e=1/6$ and $e=1/3$ belong to this uniform family.  No assertion of historical
priority for this family is intended.

Here is an explicit formula, including the coefficient that certifies
that the fractional-linear transformation is nonconstant.  Write
\[
 B_e=B(2e,\tfrac12),\qquad
 A_e=\frac{2(2e+1)}{(1-e^2)(4e+1)},\qquad
 q_e=\frac{(e+1)(4e+1)}{(2e-1)(2e+1)},
\]
\[
 R_{0,e}=\frac{e+3}{e(e-1)(2e+3)},\qquad
 R_{1,e}=\frac{2e^2+5e+1}{e(e-1)(2e-1)(2e+1)}.
\]
Then the exact evaluations are
\begin{equation}
 J_{d,e}=\pi(A_e B_e+R_{0,e}),\qquad
 J_{d+1,e}=\pi(q_e A_e B_e+R_{1,e}),
 \label{fields-eq:barnes-family-evaluations}
\end{equation}
and consequently
\begin{equation}
 K_e=\frac{q_e A_e B_e+R_{1,e}}{A_e B_e+R_{0,e}},
 \qquad
 q_eR_{0,e}-R_{1,e}
 =\frac{1}{(2e-1)(2e+1)(2e+3)}\ne0.
 \label{fields-eq:barnes-family-ratio}
\end{equation}

\paragraph{Barnes reduction.}
It is useful first to allow $0<b<1$, $e>0$, $0<d<2-b$, and set
\[
 \mathcal J_d(b,e)=\int_{[0,1]^3}
 \frac{x^b(1-x)^{1-b}y^e(1-z)^{1-b}}
 {(1-z+xyz)^d}\,dx\,dy\,dz.
\]
For $d\ne1$ and $e\ne1$, the Mellin--Barnes representation and
\[
 \frac{\Gamma(1+s)\Gamma(-s)}{\Gamma(3+s)}
 =\Gamma(-2-s)
\]
give
\begin{align*}
 \mathcal J_d(b,e)
 &=\frac{\Gamma(2-b)}{\Gamma(d)\Gamma(3-b-d)}\frac{1}{2\pi i}
   \int_{\Re s=\sigma} H_d(s)\,ds,\\
 H_d(s)&=
 \frac{\Gamma(d+s)\Gamma(1+b+s)\Gamma(1+e+s)
       \Gamma(-2-s)\Gamma(2-b-d-s)}{\Gamma(2+e+s)}.
\end{align*}
For the two exponents $d,d+1$ one may use a common line
$\sigma=-\delta$, where
\[
 \max(0,b+d-1)<\delta<\min(d,1).
\]
This interval is nonempty since $0<b<1$ and $0<d<2-b$.
Barnes's second lemma, with the residues at $s=-2,-1$ included,
evaluates this integral as
\begin{equation}
 \mathcal J_d(b,e)=\frac{\pi}{\sin\pi b}\bigl(R_d+X_d\bigr),
 \label{fields-eq:barnes-general-reduction}
\end{equation}
where
\[
 R_d=\frac{3-b-d}{(d-1)(d-2)(e-1)}+
     \frac{1-b}{e(d-1)},
\]
\[
 X_d=-\frac{\Gamma(e-1)\Gamma(2-b)\Gamma(3-d)
                   \Gamma(3+e-b-d)}
 {(d-1)(d-2)\Gamma(3-b-d)\Gamma(1+e-b)\Gamma(2+e-d)}.
\]
The same expression, evaluated at $d+1$, is valid even when
$d+1\ge 2-b$: the contour condition just given and the convergence
conditions for the original integral are the ones used for that
evaluation.  Equivalently, this is the continuation of the displayed
formula to the second exponent.  The two elementary residue
contributions after multiplying by the prefactor are exactly
\[
 \frac{\pi}{\sin\pi b}
 \frac{3-b-d}{(d-1)(d-2)(e-1)},\qquad
 \frac{\pi}{\sin\pi b}\frac{1-b}{e(d-1)}.
\]

For completeness, the contour deformation does not presuppose a
separating line at the final parameters.  Set
$(A,B,C,E)=(d,1+b,1+e,2-b-d)$ and
$S=A+B+C+E-2$.
Initially take $\Re A,\Re B,\Re C>9/4$ and $\Re E>0$.
Barnes's lemma applies on the separating line $\Re s=-9/4$ to
\[
 \frac{\Gamma(A+s)\Gamma(B+s)\Gamma(C+s)
             \Gamma(-2-s)\Gamma(E-s)}{\Gamma(S+s)}.
\]
Moving the line to $-\delta\in(-1,0)$ crosses just $s=-2,-1$;
Stirling's formula makes the horizontal sides vanish.  With lines
oriented upwards the value on the new line is the Barnes value plus
these two residues.  Both sides then continue meromorphically to
the connected domain
$\Re A,\Re B,\Re C>\delta$, $\Re E>-\delta$.
The fixed-line integral is holomorphic there.  This domain contains
both required parameter points.  The exclusions above avoid the
individual gamma poles in the displayed formula; they are automatic
in the one-parameter family.

Gamma recurrence gives, identically,
\[
 X_{d+1}=q_dX_d,\qquad
 q_d=-\frac{(2-b-d)(1+e-d)}{d(2+e-b-d)},
\]
\[
 q_dR_d-R_{d+1}
 =\frac{b(b-1)}{d(d-2)(b+d-e-2)}.
\]
These expressions also explain why the generic reduction alone does
not prove transcendence: a nonconstant fractional-linear expression
still requires an arithmetic theorem for its gamma-value parameter.

\paragraph{The beta value on the chosen line.}
Put $r=1-b$ and $t=2-d$ and introduce
\[
 S(e,r,t)=\frac{\Gamma(e)\Gamma(r)\Gamma(t)\Gamma(e+r+t)}
 {\Gamma(e+r)\Gamma(r+t)\Gamma(t+e)}.
\]
Then $X_d=rS(e,r,t)/((e-1)(d-1))$.
On our line $b=1/2$, $d=1/2-e$, hence $r=1/2$ and $t=e+3/2$.
Gamma recurrence yields
\[
 S(e,\tfrac12,e+\tfrac12)=2B(2e,\tfrac12),
\qquad
 \frac{S(e,\tfrac12,e+\tfrac32)}{S(e,\tfrac12,e+\tfrac12)}
 =\frac{(2e+1)^2}{(e+1)(4e+1)}.
\]
Substitution gives $X_d=A_eB_e$, $R_d=R_{0,e}$,
$R_{d+1}=R_{1,e}$, and $q_d=q_e$, establishing
\eqref{fields-eq:barnes-family-evaluations}.

To check convergence directly, it suffices to examine $z$ near $1$.
Writing $u=1-z$ there, the denominator is comparable to $u+xy$.
For a general exponent $p$, integration in $u$ is bounded by a
constant, a logarithm, or a multiple of $(xy)^{2-b-p}$.
The resulting endpoint integrals are finite when
$p<\min(3-b,3,3+e-b)$.  Both $p=d$ and $p=d+1$ satisfy this
condition in the stated domain.  Positivity is immediate from the
integrands and proves that the denominators in the ratio formula
do not vanish.

\paragraph{Transcendence.}
For rational $0<e<1/2$, the positive rational arguments $2e$ and
$1/2$ are both nonintegers.  Their sum is nonintegral unless
$e=1/4$.  Schneider's beta-value theorem therefore proves that
$B_e$ is transcendental for $e\ne1/4$; at $e=1/4$ it is simply
$B(1/2,1/2)=\pi$.  All coefficients in
\eqref{fields-eq:barnes-family-ratio} are rational, $A_e\ne0$, and the
displayed determinant is nonzero.  If $K_e$ were algebraic, inversion
of this fractional-linear map would make $B_e$ algebraic, a
contradiction.  This proves the stated infinite-family result.

As checks on the normalization, at $e=1/6$, writing
$V=B(1/3,1/2)$, the formula becomes
\[
 K_{1/6}=\frac{35(17-4V)}{2(32V-133)}.
\]
At $e=1/3$, writing $H=B(2/3,2/3)/(3\,2^{2/3})$ and using
$B(2/3,1/2)=6H$, it becomes
\[
 K_{1/3}=\frac{2002-4620H}{825H-350}.
\]



\subsection{The resonant line and a two-parameter logarithmic family}
\label{fields-sec:barnes-logarithmic-family}

The pole collision at $e=1$ in the gamma-value formula admits a
different arithmetic treatment.  This section uses Baker's theorem
on logarithms \cite{fields-Baker} rather than an unjustified classification of arbitrary
gamma products.

Let $b,d\in\mathbb Q$ satisfy
\[
 0<b<1,\qquad 0<d<2-b,\qquad d\ne1,
\]
and define
\[
 J_p=\int_{[0,1]^3}
 \frac{x^b(1-x)^{1-b}y(1-z)^{1-b}}
      {(1-z+xyz)^p}\,dx\,dy\,dz,
 \qquad K(b,d)=\frac{J_{d+1}}{J_d}.
\]
Then $K(b,d)$ is transcendental for every rational pair in this
domain.  Both integrals are finite and positive.  This is an
extension proved here; it is not asserted that the members of this
family were individually published as open problems or that the
general statement has no prior occurrence.

Put $r=1-b$, $t=2-d$, so that
\[
 0<r<1,\qquad 1-r<t<2,\qquad t\ne1.
\]
Write $\psi(z)=\Gamma'(z)/\Gamma(z)$ and
\[
 \Psi=\psi(1)+\psi(1+r+t)-\psi(1+r)-\psi(1+t).
\]
The exact evaluations are
\begin{align}
 \frac{J_d}{\pi/\sin\pi b}
 &=\frac{r+(r+t)\Psi/t}{d-1}=\alpha\Psi+\beta,
 \label{fields-eq:resonant-J0}\\
 \frac{J_{d+1}}{\pi/\sin\pi b}
 &=\frac{r+\dfrac{r+t-1}{t-1}
           \left(\Psi+\dfrac{r}{t(r+t)}\right)}{d}
   =\alpha_1\Psi+\beta_1.
 \label{fields-eq:resonant-J1}
\end{align}
All four coefficients are rational, and their determinant is
\begin{equation}
 \alpha_1\beta-\alpha\beta_1
 =\frac{r(r-1)}{t^2(t-2)(t-1)}\ne0.
 \label{fields-eq:resonant-determinant}
\end{equation}
Thus it suffices to prove that $\Psi$ is transcendental.

\paragraph{The limiting formula.}
For $e$ near $1$, replace $y$ by $y^e$ and call the resulting integral
$J_p(e)$. The preceding Barnes reduction gives
\[
 J_p(e)=\frac{\pi}{\sin\pi b}\{R_p(e)+X_p(e)\},
\]
with $R_p,X_p$ defined in the preceding subsection. Its fixed contour and
meromorphic continuation apply to $p=d,d+1$ in the present domain.
This explicitly retains the residues at $-2,-1$ before taking the
coalescing-pole limit.

Define
\[
 S(e,r,t)=\frac{\Gamma(e)\Gamma(r)\Gamma(t)\Gamma(e+r+t)}
 {\Gamma(e+r)\Gamma(r+t)\Gamma(t+e)}.
\]
Gamma recurrence rewrites $X_d(e)=rS(e,r,t)/((e-1)(d-1))$.
In particular,
\[
 S(1,r,t)=\frac{r+t}{rt},\qquad
 \left.\frac{\partial}{\partial e}\log S(e,r,t)\right|_{e=1}
 =\Psi.
\]
The pole parts in
\[
 R_d+X_d
 =-\frac{r+t}{(d-1)t(e-1)}+
   \frac{r}{e(d-1)}+
   \frac{rS(e,r,t)}{(e-1)(d-1)}
\]
cancel.  Taking the finite part gives
\eqref{fields-eq:resonant-J0}.  For the exponent $d+1$, replace $t$ by
$t-1$.  Gamma recurrence, or the digamma recurrence, gives
\[
 \Psi(r,t-1)=\Psi(r,t)+\frac{r}{t(r+t)},
\]
which proves \eqref{fields-eq:resonant-J1}.  Continuity of the integrals
at $e=1$ follows by domination with $y^{1-\varepsilon}$ for a fixed
sufficiently small positive $\varepsilon$.  All required convergence
inequalities are strict in the stated domain.

\paragraph{Arithmetic of the digamma combination.}
For rational $0<x<1$, Gauss's digamma formula \cite{fields-DLMF} expresses
$\psi(x)+\gamma$ as an algebraic multiple of $\pi$ plus an
algebraic linear form in real logarithms of positive algebraic
numbers.  Integer shifts add rational numbers.  Baker's theorem
says that a nonzero algebraic linear form in logarithms of
algebraic numbers is transcendental; it also says that logarithms
which are linearly independent over $\mathbb Q$ are linearly
independent over $\overline{\mathbb Q}$.

Suppose first that $r+t\notin\mathbb Z$.  The three numbers
$r,t,r+t$ are then nonintegral.  The coefficient of $\pi$ in
Gauss's expression for $\Psi$ is
\[
 c=-\frac12\bigl(\cot\pi(r+t)-\cot\pi r-\cot\pi t\bigr).
\]
This is a nonzero algebraic number.  In fact, writing
$u=\cot\pi r$, $v=\cot\pi t$, we have $u+v\ne0$ and
\[
 \cot\pi(r+t)-\cot\pi r-\cot\pi t
 =-\frac{u^2+uv+v^2+1}{u+v}\ne0.
\]
After a rational constant is subtracted, $\Psi$ is therefore a
linear form containing $c\pi=-ic\log(-1)$ with nonzero
coefficient, together with logarithms of positive real algebraic
numbers.  The logarithm $\log(-1)=i\pi$ is linearly independent
over $\mathbb Q$ of those real logarithms.  Choosing a rational
basis for the latter and applying Baker's linear-independence
theorem proves that this whole linear form cannot vanish.
Baker's transcendence theorem then proves that $\Psi$ is
transcendental.

If $r+t\in\mathbb Z$, the inequalities imply $r+t=2$.
Digamma recurrence shows that $\Psi$ differs by a rational number
from
\[
 F(r)=-2\gamma-\psi(r)-\psi(1-r).
\]
The cotangent terms in Gauss's formula cancel, so $F(r)$ is an
algebraic linear form in logarithms of positive algebraic
numbers, with no added algebraic constant.  It is nonzero:
$\psi$ is strictly increasing on the positive real axis, hence
$\psi(r)<\psi(1)=-\gamma$ and
$\psi(1-r)<\psi(1)=-\gamma$, giving $F(r)>0$.
Baker's theorem again proves transcendence.  This settles every
case in the stated domain.  Finally,
\eqref{fields-eq:resonant-determinant} transfers this transcendence to
$K(b,d)$.

For example, the choice $b=d=1/2$ gives
\[
 \Psi=4\log2-\frac{19}{6},\qquad
 J_{1/2}=\frac{\pi}{9}(67-96\log2),\qquad
 J_{3/2}=\pi(16\log2-11),
\]
and therefore
\[
 K(\tfrac12,\tfrac12)
 =\frac{144\log2-99}{67-96\log2}
 =1.77603428343791655449\ldots.
\]
The excluded line $d=1$ involves a further collision and is not
covered by this digamma argument.



\section{Embedded exact certificate for the two Gamma fields}
\label{fields-sec-certificate}
The finite enumeration and field-containment steps used above involve
integer arithmetic and rational linear algebra. The following executable
certificate is included to make every finite assertion reproducible from
this single source. It requires Python~3 and SymPy but reads no input
files and makes no numerical approximations. Assertions verify the
counts, every classical Gamma exponent identity, all cyclotomic
containment conditions, and all thirty degree-$10$ Galois orbits.
The mathematical degree-$36$ seed proof and the degree-$110$ extension
argument are the conceptual proofs given above; the program does not
replace the beta-ratio theorem by a numerical test.

{\scriptsize
\begin{verbatim}
# Exact finite certificates; Python 3 with SymPy. No input files.
from collections import defaultdict, Counter
from itertools import combinations
from math import gcd
import sympy as S

def classified(N):
    units=[u for u in range(1,N) if gcd(u,N)==1]
    groups=defaultdict(list)
    for a in range(1,N):
        for b in range(a,N):
            if a+b!=N: groups[a+b].append((a,b))
    def sig(a,b):
        return tuple(((u*a)%N+(u*b)%N-(u*(a+b))%N)//N
                     for u in units)
    total=exact=0; alg=[]; exact_alg=0
    for group in groups.values():
        for (a,b),(c,d) in combinations(group,2):
            q=(a,b,c,d)
            primitive=gcd(N,gcd(a,gcd(b,gcd(c,d))))==1
            total+=1; exact+=primitive
            if sig(a,b)==sig(c,d):
                alg.append(q); exact_alg+=primitive
    return sorted(alg), (total,len(alg),exact,exact_alg)

def target(q,N,power=1):
    v=[S.S(0)]*(N-1)
    for k in q[:2]: v[k-1]+=power
    for k in q[2:]: v[k-1]-=power
    return v

# N=21: exact row reduction of classical identities.
A21,counts=classified(21)
assert counts==(570,54,560,54)
L=[]; W=[]
for k in range(1,11):
    v=[S.S(0)]*20; w=[S.S(0)]*14
    v[k-1]=v[20-k]=1; w[3]=1; w[3+k]=-1
    L.append(v); W.append(w)
for step,m,K in [(7,3,range(1,7)),(3,7,range(1,3))]:
    for k in K:
        v=[S.S(0)]*20; w=[S.S(0)]*14
        for j in range(m): v[k+step*j-1]+=1
        v[m*k-1]-=1
        w[0]=w[3]=S.Rational(m-1,2)
        w[1 if m==3 else 2]=S.Rational(1,2)-S.Rational(m*k,21)
        L.append(v); W.append(w)
L=S.Matrix(L); W=S.Matrix(W)
def express(q):
    v=S.Matrix(target(q,21,7))
    c,parameters=L.T.gauss_jordan_solve(v)
    c=c.subs({p:0 for p in parameters})
    assert L.T*c==v
    return list(c.T*W)
base=express((1,4,2,3)); split=Counter()
assert base==[7,S.Rational(1,4),S.Rational(7,12),0,
             0,S.Rational(7,2),0,-S.Rational(7,2),0,
             S.Rational(7,2),0,S.Rational(7,2),0,0]
for q in A21:
    v=express(q); epsilon=v[2]/base[2]
    assert epsilon in (-1,0,1)
    rem=[a-epsilon*b for a,b in zip(v,base)]
    C,h=rem[0],2*rem[1]; M=sum(rem[4:])
    assert rem[2:4]==[0,0]
    assert all(x.q==1 for x in [C,h]+rem[4:])
    assert (h-M)%2==0
    split[int(epsilon)]+=1
assert split=={-1:18,0:18,1:18}
print('N21:',counts,'containment classes:',dict(split))

# N=33: derive all identities, then count exact Galois orbits.
A33,counts=classified(33)
assert counts==(2480,30,2420,30)
z=S.Symbol('z'); phi=S.Poly(S.cyclotomic_poly(33,z),z,domain=S.QQ)
assert S.expand(phi.as_expr()*(z**11-1)*(z**3-1)
                -(z**33-1)*(z-1))==0
one=S.Poly(1,z,domain=S.QQ)
def mul(a,b): return (a*b).rem(phi)
def power(a,n):
    if n<0: a=S.invert(a,phi); n=-n
    out=one
    while n:
        if n%2: out=mul(out,a)
        a=mul(a,a); n//=2
    return out
Z=[S.Poly(z**k,z,domain=S.QQ).rem(phi) for k in range(33)]
ts={j:mul(-(-1)**j*(one+2*Z[11]),
          Z[(17*j)%33]-Z[(-17*j)%33]) for j in range(1,17)}
pows={(j,e):power(t,11*e) for j,t in ts.items() for e in (-1,1)}
units=[u for u in range(1,33) if gcd(u,33)==1]
for q in A33:
    v=target(q,33); found=[]
    for j in range(1,11):
        for sign in (-1,1):
            w=v.copy()
            for k in (j,j+11,j+22): w[k-1]-=sign
            w[3*j-1]+=sign
            if all(w[k-1]==w[32-k] for k in range(1,17)):
                refs={k:w[k-1] for k in range(1,17) if w[k-1]}
                found.append((j,sign,refs))
    assert found
    j,sign,refs=found[0]; m=-sign*j
    assert sum(refs.values())==-sign and 1<=abs(m)<=5
    assert all(v in (-1,1) for v in refs.values())
    y=S.Poly(S.Rational(3)**m,z,domain=S.QQ)
    for k,e in refs.items(): y=mul(y,pows[k,-e])
    images=[]
    for u in units:
        image=S.Poly(0,z,domain=S.QQ)
        for (k,),coef in y.terms(): image+=coef*Z[(u*k)%33]
        images.append(tuple(image.all_coeffs()))
    assert len(set(images))==10
print('N33:',counts,'each Q^11 has 10 images; each Q degree 110')
\end{verbatim}
}
The first output has counts $(570,54,560,54)$ and containment multiplicities
$18,18,18$. The second has counts $(2480,30,2420,30)$ and ten distinct
Galois images for each eleventh power. In each count tuple the entries are
(total expressions, algebraic expressions, exact-denominator expressions,
exact-denominator algebraic expressions).

\clearpage

\part{Incomplete periods and arithmetic quantiles}
\section{Closed periods and the length of polynomial lemniscates}
\label{beta-sec-length}

The passage from an integral representation to a transcendence statement
requires attention to the integration cycle.  In particular, a sum or a
limit of transcendental lengths need not be transcendental.  The following
construction puts the entire length on one connected algebraic curve.

\begin{theorem}[Polynomial lemniscate length]\label{beta-thm-length}
Let $P\in\overline{\mathbb Q}[z]$ be nonconstant.  The Euclidean length of
$\Lambda_P=\{z\in\mathbb C:|P(z)|=1\}$ is transcendental.  The statement
allows several components, critical levels, and repeated polynomial
factors.  Length is measured as the length of the set, without algebraic
multiplicity.  Each geometric connected component also has
transcendental length.
\end{theorem}

The question motivating this statement appears in
\cite[Remark~8.16]{beta-tamiozzo}.  The proof below is a proposed resolution
of that stated mathematical question; historical priority and external
refereeing are not asserted.  The transcendence input is the established
closed-period theorem: a nonzero integral of an algebraic rational
differential over a closed integral cycle avoiding its poles, on a smooth
connected projective curve over $\overline{\mathbb Q}$, is transcendental
\cite[Theorem~13.9 and Corollary~13.13]{dbkz-hw}.

\begin{proof}
Write $P=Q^d$, choosing algebraic constants so that the multiplicities
$m_1,\ldots,m_r$ of the distinct zeros of $Q$ have greatest common divisor
one.  Then $\Lambda_P=\Lambda_Q$.  Put $n=\deg Q$ and let $\overline Q$
denote coefficientwise complex conjugation.  We first show that
\begin{equation}\label{beta-eq-complexification}
Q(z)\overline Q(w)=1
\end{equation}
is geometrically irreducible.  Over $\overline{\mathbb Q}(z)$ consider
$\overline Q(w)-1/Q(z)$.  At a zero of $Q$ of multiplicity $m_j$, every
root $\xi$ in a splitting field has valuation $-m_j/n$: it cannot have
nonnegative valuation, and for negative valuation the leading term of
$\overline Q(\xi)$ uniquely has smallest valuation.  A monic factor of
degree $e$ would consequently have constant coefficient of valuation
$-em_j/n$, an integer.  Thus $n\mid em_j$ for every $j$, whence $n\mid e$.
No proper factor exists.  Gauss's lemma gives the assertion for
\eqref{beta-eq-complexification}.

Set $z=x+iy$, $w=x-iy$, and normalize the projective closure to obtain
a connected smooth projective curve $X$ with a real algebraic model.
The highest homogeneous part of its plane equation is
$|a_n|^2(x^2+y^2)^n$.  Thus there are no real points at infinity.
The real locus $X(\mathbb R)$ is a finite union of circles.  It is
nonempty, since a regular value of $Q$ on the unit circle has a nonempty
inverse image.

We check that the normalization map is an immersion along this real
locus, including above critical points.  Near a point $z_0$ with
$|Q(z_0)|=1$, choose a local logarithm of $Q(z)/Q(z_0)$.  A local
biholomorphic coordinate $u=h(z)$ makes this logarithm $u^m$, for some
$m\ge1$.  The level set is then $\operatorname{Re}(u^m)=0$, a union of
$m$ distinct real lines.  The normalization separates these lines, and
$h^{-1}$ has nonzero derivative.  Consequently
\[
q=(dx)^2+(dy)^2=dz\,dw
\]
is positive and nonvanishing on every real tangent line of $X(\mathbb R)$.
Normalization is one-to-one off finitely many crossing points, which
have length zero, so adding the lengths of these circles counts the
original set exactly once.

Choose a nonzero algebraic rational differential $\tau$ on $X$ and write
$q=f\tau^2$.  If $f$ is not a square in $\overline{\mathbb Q}(X)$, let
$Y$ be the smooth connected curve with function field
$\overline{\mathbb Q}(X)(\sqrt f)$ and put
$\omega=\sqrt f\,\tau$.  If $f$ is a square, take $Y=X$ with either global
square root.  In both cases $\omega^2=q$ after pullback.

For the nonsquare case, orient a real circle and parametrize it smoothly
by $c(t)$.  There is a unique lift at each point for which
\[
\omega(\dot{\widetilde c}(t))
=\sqrt{\dot x(t)^2+\dot y(t)^2}>0.
\]
The square-root cover is unramified here because $q$ does not vanish.
This choice is continuous, and after one circuit both the position and
the tangent vector have returned to their starting values.  Uniqueness
therefore makes the lift closed.  In the square case the global square
root has a constant sign on the tangent direction of each real circle;
orient the circles separately to make their integrals positive.
There are no poles on the resulting lifted cycles.  Their sum $\Gamma$
is a closed integral cycle on the single connected curve $Y$, and
\[
\int_\Gamma\omega=\operatorname{length}(\Lambda_Q)>0.
\]
The closed-period theorem proves transcendence.  Selecting only the
normalization circles belonging to a specified geometric component gives
the componentwise statement.
\end{proof}

The same conclusion holds for $|P(z)|=\rho$ with algebraic $\rho>0$, by
replacing $P$ with $P/\rho$.  The nonvanishing of the real tangent vector
is essential: the algebraic curve parametrized by
$z(t)=(1+e^{it})^2$, $-\pi\le t\le\pi$, has a cusp and length
$\int_{-\pi}^{\pi}4\cos(t/2)\,dt=16$.  It is not a counterexample to
Theorem~\ref{beta-thm-length}, because its normalization map is not an
immersion at the cusp.

\subsection{A length certificate for a cubic lemniscate}

For $L_3=\operatorname{length}\{|z^3-z|=1\}$, the theorem yields
transcendence.  One useful positive series is obtained from
\[
y(t)^3-ty(t)=1,\quad y(0)=1,\qquad
G(t)=\left(\frac{3}{3y(t)^2-t}\right)^{1/2}
=\sum_{j\ge0}g_jt^j,\quad G(0)=1.
\]
The coefficients are rational.  The map
$\phi(u)=u\,y(u^{-2})$ parametrizes the lemniscate once on $|u|=1$;
its three branches above $u^3$ remain distinct because the discriminant
does not vanish there.  Since $|\phi'(u)|=|G(u^{-2})|^2$, Fourier
orthogonality gives
\begin{equation}\label{beta-eq-cubic-length}
L_3=2\pi\sum_{j\ge0}g_j^2,\qquad
\sum_{j>M}g_j^2\le\frac{1089}{40}(4/9)^{M+1}.
\end{equation}
For completeness, the bound follows by using $|t|=3/2$.  Every root of
$y^3-ty-1$ has modulus less than $3/2$, hence
$|3y_j^2-t|\le33/4$, while
$\prod_j(3y_j^2-t)=27-4t^3$ has modulus at least $27/2$.
Thus $|G|^2\le121/8$, and Cauchy's estimate gives
\eqref{beta-eq-cubic-length}.  It yields
\[
L_3=6.46470632781347726562312673284291603278\ldots.
\]
This proves neither transcendence of $L_3/\pi$ nor algebraic independence
of different lemniscate lengths.  It is unrelated to maximizing length
among monic polynomials of a fixed degree.

\section{A logarithm lemma for cyclotomic abelian components}
\label{beta-sec-loglemma}

We use the analytic subgroup theorem in its strengthened form
\cite[Theorem~6.2]{dbkz-hw}: if an algebraic linear subspace of the Lie
algebra of a connected commutative algebraic group contains a logarithm
of an algebraic point, it contains the Lie algebra of a connected
algebraic subgroup containing that logarithm.  The identity exponential
image is allowed.  The following formulation also accommodates the
additive directions needed for second-kind differentials.

\begin{lemma}[Abelian projection of logarithm relations]
\label{beta-lem-affine-log}
Let
\[
0\longrightarrow V\longrightarrow G\xrightarrow{\rho}A\longrightarrow0
\]
be an extension over $\overline{\mathbb Q}$ of an abelian variety by a
vector group.  Suppose an automorphism $T$ of $G$ descends to $A$, its
characteristic polynomial on $H_1(A,\mathbb Q)$ is the single
cyclotomic polynomial $\Phi_M$, and an algebraic invariant differential
$\eta$ on $G$ satisfies $T^*\eta=\xi\eta$, with $\xi$ a primitive
$M$th root of unity.  Assume $\eta$ has at least one nonzero period.
Identify $K=\mathbb Q[T|_A]$ with $\mathbb Q(\xi)$ using this eigenvalue.
Then the following statements hold.
\begin{enumerate}
\item If $w$ is a logarithm of an algebraic point of $G$ and
$\eta(w)\in\overline{\mathbb Q}$, then $d\rho(w)=0$.
\item If $u,v$ are such logarithms, $d\rho(v)\ne0$, and
$\eta(u)-\lambda\eta(v)=c$ with $\lambda,c\in\overline{\mathbb Q}$,
then $\lambda\in K$ and
\[
d\rho(u)=da\bigl(d\rho(v)\bigr),
\]
where $a\in K\subset\operatorname{End}^0(A)$ corresponds to $\lambda$.
In particular, if $d\rho(v)$ is a rational period, so is $d\rho(u)$.
\end{enumerate}
\end{lemma}

\begin{proof}
The vector group is contractible, so $H_1(G,\mathbb Q)=H_1(A,\mathbb Q)$.
The kernel of the period map $\mathcal P_\eta$ is $T$-invariant.  By
irreducibility of $\Phi_M$ and the assumed nonzero period, it is zero.
Moreover, since $H_1(A,\mathbb Q)$ is one-dimensional over $K$,
\begin{equation}\label{beta-eq-period-line}
\operatorname{im}\mathcal P_\eta=\Omega K\quad(\Omega\ne0).
\end{equation}
These are assertions about rational homology, not its complexification.

For the first assertion, apply the analytic subgroup theorem to $(w,1)$
in $G\times\mathbb G_a$ and the hyperplane
$\eta(Z)-\eta(w)t=0$.  Let $H$ be the subgroup obtained, and let $B$ be
its image in $A$.  The kernel of $H\to B$ is a vector group; hence a
nonzero rational cycle of a positive-dimensional $B$ lifts to $H$.
Its image in $H_1(A,\mathbb Q)$ is nonzero.  But the hyperplane equation
forces its $\eta$-period to vanish, since $dt$ has zero periods.
This contradicts injectivity.  Thus $B=0$ and $d\rho(w)=0$.

For the second assertion apply the theorem to $(u,v,1)$ in
$G^2\times\mathbb G_a$ and the corresponding homogeneous hyperplane.
Its subgroup has positive-dimensional second abelian projection because
$d\rho(v)\ne0$.  A rational cycle with nonzero second projection gives
\[
\lambda=
\frac{\mathcal P_\eta(\gamma_1)}{\mathcal P_\eta(\gamma_2)}\in K
\]
by \eqref{beta-eq-period-line}.  Choose $A_0(X)\in\mathbb Q[X]$ with
$A_0(\xi)=\lambda$.  The quasi-endomorphism $A_0(T)$ exists on $G$ and
descends to $a$ on $A$.  Clear denominators and apply the first assertion
to $m u-d(mA_0(T))v$, whose $\eta$-value is $mc$.  This gives the desired
abelian projection identity.  Quasi-endomorphisms preserve rational
periods.
\end{proof}

When $G=A$ and $\eta$ is a nonzero holomorphic eigendifferential, the
nonzero-period hypothesis follows automatically: a holomorphic
differential vanishing on every lattice period is zero.  Thus in this
case $\eta(w)=0$ forces $w=0$.  If $v$ is a nonzero logarithm of a torsion
point, the same lemma also gives the useful equivalence
\begin{equation}\label{beta-eq-ratio-torsion}
\frac{\eta(u)}{\eta(v)}\in\overline{\mathbb Q}
\quad\Longleftrightarrow\quad
\exp_A(u)\text{ is torsion},
\end{equation}
provided $u$ is a logarithm of an algebraic point.  The reverse implication
uses \eqref{beta-eq-period-line}.  An arbitrary algebraic scalar has not
been treated as an endomorphism: membership in $K$ was proved first.

\section{Incomplete hyperelliptic integrals and leaf division}
\label{beta-sec-leaves}

For $N\ge1$ put
\[
I_N(s)=\int_0^s\frac{dt}{\sqrt{1-t^N}},\qquad
R_N(s)=\frac{I_N(s)}{I_N(1)},\quad 0\le s\le1,
\]
with the positive real square root.  Let $s_{N,q}$ be the unique solution
of $R_N(s)=q$ in $(0,1)$.  These quantities also satisfy
\begin{equation}\label{beta-eq-hyper-beta}
R_N(s)=\mathrm I_{s^N}(1/N,1/2),
\end{equation}
where $\mathrm I_x(a,b)=B_x(a,b)/B(a,b)$ is the regularized incomplete
beta function.

\subsection{The primitive component and rational division}

For $N\ge3$ consider $H_N:y^2=1-x^N$.  Its holomorphic differentials are
$\omega_k=x^{k-1}dx/y$, $1\le k\le\lfloor(N-1)/2\rfloor$.  Under
$T(x,y)=(\zeta_Nx,y)$ the eigenvalue is $\zeta_N^k$, and the
characteristic polynomial on rational homology is
\[
\chi_N(X)=
\begin{cases}(X^N-1)/(X-1),&N\text{ odd},\\
(X^N-1)/(X^2-1),&N\text{ even}.
\end{cases}
\]
The rational idempotent selecting the factor $\Phi_N$ gives an algebraic
quotient $\pi:J_N\to A_N$, after clearing denominators, with
\[
H_1(A_N,\mathbb Q)\simeq\mathbb Q(\zeta_N),\qquad
\dim A_N=\varphi(N)/2.
\]
This is the usual rational-idempotent construction
\cite{beta-carocca}.  For $(k,N)=1$ and $1\le k<N/2$, there are
differentials $\eta_k$ on $A_N$ with $\pi^*\eta_k=\omega_k$.
No simplicity of $A_N$ is assumed.  Lemma~\ref{beta-lem-affine-log}
applies to $\eta_1$.

Let $O=(0,1)$ and $P(s)=(s,\sqrt{1-s^N})$.  The actual logarithm vector
$V(s)$ along the real path satisfies
\[
\exp_{A_N}V(s)=\pi([P(s)-O]),\qquad
\eta_k(V(s))=I_{N,k}(s):=\int_0^s\frac{t^{k-1}dt}{\sqrt{1-t^N}}.
\]

\begin{theorem}[All rational division ratios]\label{beta-thm-rational-division}
For every $q\in\mathbb Q\cap(0,1)$,
\[
s_{N,q}\in\overline{\mathbb Q}
\quad\Longleftrightarrow\quad N\in\{1,2,3,4,6\}.
\]
More generally, if $\varphi(N)\ge4$ and $0\le a<b\le1$ are algebraic,
the interior solution of
$\int_a^s(1-t^N)^{-1/2}dt=q\int_a^b(1-t^N)^{-1/2}dt$
is transcendental for every rational $0<q<1$.
\end{theorem}

\begin{proof}
Suppose $\varphi(N)\ge4$ and $q=a_0/b_0$ with integers
$0<a_0<b_0$.  If $s$ were algebraic, the vector
$U=b_0V(s)-a_0V(1)$ would exponentiate to an algebraic point, and
$\eta_1(U)=0$.  The lemma gives $U=0$.
There is a primitive index $1<k<N/2$.  With
$w(t)=(1-t^N)^{-1/2}$, the weighted mean
\[
M_k(s)=\frac{\int_0^s t^{k-1}w(t)dt}{\int_0^s w(t)dt}
\]
has derivative
\[
M_k'(s)=\frac{w(s)}{I_N(s)^2}
\int_0^s(s^{k-1}-t^{k-1})w(t)dt>0.
\]
Thus $\eta_k(U)=b_0I_N(s)(M_k(s)-M_k(1))<0$, a contradiction.
The same argument with base point $P(a)$ proves the truncated-interval
statement.

For $N=1$, $s_{N,q}=2q-q^2$.  For $N=2$,
$s_{N,q}=\sin(\pi q/2)$ is algebraic by a root-of-unity expression.
For $N=3,4$ the curves are elliptic, and $[P(1)-O]$ is torsion.
One uniform verification, also valid for later use, is to put
$O'=(0,-1)$ and $W_r=(r,0)$, $r^N=1$, and use
\begin{equation}\label{beta-eq-endpoint-divisors}
\operatorname{div}\frac{y-1}{y+1}=N(O-O'),\qquad
\operatorname{div}\frac{x-r}{x}=2W_r-O-O'.
\end{equation}
Hence $2N[W_r-O]=0$.  A rational multiple of a torsion logarithm on an
elliptic curve exponentiates to an algebraic torsion point, proving the
claim for $N=3,4$.  For $N=6$ use the finite map
$(x,y)\mapsto(u,v)=(x^2,xy)$ to $v^2=u-u^4$, for which
$du/v=2dx/y$ and both endpoint images are branch points.  The same
elliptic argument and algebraicity of finite fibers finish the proof.
\end{proof}

The primitive-component argument subsumes the earlier prime-leaf
argument using the absolute simplicity of the Jacobian of
$y^2=x-x^{p+1}$ for primes $p\ge5$.  In that proof the increasing
weight was $t^2$.  Here irreducibility of the rational cyclotomic action
replaces simplicity and includes $N=8$ and all composite exponents with
$\varphi(N)\ge4$.

\subsection{Geometric division and its boundary exceptions}

For $C_n=\{|z^n-1|=1\}$, a half-leaf has polar equation
$r^n=2\cos(n\theta)$ and length from the origin to radius $r$ equal to
\begin{equation}\label{beta-eq-leaf-length}
\ell(r)=2^{1/n}I_{2n}(r/2^{1/n}).
\end{equation}
This is also the geometric setting of \cite[Section~3.4 and
Remark~5.7]{beta-tamiozzo}.  For every $n\ge4$, every interior rational
division of a half-leaf has transcendental radius.  Each Cartesian
coordinate is separately transcendental: if one were algebraic, the
equation
\[
(X^2+Y^2)^n-2\operatorname{Re}(X+iY)^n=0
\]
would be a monic polynomial equation for the other, making the radius
algebraic too.  The two coordinates are not algebraically independent,
since they satisfy this curve equation.

For a traversal starting at the origin and covering all $2n$ half-leaves,
the $j$th point in an $M$-fold equal-length division is a half-leaf
boundary precisely when $M\mid2nj$.  These boundary points are algebraic
origins or tips; all other division points have both coordinates
transcendental.  Consequently all division points are algebraic exactly
when $M\mid2n$.  Repeated visits to the origin count as parameter times,
not as different geometric points.

\subsection{Algebraic ratios and exceptional torsion}

Equation~\eqref{beta-eq-endpoint-divisors} and
\eqref{beta-eq-ratio-torsion} give, for algebraic $0<s<1$,
\begin{equation}\label{beta-eq-projected-torsion}
R_N(s)\in\overline{\mathbb Q}
\quad\Longleftrightarrow\quad
\pi([P(s)-O])\text{ is torsion}.
\end{equation}
It is the projected point that is torsion; the entire Jacobian point
need not be torsion for a general composite $N$.

\begin{theorem}[Algebraic ratios]\label{beta-thm-algebraic-ratios}
\begin{enumerate}
\item If $N=8$, or $N=p$ or $2p$ for a prime $p\ge5$, then $R_N(s)$ is
transcendental at every algebraic $s\in(0,1)$.  Equivalently,
$s_{N,q}$ is transcendental for every algebraic $q\in(0,1)$.
\item If $\varphi(N)\ge4$, the set of algebraic $s\in(0,1)$ for which
$R_N(s)$ is algebraic is finite.  At any such exception, $q=R_N(s)$
satisfies
\[
\mathbb Q(q)=\mathbb Q(\zeta_N+\zeta_N^{-1}),\qquad
\deg q=\varphi(N)/2,
\]
and all its conjugates are distinct and lie in $(0,1)$.
\end{enumerate}
\end{theorem}

\begin{proof}
For $N=p$, the full Jacobian is the primitive component.  The change
$(X,Y)=(-x,y)$ gives $Y^2=X^p+1$, and
$\operatorname{div}(y-1)=p(O-\infty)$ shows that the torsion condition
uses the same base point as the known torsion-packet classification
\cite[Theorem~6.2.21]{beta-arul}.  In the interval $-1<X<0$ there are
no torsion points: the branch points and $X=0$ are absent, and the
additional $p=5$ orbit has $X^5=4$.

For $N=2p$ use instead $y^2=x-x^{p+1}$ and
$P(s)=(s^2,s\sqrt{1-s^{2p}})$.  Pullback of $dx/y$ is
$2ds/\sqrt{1-s^{2p}}$, and the full homology has polynomial
$\Phi_{2p}$.  Both reference endpoints are branch points.  The map
\[
X=-1/x,\qquad Y=i\,y/x^{(p+1)/2}
\]
sends the curve to $Y^2=X^p+1$, sends the base point to infinity, and
places $P(s)$ at $X<-1$.  The same torsion list excludes it.

For $N=8$ use $y^2=x-x^5$ and
$P(s)=(s^2,s\sqrt{1-s^8})$.  The automorphism
$(x,y)\mapsto(ix,\zeta_8y)$ acts with characteristic polynomial
$\Phi_8$ on the full genus-two Jacobian.  If $R_8(s)$ were algebraic,
$[P(s)-(0,0)]$ would be torsion.  Put
$X=\zeta_8x$, $Y=\delta y$ with $\delta^2=\zeta_8$; the curve becomes
$Y^2=X^5+X$.  By \cite[Theorem~(0.3)]{beta-boxall-grant}, its torsion
packet based at infinity consists of infinity and points with $X$ a
root of $X^5+X$ or $X^4\pm4X^2+1$.  The base-point change is torsion.
For $0<x<1$ these polynomials evaluate to
$\zeta_8x(1-x^4)$ and $1-x^4\pm4ix^2$, all nonzero.  This proves (1).

For (2), Lemma~\ref{beta-lem-affine-log} gives
$V(s)=da(V(1))$ for $a\in\mathbb Q(\zeta_N)$.  Since the first
eigenvalue $q$ is real, it belongs to the maximal real subfield.
For each primitive $1\le k<N/2$, its conjugate is
\[
\sigma_k(q)=\frac{I_{N,k}(s)}{I_{N,k}(1)}\in(0,1).
\]
These values strictly decrease with $k$.  Indeed, for $k<l$ the
difference of cross products is
\[
\int_0^s\int_s^1
\frac{t^{k-1}r^{l-1}-t^{l-1}r^{k-1}}
{\sqrt{1-t^N}\sqrt{1-r^N}}\,dr\,dt>0.
\]
Thus $q$ has full degree in the real cyclotomic field.
Finally, the curve $Z=\pi([H_N-O])$ contains zero and generates $A_N$.
It cannot be a torsion translate of an elliptic subvariety when
$\dim A_N\ge2$.  The Manin--Mumford theorem \cite{beta-roessler}
makes its torsion intersection finite.  The nonconstant map from $H_N$
to $Z$ has finite fibers, proving finiteness of the exceptional $s$.
\end{proof}

For example, $q=(\sqrt5-1)/2$ and $q=1/\sqrt2$ have negative conjugates.
The theorem excludes algebraic $s_{N,q}$ for $\varphi(N)\ge4$.
For $N=3,4,6$ an algebraic real ratio would belong to the real part of
$\mathbb Q(\zeta_N)$, which is $\mathbb Q$.  For $N=2$ the conclusion
follows from Gelfond--Schneider applied to $e^{i\pi q}$.  Thus for each
of these two probabilities, $s_{N,q}$ is transcendental for every $N\ge2$.
For $N=1$ the inverse remains $2q-q^2$ and is algebraic.

\section{Local nonvanishing and uniform degree bounds}
\label{beta-sec-local}

The finite exceptional set in Theorem~\ref{beta-thm-algebraic-ratios}
does not by itself decide an individual composite-exponent input.  A
local logarithm supplies a different test.  We normalize valuations by
$v_\ell(\ell)=1$.

\begin{lemma}[A ramified local test]\label{beta-lem-local}
Let $N\ge3$, let $\ell\ge5$ be prime with $\ell\nmid N$, and let an
algebraic $s\in(0,1)$ have an $\ell$-adic embedding for which it is a
unit and
\[
h=v_\ell(1-s^N)>2/(\ell-1).
\]
Then $R_N(s)$ is transcendental.  The local field containing $s$ may
be ramified.
\end{lemma}

\begin{proof}
Hensel's lemma gives $r^N=1$ with the same residue as $s$.  Set
$z=\sqrt{1-s^N}$, so $v_\ell(z)=h/2$.  The curve $H_N$ has good
reduction here, and at $W_r=(r,0)$ the local parameter $y$ gives
$x=r(1-y^2)^{1/N}$.  The local abelian integral is
\begin{equation}\label{beta-eq-local-log}
\mathcal L=\int_{W_r}^{P(s)}\frac{dx}{y}
=-\frac{2r}{N}\sum_{j\ge0}
\frac{(1-1/N)_j}{j!(2j+1)}z^{2j+1}.
\end{equation}
The Abel--Jacobi map is integral on the residue disk, and
$v_\ell(z)>1/(\ell-1)$ places the evaluation in the small convergence
domain of the formal logarithm.  Differentiation and termwise
integration identify \eqref{beta-eq-local-log} with its logarithm
coordinate; functoriality is the standard one for $\ell$-adic abelian
logarithms \cite{beta-zarhin}.

The coefficients $(1-1/N)_j/j!$ belong to $\mathbb Z_\ell$.
If $k=v_\ell(2j+1)$, then
$\ell^k\le2j+1$ and $\ell^k\ge1+k(\ell-1)$, so
$v_\ell(2j+1)\le2j/(\ell-1)$.  Relative to the first term, each
subsequent term has valuation at least
\[
jh-v_\ell(2j+1)\ge j\bigl(h-2/(\ell-1)\bigr)>0.
\]
These differences tend to infinity.  The first term uniquely has the
smallest valuation, and $v_\ell(\mathcal L)=h/2<\infty$.
If the projected point in \eqref{beta-eq-projected-torsion} were torsion,
then so would be $\pi([P(s)-W_r])$, by
\eqref{beta-eq-endpoint-divisors}.  Its logarithm is zero.  Since
$dx/y$ is pulled back from $A_N$, this would force $\mathcal L=0$,
a contradiction.
\end{proof}

What transfers between the complex and local arguments is the torsion
property of an algebraic point.  The real integral has not been
identified with a $p$-adic numerical value.

\begin{theorem}[All rational radii]\label{beta-thm-rational-radii}
For every integer $N\ge3$ and every rational $s\in(0,1)$, $R_N(s)$ is
transcendental.  If $N=2^hM$, $M\ge3$, and $s^{2^h}\in\mathbb Q\cap(0,1)$,
the same conclusion holds.
\end{theorem}

\begin{proof}
Write $s=a/b$ in lowest terms.  Bang--Zsigmondy supplies a primitive
prime divisor $\ell$ of $b^N-a^N$ for $N\ge3$, except at
$(b,a,N)=(2,1,6)$; see \cite{beta-voutier} for the classical input.
The residue of $a/b$ has order $N$, hence $N\mid\ell-1$,
$\ell\ge5$, and $\ell\nmid Nab$.  Now
$v_\ell(1-s^N)\ge1>2/(\ell-1)$, so
Lemma~\ref{beta-lem-local} applies.  At the exceptional triple use
$\ell=7$, since $2^6-1=63$ and the required valuation is one.
For the extension put $s^{2^h}=a/b$ and choose the same good prime for
$b^M-a^M$.  The roots of $X^{2^h}-a/b$ lie in an unramified extension
at the odd prime $\ell$, and $1-s^N=1-(a/b)^M$ has positive integral
valuation.  Apply the local lemma again.
\end{proof}

This theorem evaluates rational positions; Theorem~\ref{beta-thm-rational-division}
instead inverts rational ratios.  Neither direction should be substituted
for the other.  Some illustrative transcendental values are
\[
\begin{aligned}
R_{12}(1/2)&=0.45001710909980803611\ldots,\\
R_{12}(2^{-1/4})&=\mathrm I_{1/8}(1/12,1/2)\\
&=0.76065803913457167869\ldots.
\end{aligned}
\]

\begin{theorem}[Height-independent bounded-degree criterion]
\label{beta-thm-degree}
For $D\ge1$ set
\[
C(D)=\max\{3^D,\;4000(D\log(3D))^{12}\}.
\]
If $s\in\overline{\mathbb Q}\cap(0,1)$, $[\mathbb Q(s):\mathbb Q]\le D$,
and $N>C(D)$, then $R_N(s)$ is transcendental.  In particular,
$D=2$ and $N>20\,000\,000\,000$ suffice.
\end{theorem}

\begin{proof}
The precise primitive-divisor input is Voutier's
\cite[Theorem~2(i)]{beta-voutier}: for coprime nonzero algebraic
integers $A,B$ with $A/B$ not a root of unity and of degree $d$, a
prime ideal dividing $B^N-A^N$ but no earlier $B^j-A^j$ exists when
\[
N>\max\{2(2^d-1),4000(d\log(3d))^{12}\}.
\]
Its bound uses the degree of the ratio, not that of an auxiliary field.

To handle denominators, in $K=\mathbb Q(s)$ write
$(s)=\mathfrak a\mathfrak b^{-1}$ with coprime integral ideals.
Choose $h$ with $\mathfrak b^h=(b_0)$, adjoin $B$ with $B^h=b_0$, and
put $A=sB$.  In that extension, $(B)=\mathfrak b$ and $(A)=\mathfrak a$
after extension of ideals, since the fractional ideal group is
torsion-free.  Thus $A,B$ are coprime algebraic integers and $A/B=s$.
Apply Voutier and contract the resulting prime to a prime
$\mathfrak p$ of $K$.  The residue of $s$ still has exact order $N$,
because the residue-field inclusion is injective.

Let $\mathfrak p\mid\ell$ have ramification index $e$ and residue
degree $f$.  Then
\[
N\mid\ell^f-1,\qquad \ell\nmid N,\qquad
v_\ell(1-s^N)\ge1/e,\qquad ef\le d\le D.
\]
If $\ell\le2e+1$, then
$N+1\le\ell^f\le(2e+1)^f\le3^{ef}\le3^D$, a contradiction.
Thus $\ell>2e+1$ and $1/e>2/(\ell-1)$; the local lemma applies.
The initial Voutier threshold is covered because
$2(2^d-1)\le3^d\le3^D$.  Finally,
$\log6<9/5$ gives
$4000(2\log6)^{12}<4000(18/5)^{12}<20\,000\,000\,000$.
\end{proof}

The bound is explicit but is not claimed to be optimal.  For a fixed
algebraic $s$ it confines possible exceptional exponents to a finite
interval independently of its height.  Conversely, any exceptional
sequence with $N\to\infty$ must have input degrees tending to infinity.

A small ramified example does not need the large bound.  Let
$s_0=3-\sqrt7$.  For $N=6m$ with $7\nmid m$, expansion at $\sqrt7$
shows
\[
v_7(1-s_0^N)=1/2>1/3=2/(7-1):
\]
the constant term $1-3^N$ has valuation at least one, the linear term
$N3^{N-1}\sqrt7$ has valuation exactly $1/2$, and all remaining terms
have valuation at least one.  Therefore every $R_{6m}(s_0)$ in this
range is transcendental.  For example
\[
R_{12}(3-\sqrt7)=0.31883299555500762110\ldots.
\]

\section{Second-kind beta integrals and polynomial weights}
\label{beta-sec-secondkind}

For positive rational parameters write
$B_x(a,b)=\int_0^xt^{a-1}(1-t)^{b-1}dt$ and $B(a,b)=B_1(a,b)$.
The restriction that a beta differential be holomorphic can be removed
by preserving its pole data in a generalized Jacobian.

\begin{theorem}[Affine independence of beta periods]
\label{beta-thm-beta-independent}
Let $p,r$ be distinct primes with $\max(p,r)\ge5$, and let
$a=i/p$, $b=j/r$, where $1\le i<p$, $1\le j<r$.  For every algebraic
$x\in(0,1)$, the three numbers
\[
1,\quad B(a,b),\quad B_x(a,b)
\]
are linearly independent over $\overline{\mathbb Q}$.  There is no
restriction $a+b<1$.
\end{theorem}

\begin{proof}
On the smooth projective model of $C:U^p+V^r=1$, with its unique point
$\infty$ at infinity, consider
\[
\eta=U^{i-1}V^{j-r}dU.
\]
It is regular at all finite points.  At infinity its order is
\begin{equation}\label{beta-eq-beta-pole}
\operatorname{ord}_{\infty}\eta=pr-ri-pj-1.
\end{equation}
Since the prime denominators are distinct, $a+b\ne1$.  Thus $\eta$ is
holomorphic if $a+b<1$, and otherwise has a pole of order at least two.
Its residue is zero by the residue theorem, since it has no other pole.

Choose $M\ge2$ at least its pole order.  The generalized Jacobian for
the modulus $M\infty$ is an extension
\[
0\longrightarrow\mathbb G_a^{M-1}\longrightarrow G
\xrightarrow{\rho}J\longrightarrow0.
\]
There is an algebraic Abel map $j_O:C\setminus\{\infty\}\to G$ based
at $O=(0,1)$ and an invariant differential $\widetilde\eta$ with
$j_O^*\widetilde\eta=\eta$.
Here is a local check of the differential condition.  The pinched curve
for this modulus has completed local ring
$\overline{\mathbb Q}+t^M\overline{\mathbb Q}[[t]]$.
The Rosenlicht condition is $\operatorname{Res}(f\eta)=0$ for every
element $f$ of this ring.  The constant term contributes zero by the
zero residue, and the remaining terms remove the pole.  Thus $\eta$
belongs to its dualizing differentials and corresponds to an invariant
differential of $G$.  This standard algebraic construction, including
the vector extension and Abel-map interpretation, is described in
\cite{beta-zannier,beta-agostini}.

The rotation $T(U,V)=(\zeta_pU,\zeta_rV)$ preserves the modulus.  The
holomorphic characters $u/p+v/r<1$, together with their conjugates,
give every primitive $pr$th root exactly once on $H_1(J,\mathbb Q)$.
Its characteristic polynomial is $\Phi_{pr}$, and
$T^*\widetilde\eta=\xi\widetilde\eta$ with
$\xi=\zeta_p^i\zeta_r^j$ primitive.
Unlike the holomorphic case, nonzeroness of a second-kind differential
does not suffice to assert a nonzero period.  We exhibit one.  Let
$\ell$ be the positive real arc from $O=(0,1)$ to $W=(1,0)$, and put
$S(U,V)=(\zeta_pU,V)$, $R(U,V)=(U,\zeta_rV)$.  Then
$(1-R_*)(1-S_*)\ell$ is closed and
\begin{equation}\label{beta-eq-beta-period}
\int_{(1-R_*)(1-S_*)\ell}\eta
=(1-\zeta_r^j)(1-\zeta_p^i)\frac{B(a,b)}p\ne0.
\end{equation}
Its image in $G$ is a period, so
Lemma~\ref{beta-lem-affine-log} applies.

For $P_x=(x^{1/p},(1-x)^{1/r})$, let $u_x$ and $v$ be the logarithms
of $j_O(P_x)$ and $j_O(W)$ along these real paths.  Then
\[
\widetilde\eta(u_x)=B_x(a,b)/p,\qquad
\widetilde\eta(v)=B(a,b)/p.
\]
The divisors of $V-1$ and $U-1$ are respectively
$p(O-\infty)$ and $r(W-\infty)$, so $[W-O]$ is torsion in $J$.
Its path logarithm $d\rho(v)$ is nonzero: evaluate it on the
holomorphic differential with indices $i=j=1$, whose real path integral
is strictly positive and whose holomorphicity follows from
$1/p+1/r<1$.

Suppose $c_0+c_1B(a,b)+c_2B_x(a,b)=0$ with algebraic coefficients.
If $c_2\ne0$, the second part of the lemma forces $[P_x-O]$ to be
torsion in $J$, hence $[P_x-\infty]$ is torsion.  For distinct odd
primes, $(X,Y)=(-U,V)$ identifies the curve with $Y^r=X^p+1$.
The torsion classification \cite[Theorem~6.2.21]{beta-arul} excludes
finite torsion points with both coordinates nonzero.  If one prime is
two, reorder the variables and write $Y^2=X^p+1$ with
$X\in(-1,0)$.  The only extra case, $p=5$, has $X^5=4$,
also impossible on this interval.  Thus $c_2=0$.
If $c_1\ne0$, the first part of the lemma applied to $v$ would force
$d\rho(v)=0$.  Hence $c_1=0$ and then $c_0=0$.
\end{proof}

Two limitations of the proof are material.  It only forces torsion
after projection to the ordinary Jacobian; additive logarithm
directions in $G$ need not vanish.  Also, the cyclotomic field was not
assumed to embed in the entire endomorphism algebra of $G$: the proof
uses a polynomial in the actual automorphism $T$ and then descends it.

\subsection{An exact classifier for all polynomial weights}

Let $W(t)=\sum_k w_kt^k\in\overline{\mathbb Q}[t]$ and define
\begin{equation}\label{beta-eq-weight-coefficient}
c_W=\sum_k w_k\frac{(a)_k}{(a+b)_k},\qquad
D(x)=x^a(1-x)^b.
\end{equation}
There is a polynomial $P_W\in\overline{\mathbb Q}[x]$ with
\begin{equation}\label{beta-eq-weight-reduction}
\int_0^x W(t)t^{a-1}(1-t)^{b-1}dt
=c_WB_x(a,b)+D(x)P_W(x).
\end{equation}
To compute it, start with $c_0=1$, $P_0=0$ for $W=1$, and use
\[
c_{k+1}=\frac{a+k}{a+b+k}c_k,\qquad
P_{k+1}=\frac{a+k}{a+b+k}P_k-\frac{x^k}{a+b+k}
\]
for $W=t^{k+1}$, then combine linearly.  These are beta recurrence
identities \cite{beta-dlmf}; differentiation gives the useful exact
check
\begin{equation}\label{beta-eq-weight-ode}
x(1-x)P_W'+[a-(a+b)x]P_W=W-c_W.
\end{equation}
Because $D(0)=D(1)=0$, the complete weighted integral is $c_WB(a,b)$.

\begin{corollary}[Polynomial-weight classification]
\label{beta-cor-weight}
Under the hypotheses of Theorem~\ref{beta-thm-beta-independent}, and
for every algebraic $x\in(0,1)$,
\[
\int_0^xW(t)t^{a-1}(1-t)^{b-1}dt\in\overline{\mathbb Q}
\quad\Longleftrightarrow\quad c_W=0.
\]
If $c_W\ne0$, its quotient by the complete weighted integral is
transcendental.  If $W$ has real algebraic coefficients and is positive
on $(0,1)$, every interior algebraic-probability quantile of the
resulting density is transcendental.  In particular this holds for
every $\operatorname{Beta}(n+a,m+b)$ with integers $n,m\ge0$.
\end{corollary}

\begin{proof}
At algebraic $x$, the second term in
\eqref{beta-eq-weight-reduction} is algebraic.  If $c_W=0$, this proves
algebraicity.  If $c_W\ne0$, an algebraic incomplete integral would
contradict linear independence of $1,B_x$; an algebraic normalized
integral would contradict that of $1,B,B_x$.  Positivity gives
$c_WB(a,b)>0$, and strict increase of the cumulative distribution
provides its unique interior quantiles.  An algebraic quantile would
contradict the normalized-value statement.  Finally choose
$W=t^n(1-t)^m$.
\end{proof}

Thus, in the space of all polynomial weights of degree at most any fixed
$D\ge0$, the algebraic incomplete integrals correspond exactly to the
hyperplane $c_W=0$.
The exceptional hyperplane is nonempty.  For $a=2/3$, $b=2/5$,
$W(t)=8t-5$ has $c_W=0$ and
\[
\int_0^x(8t-5)t^{-1/3}(1-t)^{-3/5}dt
=-\frac{15}{2}x^{2/3}(1-x)^{2/5}.
\]
Discarding the nonzero-coefficient condition would therefore give a
false statement.

For the positive weight $W=t(1-t)$ at these same parameters,
$c_W=15/124$ and $P_W/c_W=4x-3/2$.  Consequently
\begin{equation}\label{beta-eq-new-median}
\begin{aligned}
B(2/3,2/5)\,\mathrm I_x(5/3,7/5)
&=B_x(2/3,2/5)\\
&\quad +(4x-3/2)x^{2/3}(1-x)^{2/5}.
\end{aligned}
\end{equation}
The median of $\operatorname{Beta}(5/3,7/5)$ is therefore
transcendental; its approximate value is
\[
0.55398928402851700562\ldots.
\]
The base parameters sum to $16/15>1$, and the underlying differential
has a double pole.  This example specifically uses the second-kind
extension rather than a holomorphic argument.

The excluded prime pair $\{2,3\}$ is elliptic and has algebraic
rational-probability division points.  Equal prime denominators also
cannot be admitted indiscriminately, since symmetric distributions
have the algebraic median $1/2$.  Neither algebraic independence of
different quantiles nor a complete arbitrary-composite-denominator
extension is concluded here.

\subsection{Numerical records and their role}

All decimals in these sections denote the exact integral or inverse
value specified before the decimal expansion.  A finite decimal itself
is rational.  The previous computation records used exact rational
intervals, integer root bounds, and the positive expansions
\[
I_{N,k}(s)=\sum_{j\ge0}
\frac{\binom{2j}{j}s^{Nj+k}}{4^j(Nj+k)},\qquad
B_x(a,b)=x^a\sum_{j\ge0}\frac{(1-b)_jx^j}{j!(a+j)}.
\]
For $0<a,b<1$, successive terms in the latter series have ratio less
than $x$, so the first omitted term divided by $1-x$ bounds the tail.
Complete integrals are evaluated by the two half-interval identity
$B(a,b)=B_{1/2}(a,b)+B_{1/2}(b,a)$.  These methods supply reproducible
sign tests at rational brackets for quantiles.  Their role is numerical
certification; the transcendence conclusions use the proofs above and
the explicitly cited established theorems, not a finite computation.

\section{Arithmetic classifications of statistical quantiles}
\label{beta-stats-sec-quantiles}

All probabilities in this section refer to exact real numbers, rather than
rounded entries in statistical tables. Write $\overline{\mathbb Q}$ for the
field of algebraic numbers. The transcendence input is Baker's theorem:
an algebraic linear combination of fixed logarithms of nonzero algebraic
numbers is either zero or transcendental \cite{beta-stats-baker67}.
We also use Hermite--Lindemann and Gelfond--Schneider in their classical
forms; the latter says that every value of $\alpha^\beta$ is
transcendental when $\alpha\ne0,1$ is algebraic and $\beta$ is algebraic
irrational \cite{beta-stats-baker66}. These are established inputs, not
new results. The use of logarithmic forms for rational-function integrals
also has a classical literature \cite{beta-stats-vdp}.

\subsection{Student distributions: a complete integer-degree classification}

Let $Q_\nu(p)$ denote the quantile of the standard Student distribution
with positive integer degrees of freedom $\nu$.

\begin{theorem}[Student quantiles]
\label{beta-stats-thm-student}
For $p\in\overline{\mathbb Q}\cap(0,1)$:
\begin{enumerate}
\item if $\nu$ is even, $Q_\nu(p)$ is algebraic;
\item if $\nu\ge3$ is odd, $Q_\nu(p)$ is transcendental unless
      $p=1/2$, when it equals zero;
\item if $\nu=1$, $Q_1(p)$ is algebraic if and only if $p$ is rational.
\end{enumerate}
\end{theorem}

\begin{proof}
For $\nu=2m+1$, $m\ge1$, put $u=x/\sqrt\nu$. Integration by parts gives
\begin{equation}
\label{beta-stats-eq-odd-student}
F_{2m+1}(x)=\frac12+\frac{\arctan u+R_m(u)}\pi,
\qquad
R_m(u)=\sum_{k=1}^{m}
\frac{(2k-2)!!}{(2k-1)!!}\frac{u}{(1+u^2)^k}.
\end{equation}
Every coefficient in this sum is positive, so $R_m(u)\ne0$ for real
$u\ne0$. If both $x\ne0$ and $p=F_{2m+1}(x)$ were algebraic, then,
with compatible fixed logarithms,
\[
\log\frac{1+iu}{1-iu}-(2p-1)\log(-1)=-2iR_m(u),
\qquad \log(-1)=i\pi.
\]
This contradicts Baker's theorem, since the right side is nonzero
algebraic.

For $\nu=2m$, set $y=x/\sqrt{2m+x^2}$. The transformed density is
$c_m(1-y^2)^{m-1}$, where
$c_m=(2m)!/(4^m m!(m-1)!)\in\mathbb Q$. Consequently
\[
F_{2m}(x)=\frac12+c_m\sum_{k=0}^{m-1}
(-1)^k\binom{m-1}{k}\frac{y^{2k+1}}{2k+1}.
\]
Setting this nonconstant polynomial equal to algebraic $p$ makes $y$,
and hence $x=\sqrt{2m}\,y/\sqrt{1-y^2}$, algebraic.
Finally $Q_1(p)=\tan(\pi(p-1/2))$. Rational $p$ gives an algebraic value
by roots of unity. If $p$ were algebraic irrational and $Q_1(p)$
algebraic, then
\[
\frac{1+iQ_1(p)}{1-iQ_1(p)}=(-1)^{2p-1}
\]
would contradict Gelfond--Schneider, using the value defined by
$\log(-1)=i\pi$.
\end{proof}

For example, $Q_3(39/40)=3.182446305\ldots$ is transcendental, whereas
$Q_2(39/40)=19\sqrt{2/39}$ is algebraic. At this particular rational
probability, Hermite--Lindemann alone proves the first assertion:
if $u=Q_3(39/40)/\sqrt3$ were algebraic and
$r=u/(1+u^2)>0$, then $\arctan u+r=19\pi/40$ would imply
$e^{40ir}=-((1+iu)/(1-iu))^{-20}$, a contradiction.

\subsection{Half-integer beta quantiles and three possible exceptions}

For integers $n,m\ge0$, let $F_{n,m}$ be the distribution function of
$\operatorname{Beta}(n+1/2,m+1/2)$. For $N=n+m\ge1$, define
\[
K_{n,m}=\frac{\pi}{\mathrm B(n+1/2,m+1/2)}
=\frac{4^{n+m}n!m!(n+m)!}{(2n)!(2m)!}\in\mathbb Q.
\]
There is a unique $P_{n,m}\in\mathbb Q[x]$, of degree at most $N-1$,
such that
\begin{align}
2x(1-x)P'_{n,m}+(1-2x)P_{n,m}
 &=2K_{n,m}x^n(1-x)^m-2, \label{beta-stats-eq-half-correction}\\
\pi F_{n,m}(x)
 &=2\arcsin\sqrt{x}+\sqrt{x(1-x)}P_{n,m}(x).
 \label{beta-stats-eq-half-primitive}
\end{align}
For completeness, the operator on the left of
\eqref{beta-stats-eq-half-correction} is injective on polynomials of
degree at most $N-1$: a leading term $a_dx^d$ is sent to a polynomial
with leading term $-(2d+2)a_dx^{d+1}$. Its image has codimension one
among polynomials of degree at most $N$. That image is exactly the
kernel of integration against $1/(2\sqrt{x(1-x)})$, because the
integrand is the derivative of $\sqrt{x(1-x)}P(x)$. The right side
of \eqref{beta-stats-eq-half-correction} has zero such integral.
This proves existence, uniqueness, and rational coefficients;
differentiation and the value at zero prove
\eqref{beta-stats-eq-half-primitive}.

\begin{theorem}[Three-candidate theorem]
\label{beta-stats-thm-three}
Let $n,m\ge0$ and $n+m\ge1$. If $p,x\in\overline{\mathbb Q}\cap(0,1)$
and $F_{n,m}(x)=p$, then
\[
(p,x)\in\{(1/3,1/4),(1/2,1/2),(2/3,3/4)\}.
\]
A listed pair occurs if and only if $P_{n,m}(x)=0$. For $n,m\ge1$
at most one listed pair occurs; if $n=0$ or $m=0$, none occurs.
\end{theorem}

\begin{proof}
Put $\theta=\arcsin\sqrt{x}$ and
$u=\sqrt{x/(1-x)}$. If $x$ is algebraic, then
$A=(1+iu)/(1-iu)$ is algebraic and we can fix $\log A=2i\theta$.
Equation \eqref{beta-stats-eq-half-primitive} becomes
\[
p\log(-1)-\log A=i\sqrt{x(1-x)}P_{n,m}(x).
\]
Baker's theorem forces $P_{n,m}(x)=0$. Hence $p=2\theta/\pi$ and
$A=(-1)^p$; Gelfond--Schneider forces $p\in\mathbb Q$.

It remains to control the roots of $P_{n,m}$. Subtract the arcsine
distribution function $F_{0,0}(x)=2\arcsin\sqrt{x}/\pi$ and write
$H=F_{n,m}-F_{0,0}$. Then
\[
H'(x)=\frac{K_{n,m}x^n(1-x)^m-1}{\pi\sqrt{x(1-x)}}.
\]
For $n,m\ge1$, the numerator has signs $-,+,-$: the function
$K_{n,m}x^n(1-x)^m$ is strictly unimodal, vanishes at both ends,
and has average one under the arcsine probability measure.
Since $H(0)=H(1)=0$, $H$ is negative just after zero, positive just
before one, and crosses zero exactly once in its increasing region.
Thus $P_{n,m}$ has exactly one root in $(0,1)$. When one of $n,m$
is zero, the density ratio is strictly monotone; the same derivative
argument gives no interior root.

For rational $p$, put $\zeta=e^{i\pi p}$. Then
$x=(2-\zeta-\zeta^{-1})/4$. The order of $\zeta$ is at least three,
so every conjugate of $x$ is real and lies strictly between zero and
one. Every conjugate is also a root of $P_{n,m}\in\mathbb Q[x]$.
The unique-root property therefore implies $x\in\mathbb Q$.
Now $2\cos(\pi p)=2-4x$ is both a rational number and an algebraic
integer, lying in $(-2,2)$. It is consequently $-1$, $0$, or $1$.
These give precisely the three pairs stated. Their sufficiency
follows by substituting $P_{n,m}(x)=0$ in
\eqref{beta-stats-eq-half-primitive}.
\end{proof}

This is an effective classification for each specified pair $(n,m)$:
only three rational evaluations of its correction polynomial are
required. It is not a classification of all integer pairs $(n,m)$
for which $P_{n,m}(1/4)=0$. The omitted arcsine case $n=m=0$ has
quantile $\sin^2(\pi p/2)$; for algebraic $p$ this is algebraic exactly
when $p$ is rational, by the same Gelfond--Schneider argument.

For $\operatorname{Beta}(3/2,5/2)$,
\[
P_{1,2}(x)=-\frac23(4x-1)(2x-3).
\]
Its unique algebraic--algebraic interior pair is
$F_{1,2}(1/4)=1/3$. In particular its median and its $1/40$ and
$39/40$ quantiles are transcendental.

\subsection{Posterior intervals and Fisher distributions}

In a binomial sample with $N\ge1$ observations and $s$ successes,
the Jeffreys posterior is $\operatorname{Beta}(s+1/2,N-s+1/2)$.
Theorem~\ref{beta-stats-thm-three} makes both endpoints of its exact,
unmodified $95\%$ equal-tail interval transcendental, since neither
$1/40$ nor $39/40$ is a candidate probability. Its posterior median
is $1/2$ when $s=N-s$ and transcendental otherwise. For the latter
assertion, compare the density at $x$ and $1-x$ on $(0,1/2)$ to
exclude a median of $1/2$ in the asymmetric case. A convention that
replaces an endpoint by $0$ or $1$ at boundary data changes this
conclusion for that replaced endpoint.

By contrast, Clopper--Pearson endpoints are quantiles of beta
distributions with positive integer parameters, or boundary values
$0,1$. Their distribution functions are rational-coefficient
polynomials, so all such endpoints at algebraic confidence levels
are algebraic. These definitions agree with the standard statistical
conventions described in \cite{beta-stats-ratesci}; the arithmetic
distinction expresses no claim about statistical superiority.

For a standard Fisher $F(d_1,d_2)$ variable $Y$, the transformation
$X=d_1Y/(d_1Y+d_2)$ gives
$X\sim\operatorname{Beta}(d_1/2,d_2/2)$ and preserves algebraicity
in both directions. Thus for positive integer $d_1,d_2$:
\begin{itemize}
\item If either degree is even, every quantile at algebraic probability
is algebraic. If exactly one beta parameter is half-integral, a
square-root substitution reduces its integral to a polynomial.
\item If both degrees are odd and $(d_1,d_2)\ne(1,1)$, there is at
most one algebraic quantile at algebraic probability, determined by
Theorem~\ref{beta-stats-thm-three}. In particular the $95\%$ quantile
is transcendental.
\item For $(d_1,d_2)=(1,1)$, algebraic probabilities give algebraic
quantiles exactly at rational probabilities.
\end{itemize}
For example the sole algebraic--algebraic interior pair for $F(3,5)$
is its probability-$1/3$ quantile $5/9$.

\section{Shortest beta probability intervals}
\label{beta-stats-sec-hpd}

Let $f$ and $F$ be the density and distribution function of
$\operatorname{Beta}(\alpha,\beta)$. We measure interval length in
the original probability coordinate $x$, not after a nonlinear
reparameterization. Assume
\begin{equation}
\label{beta-stats-eq-hpd-hyp}
\alpha,\beta\in\mathbb Q,\qquad \alpha,\beta>1,\qquad
\alpha+\beta\in\mathbb Z,\qquad
c\in\overline{\mathbb Q}\cap(0,1).
\end{equation}
The strict unimodality of $f$ gives a unique highest-density interval
$[a,b]$ of mass $c$, with
\[
0<a<\frac{\alpha-1}{\alpha+\beta-2}<b<1,
\qquad f(a)=f(b),\qquad F(b)-F(a)=c.
\]
Indeed its density level decreases continuously as its mass increases
from zero to one. This interval is the unique shortest interval of
that mass: at its boundary level $\lambda$, replacing any part of it
by a set outside it loses density, and shortening it strictly loses
mass.

\begin{theorem}[Arithmetic dichotomy for shortest intervals]
\label{beta-stats-thm-hpd}
Under \eqref{beta-stats-eq-hpd-hyp}, if $\alpha,\beta$ are integers,
then $a,b$, and $b-a$ are algebraic. If they are nonintegers, then
each of $a,b$, and $b-a$ is transcendental, and
\[
\operatorname{trdeg}_{\overline{\mathbb Q}}
\overline{\mathbb Q}(a,b)=1.
\]
In the noninteger case the midpoint is $1/2$ for $\alpha=\beta$,
and is transcendental for $\alpha\ne\beta$.
\end{theorem}

\subsection{A correction-to-density monotonicity lemma}

In the noninteger case write uniquely
$\alpha=n+r$, $\beta=m+1-r$, where $n,m\ge1$ and
$r\in\mathbb Q\cap(0,1)$. Set $N=n+m$ and
\[
w_0(x)=x^{r-1}(1-x)^{-r},\quad
B_0=\mathrm B(r,1-r)=\frac\pi{\sin\pi r},\quad
K=\frac{B_0}{\mathrm B(\alpha,\beta)}
=\frac{N!}{(r)_n(1-r)_m}.
\]
The beta recurrence and reflection identities used here are standard
\cite{beta-stats-dlmf}. Exactly as in the polynomial image argument
above, there is a unique $P\in\mathbb Q[x]$, $\deg P\le N-1$, with
\begin{align}
x(1-x)P'(x)+(r-x)P(x)&=Kx^n(1-x)^m-1,
\label{beta-stats-eq-hpd-correction}\\
B_0F(x)&=\int_0^x w_0(t)\,dt+D(x),\quad
D(x)=x^r(1-x)^{1-r}P(x).
\label{beta-stats-eq-hpd-primitive}
\end{align}
For that argument the relevant annihilating functional is integration
against $w_0(x)$, and
$w_0[x(1-x)P'+(r-x)P]=(x^r(1-x)^{1-r}P)'$.

\begin{lemma}
\label{beta-stats-lem-monotone}
The rational function
$R(x)=P(x)/(x^{n-1}(1-x)^{m-1})$ satisfies $R'(x)>0$ for $0<x<1$.
\end{lemma}

\begin{proof}
Expand in the unnormalized Bernstein basis,
$P(x)=\sum_{k=0}^{N-1}p_kx^k(1-x)^{N-1-k}$.
With $p_{-1}=p_N=0$, comparison of degree-$N$ coefficients in
\eqref{beta-stats-eq-hpd-correction} yields
\[
(j+r)p_j-(N-j+1-r)p_{j-1}
=K\mathbf1_{j=n}-\binom Nj\qquad(0\le j\le N).
\]
Starting with $p_0=-1/r$, forward recurrence gives $p_k<0$ for
$k<n$. Starting with $p_{N-1}=1/(1-r)$, backward recurrence gives
$p_k>0$ for $k\ge n$. The corresponding summand of $R$ is
$p_kx^j(1-x)^{1-j}$, with $j=k-n+1$, whose derivative equals
\[
p_kx^{j-1}(1-x)^{-j}(j-x).
\]
It is strictly positive both for $k<n$ ($p_k<0$, $j\le0$) and
for $k\ge n$ ($p_k>0$, $j\ge1$).
\end{proof}

Since $D(x)=(B_0/K)f(x)R(x)$, the equal-density condition implies
the decisive noncancellation
\begin{equation}
\label{beta-stats-eq-noncancel}
D(b)-D(a)=\frac{B_0}{K}f(a)[R(b)-R(a)]>0.
\end{equation}
Knowing only the algebraic mass $F(b)-F(a)$ does not make either
$F(a)$ or $F(b)$ algebraic. This is why the single-quantile theorem
alone would not suffice for shortest intervals.

\subsection{Proof of the dichotomy and its geometric consequences}

Suppose first that the noninteger parameters satisfy
\eqref{beta-stats-eq-hpd-hyp} and that both endpoints are algebraic.
Then $D(b)-D(a)$ is nonzero algebraic. Write $r=p/q$ in lowest terms
and make the positive-real substitution
$y=(x/(1-x))^{1/q}$. It gives
\[
w_0(x)\,dx=\frac{q y^{p-1}}{1+y^q}\,dy.
\]
This is a proper rational function with distinct algebraic poles.
Its integral between algebraic endpoints is therefore an algebraic
linear combination of logarithms of nonzero algebraic numbers;
logarithms are chosen continuously along the integration path.
The same is true of $cB_0$, because $\sin\pi r$ is nonzero
algebraic and $\pi=\log(-1)/i$. But
\eqref{beta-stats-eq-hpd-primitive} implies
\[
D(b)-D(a)=cB_0-\int_a^b w_0(x)\,dx,
\]
contradicting Baker's theorem and
\eqref{beta-stats-eq-noncancel}.

Choose a common positive denominator $q$ so that
$A=q(\alpha-1)$ and $B=q(\beta-1)$ are positive integers, and put
$h(X)=X^A(1-X)^B$. Equal density gives the polynomial relation
$h(a)=h(b)$. If either endpoint were algebraic, this relation would
make the other algebraic too. Thus both are transcendental; the
nonzero relation also proves that their joint transcendence degree
is exactly one.

If $w=b-a$ were algebraic, the transcendental number $a$ would
satisfy $h(a+w)-h(a)=0$. For $w>0$ the polynomial
$h(X+w)-h(X)$ is nonzero, contradicting transcendence. If the
midpoint $z=(a+b)/2$ were algebraic, the same reasoning applied to
$h(2z-X)-h(X)$ would force this polynomial to vanish identically.
Reflection would then preserve the roots $0,1$ and their
multiplicities. Necessarily $z=1/2$ and $A=B$, hence
$\alpha=\beta$. In this symmetric case uniqueness of the interval
indeed gives $z=1/2$.

For the integer case let $A=\alpha-1$, $B=\beta-1$. Now $F$ is
a rational-coefficient polynomial. The equal-density pair is
parameterized by $t>1$ as
\[
a(t)=\frac{t^A-1}{t^{A+B}-1},\qquad
b(t)=t^B a(t).
\]
To derive this, use $t=(b/a)^{1/B}$ in
$a^A(1-a)^B=b^A(1-b)^B$. The mass
$C(t)=F(b(t))-F(a(t))$ is a rational function in $\mathbb Q(t)$,
and it is nonconstant because its limits as $t\downarrow1$ and
$t\to\infty$ are zero and one. The equation $C(t)=c$ therefore
gives a nonzero algebraic-coefficient polynomial equation for $t$.
Thus $t,a,b$, and $b-a$ are algebraic, completing the proof of
Theorem~\ref{beta-stats-thm-hpd}.

\subsection{Examples and scope}

For $\operatorname{Beta}(3/2,5/2)$ and $c=19/20$, the unique
shortest interval is approximately
\[
[a,b]=[0.009556870873\ldots,\;0.770757269257\ldots].
\]
Both endpoints, its length, and its midpoint are transcendental.
This conclusion also holds for every algebraic mass in $(0,1)$.
The result is not confined to half-integer parameters: for
$\operatorname{Beta}(4/3,5/3)$ one has $r=1/3$, $n=m=1$,
$K=9$, and $P(x)=-3+(9/2)x$, so the same four conclusions hold.

With at least one success and one failure, a Jeffreys beta posterior
has transcendental shortest-interval endpoints and length at every
algebraic credibility level; its midpoint is transcendental unless
the counts are equal. A uniform-prior posterior instead has integer
parameters and algebraic endpoints and length. The theorem does
not cover zero-success or zero-failure HPD boundary cases,
U-shaped beta densities, or general rational parameters whose sum
is not an integer. No claim of priority for these particular
classifications is made here, and no conclusion about the
transcendence of a specified normal quantile or Euler's constant
follows from them.

\section{Finite Gaussian mixtures: an exact endpoint classification}
\label{beta-stats-sec-gaussian}

The logarithmic arguments above are complemented by an $E$-function
method for finite Gaussian mixtures. The normalization interval is
finite throughout this classification. This restriction is essential
to its proof and separates it from a classification of the ordinary
normal quantiles.

\subsection{Linear independence at distinct positive algebraic endpoints}

For a positive integer $d$, define
$J_d(x)=\int_0^x e^{-t^{2d}}\,dt$.

\begin{lemma}[Endpoint separation]
\label{beta-stats-lem-gaussian-separation}
If $r_1,\ldots,r_s$ are distinct positive algebraic numbers, then
\[
1,\ J_d(r_1),\ldots,J_d(r_s),\quad
e^{-r_1^{2d}},\ldots,e^{-r_s^{2d}}
\]
are linearly independent over $\overline{\mathbb Q}$.
\end{lemma}

\begin{proof}
We do not regard $J_d(z)$ itself as an $E$-function. Instead set
\[
E_{c,d}(z)=\int_0^1e^{-czt^{2d}}\,dt
=\sum_{n\ge0}\frac{(-c)^n}{2dn+1}\frac{z^n}{n!},
\qquad c\in\overline{\mathbb Q}_{>0}.
\]
It satisfies $2dzE_{c,d}'(z)+E_{c,d}(z)=e^{-cz}$. Its normalized
coefficients have exponentially bounded conjugates and common
denominators: if $Dc$ is an algebraic integer, a common denominator
through index $M$ is
$D^M\operatorname{lcm}(1,\ldots,2dM+1)$, which grows exponentially.
Thus $E_{c,d}$ and $e^{-cz}$ are $E$-functions forming, together with
$1$, a differential system whose only possible finite singularity
is zero.

For distinct $0<c_1<\cdots<c_s$, these functions are linearly
independent over $\mathbb C(z)$. Indeed, assume a nontrivial rational
relation and choose the largest $c$ appearing. Along $z=-x$,
$x\to+\infty$, the smaller exponential scales and their derivatives
are exponentially smaller than $e^{cx}$, up to polynomial factors.
If the coefficient of $E_{c,d}$ is zero, the nonzero rational
coefficient of $e^{-cz}$ would have to be exponentially small, an
impossibility. Otherwise the relation yields an $R\in\mathbb C(z)$
such that $T=e^{cz}E_{c,d}-R$ and $T'$ are exponentially small there.
Applying $\mathcal L_c=2dz\,d/dz+1-2dcz$ shows that the rational
function $1-\mathcal L_cR$ is exponentially small, hence zero.
But
\[
2dzR'+(1-2dcz)R=1
\]
has no rational solution. A nonzero finite pole away from zero
cannot survive its differentiated highest-order term. A pole of
order $m\ge1$ at zero has nonzero coefficient factor $1-2dm$.
Thus $R$ would be a polynomial, which is excluded by the highest
degree term $-2dczR$ (and $R=0$ is also excluded).

Beukers's refined Siegel--Shidlovskii theorem
\cite[Theorem~1.3]{beta-stats-beukers} lifts any algebraic linear
relation among these values at the regular algebraic point $z=1$
to a polynomial-coefficient linear relation among the functions.
The established functional independence rules that out. Finally use
$c_i=r_i^{2d}$ and $J_d(r_i)=r_iE_{r_i^{2d},d}(1)$.
\end{proof}

\subsection{A finite classification algorithm}

Let $K$ be a subfield of the real algebraic numbers. Choose
$a<b$, positive weights $w_j$, positive scales $\lambda_j$, and
centers $\mu_j$, all in $K$, for $1\le j\le N$. Merge components
with identical $(\lambda_j,\mu_j)$. On $[a,b]$ normalize the positive
density
\[
f_d(x)=\sum_{j=1}^{N}w_j
e^{-[\lambda_j(x-\mu_j)]^{2d}},\qquad
S_d(x)=\sum_{j=1}^{N}\frac{w_j}{\lambda_j}
J_d(\lambda_j(x-\mu_j)).
\]
Its distribution function is
$C_d(x)=(S_d(x)-S_d(a))/(S_d(b)-S_d(a))$.
For each positive radius $r$, write
\[
V_r(x)=\sum_{j:\,|\lambda_j(x-\mu_j)|=r}
\frac{w_j}{\lambda_j}\operatorname{sgn}(x-\mu_j).
\]
Zero radii contribute nothing. Since $J_d$ is odd, endpoint
separation gives, for $x\in\overline{\mathbb Q}\cap(a,b)$ and
$p\in\overline{\mathbb Q}\cap(0,1)$,
\begin{equation}
\label{beta-stats-eq-gaussian-vector}
C_d(x)=p
\quad\Longleftrightarrow\quad
V(x)-(1-p)V(a)-pV(b)=0.
\end{equation}

\begin{theorem}[Finite endpoint classification]
\label{beta-stats-thm-gaussian}
All algebraic--algebraic interior quantile pairs $(x,p)$ for $C_d$
belong to $K^2$ and form a finite set. They can be determined by
testing at most $5N^2$ possible locations and performing algebraic
arithmetic. The set of pairs is independent of the positive integer
$d$.
\end{theorem}

\begin{proof}
Let $\mathcal R$ be the nonzero endpoint radii
$|\lambda_j(a-\mu_j)|$ and $|\lambda_j(b-\mu_j)|$; then
$|\mathcal R|\le2N$. The possible locations are
\[
\mu_j,\qquad \mu_j\pm r/\lambda_j\quad(r\in\mathcal R),\qquad
\frac{\lambda_j\mu_j-\varepsilon\lambda_k\mu_k}
{\lambda_j-\varepsilon\lambda_k}\quad(j<k,\ \varepsilon=\pm1),
\]
with zero denominators, duplicates, and locations outside $(a,b)$
discarded. There are at most
$N+4N^2+2\binom N2=5N^2$ candidates, all in $K$.
At any other algebraic location, every moving radius is nonzero,
distinct from all other moving radii, and different from every
endpoint radius. Its coefficient in
\eqref{beta-stats-eq-gaussian-vector} is then uniquely
$\pm w_j/\lambda_j\ne0$, a contradiction.

For a surviving location put $D=V(b)-V(a)$. Positivity of the
normalizing integral implies $D\ne0$. A nonzero coordinate gives
the only possible probability,
$p=(V_r(x)-V_r(a))/D_r\in K$. Check all remaining coordinates and
$0<p<1$. The criterion does not contain $d$, proving invariance.
The bound counts candidates; it is not a polynomial bound on bit
complexity for algebraic-number computations.
\end{proof}

\subsection{Concrete quantiles and the limit of the method}

Consider the normalized density
\[
e^{-(x+1/2)^{2d}}+2e^{-(x-1/2)^{2d}},\qquad -1\le x\le1.
\]
Writing $J=J_d$ gives $S(x)=J(x+1/2)+2J(x-1/2)$ and the quantile
equation
\[
S(x)=(-1+3p)J(1/2)+(-2+3p)J(3/2).
\]
If $x,p$ are algebraic and $-1<x<1$, no moving radius can equal
$3/2$, so endpoint separation forces $p=2/3$. The only interior
location candidates are $-1/2,0,1/2$. At these points the residual
$S(x)-J(1/2)$ is respectively
$-2J(1)-J(1/2)$, $-2J(1/2)$, and $J(1)-J(1/2)$, none zero.
Thus \emph{every} quantile at an algebraic probability in $(0,1)$
is transcendental. At $d=1$ its median is
$0.111102019216\ldots$.

The same algorithm records genuine exceptions. For this mixture
on $[-3/2,3/2]$ instead, the unique algebraic--algebraic interior
pair is $(x,p)=(1/2,2/3)$: the outer endpoint radius forces $p=2/3$
and the three candidates are again $-1/2,0,1/2$, of which exactly
$1/2$ satisfies $S(x)=J(1)$. For a single density $e^{-x^{2d}}$
on $[0,1]$, every interior quantile at algebraic probability is
transcendental, directly by endpoint separation. In particular the
Gaussian truncated median defined by
$J_1(m_0)=J_1(1)/2$ is transcendental, with
$m_0=0.392697625792\ldots$.

For the untruncated standard normal law,
$\Phi(x)=1/2+J_1(x/\sqrt2)/\sqrt\pi$, the method proves only that
\emph{at most one} positive algebraic location can have an algebraic
distribution value. Two distinct such locations would eliminate
$\sqrt\pi$ and contradict endpoint separation. This does not
identify which of two specified quantiles is transcendental, nor
prove that a particular one, such as $\Phi^{-1}(39/40)$, is
transcendental. The finite-normalization theorem must not be
extended to that setting without an additional argument.

\clearpage

\part{Transfers to Euler constants and Gamma values}
\section{Gamma values and Euler's constant: transfer and obstruction}
\label{gamma-section}

The preceding beta-integral results motivate two different arithmetic
questions. One concerns an individual Gamma value, such as $\Gamma(1/5)$;
the other concerns Euler's constant
\[
 \gamma=\lim_{n\to\infty}(H_n-\log n)=-\Gamma'(1),
 \qquad H_n=\sum_{j=1}^n\frac1j,\quad H_0=0.
\]
Neither question is the functional differential transcendence of
$\Gamma(z)$, established by H\"older. In particular, the latter theorem
does not prevent the algebraic values $\Gamma(n)=(n-1)!$.
The discussion below proves explicit reductions, cancellation identities,
and a quantitative obstruction for one approximation scheme. It does not
prove that $\Gamma(1/5)$ is transcendental or that $\gamma$ is irrational.
The individual fifth-denominator question is still identified as open in
the 2026 account \cite{gamma-Sheth}. No priority claim is attached to the
auxiliary results assembled here.

\subsection{An exact fifth-denominator reduction}
\label{gamma-fifth}

Put
\[
 G_1=\Gamma(1/5),\qquad G_2=\Gamma(2/5),\qquad
 A=B(1/5,1/5),\qquad B=B(2/5,2/5).
\]
The beta--Gamma identity and reflection formula give
\begin{equation} A=\frac{G_1^2}{G_2},\qquad
 B=\frac{\sin(\pi/5)G_1G_2^2}{\pi},\qquad
 \boxed{G_1^5=\frac{\pi A^2B}{\sin(\pi/5)}}.
 \label{gamma-fifth-formula}
\end{equation}
Since $\sin(\pi/5)$ is nonzero and algebraic, $G_1$ is transcendental
if and only if $\pi A^2B$ is transcendental. The transcendence or linear
independence of separate beta values does not imply this product statement.
In particular, a theorem describing linear relations between ordinary
$1$-periods cannot be used as a theorem describing all polynomial
relations between their products. The known bound
\[
 \operatorname{trdeg}_{\mathbb Q}\mathbb Q(\pi,G_1,G_2)\ge2
\]
recorded in \cite{gamma-Rivoal} also does not select $G_1$: the bound is
consistent with $G_1$ algebraic and $\pi,G_2$ algebraically independent.

\subsection{Balanced products and the lost linear term}
\label{gamma-balanced}

The beta identity
\[
 B(a,b)=\frac{\Gamma(a)\Gamma(b)}{\Gamma(a+b)}
\]
is invariant under the formal modification
$\Gamma(z)\mapsto\Gamma_\kappa(z)=e^{\kappa z}\Gamma(z)$.
This modification is an algebraic device for comparing formulas, not a
second Gamma function satisfying all its original normalization,
recurrence, and growth conditions.

\begin{proposition}[Balanced-product invariance]
\label{gamma-balance-prop}
Suppose that $a_j(t)$ are holomorphic near $0$, their Gamma values avoid
poles there, and $m_j$ are integers. If
$\sum_jm_ja_j(t)$ is constant, then the normalized function
\[
 \frac{\prod_j\Gamma(a_j(t))^{m_j}}
      {\prod_j\Gamma(a_j(0))^{m_j}}
\]
is unchanged when $\Gamma$ is replaced by $\Gamma_\kappa$.
Consequently all its Taylor coefficients are invariant under that
modification.
\end{proposition}
\begin{proof}
The numerator and denominator acquire the respective factors
$\exp(\kappa\sum_jm_ja_j(t))$ and
$\exp(\kappa\sum_jm_ja_j(0))$, which are equal.
\end{proof}

For affine rational parameters this cancellation can be read directly
from the digamma function $\psi=\Gamma'/\Gamma$:
\begin{align}
 \frac{d}{dt}\log B(at+b,ct+d)
 &=a\psi(at+b)+c\psi(ct+d)\notag\\
 &\hspace{9mm}-(a+c)\psi((a+c)t+b+d).
 \label{gamma-beta-derivative}
\end{align}
At rational arguments away from poles, Gauss's formula and the shift
formula write $\psi(r)=-\gamma+L_r$, where $L_r$ is an explicit
combination of $\pi$, logarithms of algebraic numbers, and algebraic
coefficients. The coefficient of $\gamma$ in
\eqref{gamma-beta-derivative} is $-a-c+(a+c)=0$. For example,
\[
 \left.\frac{\partial}{\partial a}\log B(a,b)\right|_{a=b=1/3}
 =\psi(1/3)-\psi(2/3)=-\frac{\pi}{\sqrt3}.
\]
Thus differentiation of the balanced Gamma quotients treated earlier
does not by itself recover the linear term $-\gamma z$ in
$\log\Gamma(1+z)$. This is a limitation of that information, not a
theorem that $\gamma$ belongs to no field of ordinary periods.

There is also a geometric change when parameters are differentiated.
For instance,
\[
 \partial_h[x(1-x)]^h=[x(1-x)]^h\log(x(1-x)).
\]
If $f$ is a nonconstant rational function on a projective algebraic
curve, no branch of $\log f$ becomes a meromorphic function on a finite
algebraic cover. Otherwise its derivative would be an exact meromorphic
differential, which has zero residues, whereas the pullback of $df/f$
has residue $e\operatorname{ord}_P(f)\ne0$ at a point of ramification
index $e$ over a zero or pole of $f$. Hence the old algebraic-differential
argument cannot simply be reused after inserting this logarithm.

Likewise, the factor $e^{-t}$ in an Euler integral has an essential
singularity at infinity: in the coordinate $u=1/t$ it is $e^{-1/u}$.
A finite algebraic cover does not turn that essential singularity into
a finite-order pole. The generalized Jacobians with finite modulus used
for the earlier first- and second-kind differentials do not directly
represent this expression. This assertion concerns the displayed
integrands; it does not exclude every possible different representation
of $\gamma$.

\subsection{Confluence is exact, but transcendence need not survive it}
\label{gamma-confluence}

For $\Re s>0$, set
\[
 F_N(s)=N^sB(s,N)
 =\int_0^N t^{s-1}(1-t/N)^{N-1}\,dt.
\]
For $N\ge2$ the factor in parentheses, extended by zero past $N$, is
at most $e^{-t/2}$. Dominated convergence, locally uniformly with the
necessary logarithmic factors, gives
$F_N(s)\to\Gamma(s)$ and convergence of its derivatives on compact
subsets of $\Re s>0$. At $s=1$,
\begin{equation} F_N(1)=1,\qquad F_N'(1)=\log N-H_N\longrightarrow-\gamma.
 \label{gamma-confluence-derivative}
\end{equation}
An equivalent finite-part version uses
\[
 \widetilde F_N(s)=N^sB(s,N+1)
 =\frac{N^sN!}{s(s+1)\cdots(s+N)}
 =\frac1s+\log N-H_N+O(s)
\]
near $s=0$. These are genuine analytic connections between beta and
Gamma values. For $N\ge2$, each $H_N-\log N$ is transcendental by
Lindemann's theorem, but this property is not preserved by taking
limits. A separate arithmetic estimate is required.

\subsection{An antisymmetric exponential-integral transfer}
\label{gamma-antisymmetric}

A different limit arising in the earlier discussion connects an
antisymmetric iterated integral with both Euler's constant and a
dilogarithm. For real $r,R>0$ define
\[
 H_r(R)=\int_0^R\left(
 \operatorname{Ein}(rt)\frac{1-e^{-t}}t
 -\operatorname{Ein}(t)\frac{1-e^{-rt}}t\right)dt.
\]
Here $\operatorname{Ein}(z)=\sum_{n\ge1}(-1)^{n+1}z^n/(n\,n!)$
is entire.
The previously selected finite-point constant was $H_2(1)$.
For positive arguments it is convenient to define the real branch
of the dilogarithm by
$\operatorname{Li}_2(-r)=-\int_0^r\log(1+t)\,dt/t$.

\begin{proposition}[Renormalized iterated-integral limit]
\label{gamma-Hr-limit}
For every real $r>0$,
\begin{align}
 \lim_{R\to\infty}\bigl(H_r(R)-\log r\log R\bigr)
 &=\gamma\log r+\log^2r
   +2\operatorname{Li}_2(-r)+\frac{\pi^2}{6}\notag\\
 &=\gamma\log r-2\operatorname{Li}_2(-1/r)-\frac{\pi^2}{6}.
 \label{gamma-Hr-formula}
\end{align}
\end{proposition}
\begin{proof}
Differentiation under the finite integral and integration by parts give
\begin{align*}
 r\partial_r H_r(R)
 &=(1+e^{-rR})\operatorname{Ein}(R)
    +2\operatorname{Ein}(rR)-2\operatorname{Ein}((r+1)R).
\end{align*}
For completeness, the first differentiated integral is
$\operatorname{Ein}(R)+\operatorname{Ein}(rR)
-\operatorname{Ein}((r+1)R)$; the term subtracted from it is
$-e^{-rR}\operatorname{Ein}(R)
+\operatorname{Ein}((r+1)R)-\operatorname{Ein}(rR)$.
Using $\operatorname{Ein}(x)=\gamma+\log x+E_1(x)$, where
$E_1(x)=\int_x^\infty e^{-t}\,dt/t$, shows that
\[
 r\partial_r H_r(R)-\log R
 \longrightarrow\gamma+2\log r-2\log(1+r)
\]
uniformly for $r$ in any compact subinterval of $(0,\infty)$.
Since $H_1(R)=0$, integration with respect to $dr/r$ from $1$ to
$r$ gives the first limit formula, using
$\operatorname{Li}_2(-1)=-\pi^2/12$. The identity
\[
 \operatorname{Li}_2(-r)+\operatorname{Li}_2(-1/r)
 =-\frac{\pi^2}{6}-\frac12\log^2r
\]
gives the second form. It follows directly by differentiation in $r$
and evaluation at $r=1$.
\end{proof}

At $r=2$ the limit is
$\gamma\log2-2\operatorname{Li}_2(-1/2)-\pi^2/6$, approximately
$-0.3480102422994023475$. This is an exact analytic transfer, but
finite-point E-function transcendence theorems do not automatically
extend to this renormalized limit at infinity. Even transcendence of
the whole limit would not isolate $\gamma$ without additional control
of its relation with the dilogarithm and $\pi^2$ terms. Thus this
identity closes neither an individual Euler-constant problem nor a
dilogarithm independence problem.

\subsection{Polynomial logarithmic moments and a shared tail coefficient}
\label{gamma-moments}

The freedom to introduce a polynomial weight can be analyzed completely
for one logarithm and the fixed exponential weight. Let
$P(t)=\sum_{k=0}^dp_kt^k\in\overline{\mathbb Q}[t]$, and write
\[
 B_P=\sum_{k=0}^dp_kk!,\qquad
 A_P=\sum_{k=0}^dp_kk!H_k.
\]
Since $\Gamma'(k+1)=k!(H_k-\gamma)$, absolute convergence and linearity
give
\begin{equation}
 J_P:=\int_0^\infty P(t)e^{-t}\log t\,dt=A_P-B_P\gamma.
 \label{gamma-moment-full}
\end{equation}
If $B_P=0$, this integral is algebraic. If $B_P\ne0$, its algebraicity
is equivalent to the algebraicity of $\gamma$.

To expose the tail, introduce
\[
 Q_P(t)=\sum_{j=0}^dP^{(j)}(t),\qquad
 R_P(t)=\frac{Q_P(t)-Q_P(0)}t,
 \qquad C_P=\sum_{j\ge0}R_P^{(j)}(1).
\]
The last sum is finite, and is zero when $R_P=0$.
Then
\[
 Q_P-Q'_P=P,\qquad Q_P(0)=B_P,
 \qquad\int_0^\infty e^{-t}R_P(t)\,dt=A_P.
\]
The last identity follows for $P=t^k$ by summing $k!/(k-j)$ over
$j=0,\ldots,k-1$, and then follows generally by linearity.
Define
\[
 E_1(x)=\int_x^\infty\frac{e^{-t}}t\,dt,\qquad
 \eta=\operatorname{Ein}(1)=\gamma+E_1(1),
 \qquad\delta=eE_1(1).
\]
Integration by parts, using
$Pe^{-t}=-(Q_Pe^{-t})'$ and $\log1=0$, gives
\begin{align}
 T_P:=\int_1^\infty P(t)e^{-t}\log t\,dt
 &=B_PE_1(1)+\frac{C_P}{e},\label{gamma-moment-tail}\\
 L_P:=-\int_0^1P(t)e^{-t}\log t\,dt
 &=B_P\eta-A_P+\frac{C_P}{e}.
 \label{gamma-moment-finite}
\end{align}
Here the polynomial remainder is integrated by repeated integration by
parts. The identity $J_P=T_P-L_P$ provides a consistency check.

\begin{center}
\begin{tabular}{crrrl}
\toprule
$P(t)$&$A_P$&$B_P$&$C_P$&$J_P$\\
\midrule
$1$&0&1&0&$-\gamma$\\
$t$&1&1&1&$1-\gamma$\\
$t^2$&3&2&4&$3-2\gamma$\\
$t^3$&11&6&17&$11-6\gamma$\\
$t-1$&1&0&1&$1$\\
\bottomrule
\end{tabular}
\end{center}

In a finite linear combination, cancellation of the displayed
$E_1(1)$ coefficient in \eqref{gamma-moment-tail} forces the same
cancellation of the $\gamma$ coefficient in
\eqref{gamma-moment-full}. This is a statement about the explicit
coefficients and does not assume any unknown independence of $E_1(1)$
and $e$. In particular,
\[
 \int_0^\infty(t-1)e^{-t}\log t\,dt=1
\]
is a counterexample to the claim that every nonzero exponential
logarithmic integral is transcendental. The conclusion is restricted
to finite polynomial weights, one logarithm, this fixed exponential,
and linear cancellation. Rational weights or nonlinear constructions
are not covered.

\subsection{A positive kernel and a quantitative degree obstruction}
\label{gamma-positive}

Binet's formulas \cite{gamma-DLMF} yield, for $n\ge1$,
\begin{align}
 J_n&=2\gamma+2\log n-H_n-H_{n-1}\notag\\
 &=\int_0^\infty e^{-nt}\left(\coth(t/2)-\frac2t\right)dt
 =\int_0^1x^{n-1}K(x)\,dx,\label{gamma-Jn}\\
 K(x)&=\frac{1+x}{1-x}+\frac2{\log x},\qquad 0<x<1.
 \notag
\end{align}
The notation $J_n$ here refers to this kernel moment, rather than the
polynomial logarithmic moment $J_P$ of the preceding subsection.
For $P(x)=\sum_{n=1}^{N+1}c_nx^{n-1}\in\mathbb Q[x]$, set
\[
 C=P(1),\qquad A=\sum_{n=1}^{N+1}c_n(H_n+H_{n-1}),
 \qquad L(P)=\int_0^1P(x)K(x)\,dx.
\]
Exact logarithm cancellation gives
\begin{equation}
 \sum_nc_n\log n=0\quad\Longrightarrow\quad L(P)=2C\gamma-A.
 \label{gamma-log-cancel}
\end{equation}
The cancellation condition can be imposed without using $\gamma$:
after clearing rational denominators, unique factorization shows that
it is equivalent to
$\sum_nc_nv_\ell(n)=0$ for each prime $\ell\le N+1$.

\begin{proposition}[Pointwise kernel bound]
\label{gamma-kernel-bound}
For $0<x<1$, one has $K(x)>(1-x)/6$.
\end{proposition}
\begin{proof}
Let $f(x)=-\log x-12(1-x)/(5+8x-x^2)$. Direct differentiation gives
\[
 f'(x)=-\frac{(x-25)(x-1)^3}{x(x^2-8x-5)^2}<0,
 \qquad f(1)=0.
\]
Thus $f(x)>0$ for $x<1$. Rewriting this inequality and multiplying
only positive quantities gives the stated bound.
\end{proof}

\begin{theorem}[Positive-polynomial endpoint obstruction]
\label{gamma-degree-bound}
If $P$ is a real polynomial nonnegative on $[0,1]$ and $P(1)>0$, then
\[
 \frac{L(P)}{P(1)}\ge
 \begin{cases}
 \displaystyle\frac1{3(m+1)^2(m+2)^2},&\deg P\le2m,\\[5pt]
 \displaystyle\frac1{3(m+1)(m+2)^2(m+3)},&\deg P\le2m+1.
 \end{cases}
\]
If $P$ has rational coefficients and satisfies
\eqref{gamma-log-cancel}, the rational number $a_P=A/(2P(1))$
satisfies the corresponding bounds
\[
 \gamma-a_P\ge
 \begin{cases}
 \displaystyle\frac1{6(m+1)^2(m+2)^2},&\deg P\le2m,\\[5pt]
 \displaystyle\frac1{6(m+1)(m+2)^2(m+3)},&\deg P\le2m+1.
 \end{cases}
\]
\end{theorem}
\begin{proof}
The Markov--Luk\'acs representation for interval positivity
\cite{gamma-Schmudgen} gives, in even degree,
\[
 P(x)=U(x)^2+x(1-x)V(x)^2,
 \quad \deg U\le m,\quad\deg V\le m-1.
\]
For $m=0$ take $V=0$. The shifted Jacobi polynomials
$J_k^{*}(x)=P_k^{(1,0)}(2x-1)$ satisfy
\[
 J_k^{*}(1)=k+1,\qquad
 \int_0^1(1-x)J_k^{*}(x)^2\,dx=\frac1{2(k+1)}
\]
and are mutually orthogonal for this weight \cite{gamma-DLMF}.
Expanding $U$ and applying Cauchy--Schwarz therefore yields
\[
 U(1)^2\le
 \left(\sum_{k=0}^m2(k+1)^3\right)
 \int_0^1(1-x)U(x)^2\,dx.
\]
The sum is $(m+1)^2(m+2)^2/2$. Since $P(1)=U(1)^2$, dropping the
nonnegative $V$ contribution gives the required bound after
Proposition~\ref{gamma-kernel-bound}.

In odd degree write
$P=xU^2+(1-x)V^2$, with both polynomial degrees at most $m$.
Use instead $P_k^{(1,1)}(2x-1)$, whose value at $1$ is $k+1$ and
whose squared norm for the weight $x(1-x)$ is
\[
 \frac{k+1}{(2k+3)(k+2)}.
\]
The relevant endpoint sum is now
\[
 \sum_{k=0}^m(k+1)(k+2)(2k+3)
 =\frac{(m+1)(m+2)^2(m+3)}2.
\]
Since $\int(1-x)P\ge\int x(1-x)U^2$ and $P(1)=U(1)^2$, the same
argument proves the odd-degree bound. The rational approximation
statements follow by dividing \eqref{gamma-log-cancel} by $2P(1)$.
\end{proof}

The endpoint estimates for the weighted polynomial integrals are
optimal: equality in Cauchy--Schwarz is achieved by the corresponding
reproducing-kernel polynomial and the other squared term may be set
to zero. Optimality of the final estimate with the kernel $K$ is not
claimed. The proof also works if interval positivity is stated using
sums of squares instead of a single square in each term.

For a fully elementary weaker bound, Markov's inequality
$\lVert P'\rVert_\infty\le2N^2\lVert P\rVert_\infty$ on $[0,1]$
shows, when $N\ge1$, that
\begin{equation}
 \frac{L(P)}{P(1)}\ge\frac1{384N^4},
 \qquad\gamma-a_P\ge\frac1{768N^4}.
 \label{gamma-elementary-bound}
\end{equation}
Indeed, writing $M=\max P$, an interval of length
$\ell=1/(4N^2)$ fits on one side of a maximum point, and $P\ge M/2$
there. Among intervals of that length the smallest integral of $1-x$
is $\ell^2/2$. Hence
$\int_0^1(1-x)P\,dx\ge M/(64N^4)\ge P(1)/(64N^4)$.

These inequalities prohibit exponential-in-degree convergence in this
specific positive-polynomial scheme. If $a_P=p/q$ is reduced, the
criterion $q|\gamma-p/q|\to0$ necessarily requires $q/N^4\to0$.
That condition is not sufficient. Signed polynomials, rational kernels,
and other approximation methods are outside the theorem. A factor
$(1-x)^k$, $k>0$, gives $P(1)=0$ and removes the $\gamma$ coefficient
under exact logarithm cancellation; it cannot be used without checking
that loss of the target coefficient.

\subsection{A tested rational approximation family}
\label{gamma-test-family}

For integers $m\ge2$, choose
\[
 P_m(x)=2x^{m-1}-x^{m^2-1}>0\quad(0<x<1),\qquad P_m(1)=1.
\]
The identity $2\log m-\log(m^2)=0$ produces the rational number
\[
 a_m=\frac{2H_m+2H_{m-1}-H_{m^2}-H_{m^2-1}}2,
 \qquad\gamma-a_m=\frac{2J_m-J_{m^2}}2>0.
\]
Binet's second representation gives
\[
 J_n=4\int_0^\infty
 \frac{t\,dt}{(t^2+n^2)(e^{2\pi t}-1)}.
\]
Apply $1-u<(1+u)^{-1}<1$ and the standard zeta evaluations of
the first and third moments to obtain
\[
 \frac1{6n^2}-\frac1{60n^4}<J_n<\frac1{6n^2}.
\]
Consequently the error is bracketed by explicit rational numbers:
\begin{equation}
 \frac1{6m^2}-\frac1{10m^4}<\gamma-a_m
 <\frac1{6m^2}-\frac1{12m^4}+\frac1{120m^8}.
 \label{gamma-error-bracket}
\end{equation}

\begin{center}
\begin{tabular}{rrrr}
\toprule
$m$&Digits of reduced $q_m$&$\gamma-a_m$&
$\log_{10}(q_m|\gamma-a_m|)$\\
\midrule
2&2&$0.0355490$&$-0.0690$\\
4&7&$0.0100280$&$4.1600$\\
10&41&$0.00165667$&$38.0626$\\
16&110&$0.000649516$&$106.3357$\\
32&446&$0.000162665$&$441.6931$\\
\bottomrule
\end{tabular}
\end{center}

The numerators, reduced denominators, and rational brackets were
calculated exactly; the displayed decimals are numerical diagnostics.
The examples fail to produce a small integer linear form as $m$ grows
through the tested values. The table is not an asymptotic impossibility
proof. The uniform proved conclusion is
Theorem~\ref{gamma-degree-bound}.

\subsection{Exponential-integral values and the exact rank question}
\label{gamma-rank}

Recall that $\operatorname{Ein}(x)=\gamma+\log x+E_1(x)$ for $x>0$.
The Euler--Gompertz constant also has the equivalent expression
\[
 \delta=\int_0^\infty\frac{e^{-t}}{1+t}\,dt
 =\int_0^\infty e^{-t}\log(1+t)\,dt=eE_1(1).
\]
Rivoal proves that $e$ and $\eta=\operatorname{Ein}(1)$ are
algebraically independent \cite[\S5]{gamma-Rivoal}. Thus
\begin{align}
 3&\le\dim_{\overline{\mathbb Q}}
       \langle1,e,\delta,e\eta\rangle\le4,\notag\\
 \dim_{\overline{\mathbb Q}}\langle1,e,\delta,e\eta\rangle=3
 &\quad\Longleftrightarrow\quad
 \gamma\in\operatorname{span}_{\overline{\mathbb Q}}
 \{1,e^{-1},\eta\}.
 \label{gamma-rank-equivalence}
\end{align}
Indeed, $1,e,e\eta$ are linearly independent, and a rank drop is
equivalent to $\delta=A+Be+Ce\eta$ with algebraic $A,B,C$. Dividing
by $e$ gives $\gamma=(1-C)\eta-B-A/e$, and the converse reverses
this calculation. Rank four would prove $\gamma$ transcendental, but
was not established. If $\gamma$ were algebraic, the rank would be
three and $e,\delta$ would be algebraically independent, since
$\delta=e(\eta-\gamma)$ is an invertible rational change of variables.
This is fully consistent with the known transcendence-degree bound.

\subsection{Multiscale independence and an algebraic-limit counterexample}
\label{gamma-multiscale}

\begin{proposition}[A multiscale consequence of classical theorems]
\label{gamma-multiscale-prop}
If $n_1,\ldots,n_m$ are distinct positive integers, then
\[
 e,\operatorname{Ein}(n_1),\ldots,\operatorname{Ein}(n_m)
\]
are algebraically independent over $\overline{\mathbb Q}$.
\end{proposition}
\begin{proof}
Work in $K=\mathbb C(z,t)$ with $t=e^{-z}$ and derivation
$D=\partial_z-t\partial_t$. The functions $z,t$ are algebraically
independent. Put $y_j=\operatorname{Ein}(n_jz)$; then
\[
 Dy_j=\frac{1-t^{n_j}}z.
\]
The primitive case of the Kolchin--Ostrowski theorem
\cite[Theorem 2.5]{gamma-Srinivasan} says that algebraic dependence
of the $y_j$ over $K$ would give a nonzero constant linear combination
$\sum_jc_jy_j=R\in K$. There are no new differential constants: the
fields in question embed in the common field of meromorphic functions,
whose constants are $\mathbb C$. Differentiating gives
\[
 DR=\frac{\sum_jc_j(1-t^{n_j})}{z}.
\]
The numerator is a nonzero polynomial in the independent variable $t$.
Expand $R$ at the divisor $z=0$ over $\mathbb C(t)$. If $R$ has a
pole of order $d\ge1$, $\partial_zR$ has a nonzero pole of order
$d+1$, while $-t\partial_tR$ has order at most $d$. If $R$ is regular,
$DR$ is regular. A derivative therefore cannot have precisely the
simple pole displayed above. This contradiction proves independence
of the primitives over $K$.

All $y_j$ and $t$ are E-functions. Adding $1$ and $e^{-n_jz}$ gives
a homogeneous first-order system over $\mathbb Q(z)$ whose only finite
singularity is $0$. The extra exponentials are powers of $t$, so they
do not increase the functional transcendence degree. The
Siegel--Shidlovskii theorem at $z=1$, in the form recalled in
\cite{gamma-Effective}, gives degree $m+1$ for the values. Replacing
$e^{-1}$ by its inverse finishes the proof.
\end{proof}

In particular, for
\[
 Y_n=(n+1)\operatorname{Ein}(2^n)-n\operatorname{Ein}(2^{n+1}),
 \qquad n\ge1,
\]
every finite collection $e,Y_1,\ldots,Y_N$ is algebraically independent.
The linear transformation of the corresponding $\operatorname{Ein}$
values has full row rank and can be completed to an invertible one.
Logarithm cancellation gives
\begin{align*}
 Y_n&=\gamma+\varepsilon_n,\\
 \varepsilon_n&=(n+1)E_1(2^n)-nE_1(2^{n+1})
 \sim\frac{n+1}{2^n}e^{-2^n}>0.
\end{align*}
Positivity follows from $E_1(2x)<E_1(x)$, and the asymptotic follows
from $E_1(x)\sim e^{-x}/x$. These conclusions still do not determine
the arithmetic nature of the limit.

\begin{proposition}[Matching algebraic limits]
\label{gamma-limit-counterexample}
For each real algebraic $a$, let $p_n$ be the $n$th prime and define
\[
 Z_n=a+\frac{n+1}{2^n}
 \exp\left(-2^n+\frac{\sqrt{p_n}}{2^n}\right).
\]
Then $Z_n$ decreases strictly to $a$,
$Z_n-a\sim(n+1)2^{-n}e^{-2^n}$, and every finite collection
$e,Z_1,\ldots,Z_N$ is algebraically independent.
\end{proposition}
\begin{proof}
Bertrand's postulate implies $p_n\le2^n$, so the correction
$\sqrt{p_n}/2^n$ tends to zero. The prefactors $(n+1)/2^n$ decrease
strictly, and the difference of successive exponents is less than
$-2^n+1<0$. This proves the analytic assertions.

The algebraic numbers
$1,-2^1+\sqrt{p_1}/2^1,\ldots,-2^N+\sqrt{p_N}/2^N$ are linearly
independent over $\mathbb Q$: apply the independent square-root sign
changes in the multiquadratic field. Distinct monomials in their
exponentials consequently have distinct algebraic exponents.
Lindemann--Weierstrass makes those exponentials linearly independent
over $\overline{\mathbb Q}$, ruling out any nonzero polynomial
relation. Algebraic shifts and nonzero rational scalings preserve
algebraic independence.
\end{proof}

This counterexample matches finite algebraic independence, positivity,
monotonicity, and the leading error scale. It excludes the proposed
inference from those properties alone, not a method using stronger
arithmetic information about the actual $Y_n$. The related ``critical
slope'' identity
\[
 \lim_{n\to\infty}\frac{-\log|P(Y_n)|}{2^n}
 =\operatorname{ord}_{X=\gamma}P
 \qquad(0\ne P\in\mathbb C[X])
\]
is just Taylor expansion at the limit; the same identity holds for
$Z_n$ and $a$. It supplies no independent exclusion of an algebraic
limit.

\subsection{Two arithmetic blocks and a fixed differential equation}
\label{gamma-gevrey}

The following reduction makes the connection constant explicit. Set
\[
 A(z)=e^z\operatorname{Ein}(z)=\sum_{n\ge0}\frac{H_n}{n!}z^n,
 \qquad\widehat D(z)=\sum_{n\ge0}(-1)^nn!z^{-n-1}.
\]
The first is an E-function. The second is divergent for every fixed
nonzero complex $z$, but has the positive-direction Borel--Laplace sum
\[
 D_0(z)=\int_0^\infty\frac{e^{-t}}{z+t}\,dt=e^zE_1(z),
 \qquad z>0.
\]
We use the mixed arithmetic-Gevrey framework of
Fischler--Rivoal \cite{gamma-Gevrey}; in particular, formal statements
about its two parts are distinguished from statements about one
analytic sum. For algebraic $r$ define
\[
 \widehat\Psi_r=A-re^z-\widehat D.
\]
Its positive sum satisfies the exact identity
\begin{equation}
 \Psi_{r,0}(z)=e^z(\gamma-r+\log z),\qquad
 \Psi_{r,0}(1)=e(\gamma-r).
 \label{gamma-connection-value}
\end{equation}
Since
\[
 A'=A+\frac{e^z-1}{z},\qquad D'_0=D_0-\frac1z,
\]
all these sums satisfy
\begin{equation} \Psi'_{r,0}=\Psi_{r,0}+\frac{e^z}{z},\qquad
 zY''+(1-2z)Y'+(z-1)Y=0.
 \label{gamma-second-order}
\end{equation}
The second-order equation has local solution basis $e^z,e^z\log z$.
The number $\gamma-r$ is a connection constant selecting a particular
arithmetic solution; it is not a coefficient in this rational
differential equation. Existence of the equation therefore proves no
arithmetic property of that constant.

\begin{proposition}[Formal division obstruction]
\label{gamma-division-obstruction}
For algebraic $r$, the formal expression $\widehat\Psi_r$ cannot be
written as
\[
 (z-1)\left(F(z)+\sum_{n\ge1}b_nz^{-n}\right)
\]
where $F$ is an E-function and every $b_n$ is algebraic. In particular,
no such mixed arithmetic-Gevrey quotient exists.
\end{proposition}
\begin{proof}
The nonnegative powers would give
\[
 A(z)-re^z=(z-1)F(z)+b_1.
\]
Since $F$ is entire, evaluation at $1$ gives
$b_1=e\eta-re$. This contradicts the algebraic independence of
$e,\eta$: otherwise $XY-rX-b_1$ would be a nonzero algebraic
polynomial relation between them.
\end{proof}

What remains unproved is the lifting implication
\begin{equation}
 \Psi_{r,0}(1)=0\quad\Longrightarrow\quad
 \widehat\Psi_r=(z-1)\times
 \text{a mixed arithmetic-Gevrey expression}.
 \label{gamma-zero-lifting}
\end{equation}
An analytic zero permits division on that analytic branch, but does
not ensure the arithmetic coefficient and denominator conditions of
the quotient. A general zero-removal assertion of this kind is
Conjecture 3 of \cite{gamma-Gevrey}. For the displayed family,
requiring \eqref{gamma-zero-lifting} for all algebraic $r$ is already
equivalent to $\gamma$ being transcendental: one direction follows
from Proposition~\ref{gamma-division-obstruction} and
\eqref{gamma-connection-value}; in the other direction its antecedent
never occurs. Thus this implication is not an independently proved
lemma that lowers the difficulty of the target.

There is a concrete branch issue as well. Under the hypothetical
choice $r=\gamma$, the quotient of the positive sum is
$e^z\log z/(z-1)$, with a removable singularity at $1$. Continuation
around $0$ adds $2\pi i e^z/(z-1)$, which has a pole at $1$ on the
new branch. This observation explains why one-branch removability must
not be confused with a global arithmetic assertion; it is not itself
a transcendence proof.

\subsection{Disjoint direct $\ell$-adic convergence domains}
\label{gamma-padic}

The natural two formal blocks do not admit simultaneous direct
evaluation at a nonzero $\ell$-adic input. This can be checked exactly,
including their common boundary.

\begin{proposition}[Complementary convergence domains]
\label{gamma-padic-prop}
Fix a prime $\ell$, normalize $v_\ell(\ell)=1$, and fix an algebraic
$\ell$-adic embedding of $r$. For nonzero $z$ in a complete extension,
the respective convergence conditions are exactly
\[
 \sum_{n\ge0}\frac{H_n-r}{n!}z^n:
 \quad v_\ell(z)>\frac1{\ell-1},
 \qquad
 \widehat D(z):\quad v_\ell(z)<\frac1{\ell-1}.
\]
Both series diverge on the boundary.
\end{proposition}
\begin{proof}
Legendre's formula and the elementary harmonic-number estimate give
\[
 v_\ell(n!)=\frac{n-s_\ell(n)}{\ell-1},\qquad
 v_\ell(H_n)\ge-\lfloor\log_\ell n\rfloor.
\]
These imply convergence in the indicated open domains: the linear
term in $n$ dominates the logarithmic correction. To obtain the
reverse statements and test the boundary, take $n=\ell^k$. The unique
term of smallest valuation in $H_{\ell^k}$ is $1/\ell^k$, so
$v_\ell(H_{\ell^k})=-k$ and, for all sufficiently large $k$,
$v_\ell(H_{\ell^k}-r)=-k$. When
$v_\ell(z)=1/(\ell-1)$, the corresponding E-series term has valuation
\[
 -k-\frac{\ell^k-1}{\ell-1}+\frac{\ell^k}{\ell-1}
 =-k+\frac1{\ell-1}.
\]
The negative-power factorial term has valuation
\[
 \frac{\ell^k-1}{\ell-1}-\frac{\ell^k+1}{\ell-1}
 =-\frac2{\ell-1}.
\]
Neither sequence of terms tends to zero. On the wrong side of either
boundary, the same subsequence has a strictly negative linear trend,
so convergence also fails there.
\end{proof}

This result rules out direct simultaneous evaluation of these two
particular series. It does not rule out $\ell$-adic analytic
continuation, Frobenius structures, overconvergence, or another local
construction. A local nonvanishing computation needs an independently
proved comparison theorem before it can yield a conclusion about
the complex connection constant $\gamma$.

\subsection{The Euler exponential motive and an exact lifting equivalence}
\label{gamma-motive}

Fres\'an--Jossen construct an exponential motive $M(\gamma)$ and prove
\cite[Proposition 12.8.7]{gamma-FresanJossen} that
\[
 0\longrightarrow\mathbb Q(0)\longrightarrow M(\gamma)
 \longrightarrow\mathbb Q(-1)\longrightarrow0
\]
is nonsplit, with
$\operatorname{Hom}_{\rm mot}(M(\gamma),\mathbb Q(0))=0$.
Their Corollary 12.8.8 identifies its motivic Galois group with
$\mathbb G_a\rtimes\mathbb G_m$ and makes the algebraic independence
of $\gamma,\pi$ conditional on the exponential period conjecture.
Geometric nonsplitting alone is not that conjecture.

The arithmetic content of a proposed lift can be measured directly.
Consider the unfiltered de Rham--Betti comparison category, with a
$\overline{\mathbb Q}$ de Rham space, a rational Betti space, and the
comparison matrix
\[
 P_\gamma=\begin{pmatrix}1&\gamma\\0&2\pi i\end{pmatrix},
 \qquad\operatorname{comp}(d_0)=b_0,\quad
 \operatorname{comp}(d_1)=\gamma b_0+2\pi i b_1.
\]
A compatible morphism to $\mathbb Q(0)$ has rational Betti row
$(a,b)$ and de Rham row $(a,a\gamma+2\pi i b)$. It exists exactly
when the second entry is algebraic. Since $\gamma$ is real, subtracting
the complex conjugate of that entry forces $b=0$ by the transcendence
of $\pi$. A nonzero map then has $a\ne0$, forcing $\gamma$ algebraic.
Conversely, algebraic $\gamma$ gives the compatible rows $(1,0)$ and
$(1,\gamma)$. Hence
\begin{equation}
 \boxed{\operatorname{Hom}_{\rm dR,B}
 (M(\gamma),\mathbb Q(0))=0
 \quad\Longleftrightarrow\quad\gamma\text{ is transcendental}.}
 \label{gamma-DB-equivalence}
\end{equation}
Because the motivic Hom space is zero, surjectivity of the single map
\[
 \operatorname{Hom}_{\rm mot}(M(\gamma),\mathbb Q(0))
 \longrightarrow
 \operatorname{Hom}_{\rm dR,B}(M(\gamma),\mathbb Q(0))
\]
is already equivalent to the desired transcendence.

Equivalently, for algebraic $r$ the line
$\overline{\mathbb Q}(d_1-rd_0)$ compares to $\mathbb Cb_1$ exactly
when $r=\gamma$. If this pair lifted to a submotive, it would split
the extension. Nonsplitting excludes a motivic complement, not every
numerically compatible pair. Full exponential Hodge realization has
additional filtration and sheaf data; a theorem about that realization
must not be replaced by a claim about the bare numerical matrix. The
ordinary $1$-period theorems used earlier establish no such lift for
this exponential motive.

\subsection{Conditional closure and verified scope}
\label{gamma-scope}

Let $\mathcal E$ be the ring of values of E-functions at algebraic
points and $\mathcal D$ the ring of the appropriate sums of arithmetic
Gevrey series of order $+1$, with the definitions of
\cite{gamma-Gevrey}. The assertion
\begin{equation}
 \mathcal E\cap\mathcal D=\overline{\mathbb Q}
 \label{gamma-intersection-conjecture}
\end{equation}
is a conjecture. Its implication for $\gamma$ is transparent:
$\delta\in\mathcal D$, while $e,\eta\in\mathcal E$. If $\gamma$
were algebraic, $\delta=e\eta-e\gamma$ would also belong to
$\mathcal E$. Equation~\eqref{gamma-intersection-conjecture} would
make $\delta$ algebraic, contradicting the algebraic independence of
$e,\eta$. This is a conditional argument, not a proof of the
intersection assertion. Fischler--Rivoal's Theorem 2 already gives
the broader conditional transcendence of $\Gamma^{(m)}(a)$ at
positive rational $a$, except the factorial case $m=0$, $a$ a
positive integer.

Results about infinitely many Gamma derivatives do not select
$-\Gamma'(1)$. For example,
\[
 \Gamma''(1)=\gamma^2+\frac{\pi^2}{6}
\]
would be transcendental even under the hypothesis that $\gamma$ is
algebraic. The 2026 derivative results of Powers
\cite{gamma-Powers} are therefore not treated as a solution of this
individual target and are not assumptions in the proofs above.

Exact verification of the degree argument checked the kernel
derivative, both Jacobi norm formulas and endpoint extremizers for
$m=0,\ldots,5$, and the rational denominators and brackets in
\eqref{gamma-error-bracket}. The separate arithmetic-block audit
checked the polynomial identities for $58$ weights, $81$ convolution
coefficients of $A$, $81$ first-order coefficients, and $484$
second-order Laurent coefficient identities over four rational
choices of $r$. Boundary valuations were checked for primes
$2,3,5,7$ and $k=1,2,3,4$. Floating-point integral comparisons are
diagnostic checks, not interval certificates. The all-degree and
all-prime conclusions rely on the proofs, not these finite tests.

The completed conclusions are the explicit Gamma-value reduction,
the antisymmetric iterated-integral limit, balanced and weighted
cancellations, the positive-polynomial degree
bound, the multiscale independence consequence, the matching
algebraic-limit counterexample, the formal division obstruction,
the disjoint direct $\ell$-adic domains, and the exact comparison
lifting equivalence. The missing assertions are arithmetic control
of the Gamma product, a sufficiently strong new approximation
mechanism, or an applicable exponential-period relation theorem.
None is replaced here by an unproved assumption disguised as a lemma.

\section{Concluding status and reproducibility}
The positive results above are specific transcendence statements and
classifications obtained from cited arithmetic theorems and explicit
calculations. The Euler-constant chapter has a different outcome: it
identifies exact cancellations, counterexamples to proposed limit arguments,
and precise arithmetic conditions still missing from a proof.
Changing the name of a condition equivalent to Euler transcendence does
not turn it into an established intermediate theorem.

The finite catalogues, polynomial identities, Frobenius polynomials,
cyclotomic calculations, and proof certificates needed to identify the
asserted cases are incorporated into this document. Decimal values are
orientation data, not evidence for transcendence. Printed algorithms are
finite exact procedures; where their results enter a proof, the required
inputs and arithmetic tests are specified. Independent internal checking
does not constitute a formal proof-assistant verification or external peer
review.

The accompanying source is one self-contained \LaTeX\ file. It embeds the
entire body and bibliography and uses no external images, tables, code files,
or bibliography database. With a standard TeX Live or MiKTeX installation,
run \texttt{pdflatex} twice on that file to resolve the contents and references.
The only two delivered files are this PDF and that source.

\clearpage
\addcontentsline{toc}{section}{References}
\begin{thebibliography}{99}
\bibitem{dbkz-source}
R. Dougherty-Bliss, C. Koutschan and D. Zeilberger,
\emph{Tweaking the Beukers integrals in search of more miraculous
irrationality proofs \`a la Ap\'ery}, Ramanujan J. \textbf{58} (2022),
973--994, arXiv:2101.08308.
Accompanying exact output:
\url{https://sites.math.rutgers.edu/~zeilberg/tokhniot/oGenBeukersLog2.txt}.

\bibitem{dbkz-hw}
A.~Huber and G.~W\"ustholz,
\emph{Transcendence and Linear Relations of 1-Periods},
Cambridge Tracts in Mathematics~227, Cambridge University Press, 2022.
Author manuscript, Theorem~6.2, Theorem~13.9 and Corollary~13.13.
\url{https://home.mathematik.uni-freiburg.de/arithgeom/preprints/huber-einsmotive.pdf}.

\bibitem{dbkz-ww}
J. Wolfart and G. W\"ustholz,
\emph{Der \"Uberlagerungsradius gewisser algebraischer Kurven und die Werte
der Betafunktion an rationalen Stellen}, Math. Ann. 273 (1985), 1--15.
Satz 1, Satz 4, and Propositions 1--2.
\url{https://doi.org/10.1007/BF01455911}.

\bibitem{dbkz-hz}
E. W. Howe and H. J. Zhu,
\emph{On the existence of absolutely simple abelian varieties of a
given dimension over an arbitrary field},
J. Number Theory \textbf{92} (2002), 139--163.
\url{https://arxiv.org/abs/math/0002205}.

\bibitem{log1-dbkz}
R.~Dougherty-Bliss, C.~Koutschan, and D.~Zeilberger,
\emph{Generalized Beukers integrals}, associated computational output
\texttt{oGenBeukersLog1.txt},
\url{https://sites.math.rutgers.edu/~zeilberg/tokhniot/oGenBeukersLog1.txt}.

\bibitem{log1-grant-shaulis}
D.~Grant and D.~Shaulis,
\emph{The cuspidal torsion packet on hyperelliptic Fermat quotients},
J.~Th\'eor. Nombres Bordeaux \textbf{16} (2004), 577--585.
\url{https://www.numdam.org/item/10.5802/jtnb.462.pdf}.

\bibitem{log1-boxall-grant}
J.~Boxall and D.~Grant,
\emph{Examples of torsion points on genus two curves},
Trans. Amer. Math. Soc. \textbf{352} (2000), 4533--4555.
\url{https://euclid.colorado.edu/~grant/publications/examples.pdf}.

\bibitem{log1-chudnovsky}
G.~V.~Chudnovsky,
\emph{Algebraic independence of values of exponential and elliptic
functions}, Proceedings of the International Congress of Mathematicians,
Helsinki 1978, vol.~1, 339--350.
\url{https://www.mathunion.org/fileadmin/ICM/Proceedings/ICM1978.1/ICM1978.1.ocr.pdf}.

\bibitem{log1-sertoz}
E.~C.~Sert\"oz, J.~Ouaknine, and J.~Worrell,
\emph{Computing transcendence and linear relations of 1-periods},
arXiv:2505.20397 (2025),
\url{https://arxiv.org/abs/2505.20397}.

\bibitem{extra8-dbkz}
R.~Dougherty-Bliss, C.~Koutschan and D.~Zeilberger,
\emph{Tweaking the Beukers Integrals in Search of More Miraculous
Irrationality Proofs \`a la Ap\'ery}, arXiv:2101.08308 (2021),
denominator-seven list, p.~17.
\url{https://arxiv.org/abs/2101.08308}.

\bibitem{extra8-zudilin}
W. Zudilin, \emph{The birthday boy problem}, arXiv:2108.06586v2 (2023),
Mellin--Barnes representation and its Barnes-lemma specializations.
\url{https://arxiv.org/abs/2108.06586}.

\bibitem{extra8-schneider}
Th.~Schneider,
\emph{Zur Theorie der Abelschen Funktionen und Integrale},
J.~reine angew.~Math.~183 (1941), 110--128.
\url{https://doi.org/10.1515/crll.1941.183.110}.

\bibitem{extra8-waldschmidt}
M.~Waldschmidt,
\emph{Transcendence of Periods: The State of the Art},
Pure Appl.~Math.~Q.~2 (2006), 435--463, Theorem~12.
\url{https://webusers.imj-prg.fr/~michel.waldschmidt/articles/pdf/TranscendencePeriods.pdf}.

\bibitem{extra8-dhillon-saha}
S.~Dhillon and E.~Saha,
\emph{On arithmetic nature of special values of the incomplete beta
function}, Acta Arith.~218 (2025), 273--284.
\url{https://doi.org/10.4064/aa240610-6-10}.

\bibitem{extra8-barnes}
I.~D.~Mishev, \emph{An invariance group for a linear
combination of two Saalsch\"utzian $ {}_4F_3(1)$ hypergeometric series},
arXiv:0910.0093, Eq.~(5.6) (Barnes' second lemma).
\url{https://arxiv.org/abs/0910.0093}.

\bibitem{extra8-dlmf}
NIST Digital Library of Mathematical Functions,
Eq.~16.4.6 (Watson's sum).
\url{https://dlmf.nist.gov/16.4.E6}.

\bibitem{extra7-dlmf}
NIST Digital Library of Mathematical Functions,
Eq.~16.4.7 (Whipple's sum).
\url{https://dlmf.nist.gov/16.4.E7}.

\bibitem{fields-MSE}
Wolfgang, \emph{How to determine the algebraic structure of the space of
algebraic gamma quotients?}, Mathematics Stack Exchange question 4959377,
16 August 2024, revised 25 August 2024.
\url{https://math.stackexchange.com/questions/4959377/}.

\bibitem{fields-Kushel}
O. Y. Kushel, \emph{Generalized Brouncker's continued fractions and their
logarithmic derivatives}, arXiv:1301.3734, Theorem 1.4.
\url{https://arxiv.org/abs/1301.3734}.

\bibitem{fields-Kenison}
G. Kenison et al., \emph{On Positivity and Minimality for Second-Order
Holonomic Sequences}, arXiv:2007.12282, Corollary 2.9, Section 5, and
Appendix B. \url{https://arxiv.org/abs/2007.12282}.

\bibitem{fields-Baker}
A. Baker, \emph{Transcendental Number Theory}, Cambridge University Press,
1975. Linear forms in logarithms of algebraic numbers.

\bibitem{fields-DLMF}
NIST Digital Library of Mathematical Functions, Sections 5.4--5.5,
in particular Gauss's digamma formula, equation 5.4.19.
\url{https://dlmf.nist.gov/5.4#E19}; \url{https://dlmf.nist.gov/5.5}.

\bibitem{beta-tamiozzo}
M.~Tamiozzo, \emph{Algebraicity of the division points of the trifolium
and related topics}, arXiv:1912.00204, 2019.
\url{https://arxiv.org/abs/1912.00204}.

\bibitem{beta-carocca}
A.~Carocca and R.~E.~Rodr\'iguez, Jacobians with group actions and rational
idempotents, \emph{Journal of Algebra} \textbf{306} (2006), 322--343.
\url{https://arxiv.org/abs/math/0305328}.

\bibitem{beta-arul}
V.~Arul, \emph{Torsion points on Fermat quotients of the form
$y^n=x^d+1$}, arXiv:1910.14251v2, 2020, Theorem~6.2.21.
\url{https://arxiv.org/abs/1910.14251}.

\bibitem{beta-boxall-grant}
J.~Boxall and D.~Grant, Examples of torsion points on genus two curves,
\emph{Transactions of the American Mathematical Society} \textbf{352}
(2000), 4533--4555, Theorem~(0.3).
\url{https://doi.org/10.1090/S0002-9947-00-02368-0}.

\bibitem{beta-roessler}
D.~Roessler, \emph{A note on the Manin--Mumford conjecture},
arXiv:math/0409083, 2004.
\url{https://arxiv.org/abs/math/0409083}.

\bibitem{beta-zarhin}
Yu.~G.~Zarhin, \emph{$p$-adic abelian integrals and commutative Lie groups},
arXiv:alg-geom/9603006, 1996, Sections~1 and~3.
\url{https://arxiv.org/abs/alg-geom/9603006}.

\bibitem{beta-voutier}
P.~M.~Voutier, Primitive divisors of Lucas and Lehmer sequences, II,
\emph{Journal de Th\'eorie des Nombres de Bordeaux} \textbf{8} (1996),
251--274, Theorem~2(i).
\url{https://doi.org/10.5802/jtnb.168}; author manuscript
\url{https://arxiv.org/abs/1211.3107}.

\bibitem{beta-zannier}
U.~Zannier, \emph{Hyperelliptic continued fractions and generalized
Jacobians}, arXiv:1602.00934, 2016, Section~2.1.2.
\url{https://arxiv.org/abs/1602.00934}.

\bibitem{beta-agostini}
D.~Agostini, T.~\"O.~\c{C}elik and J.~B.~Little,
Polynomial analogs of theta functions and rational solutions of the KP
equation, \emph{International Mathematics Research Notices}
\textbf{2024} (2024), no.~22, 13964--13987, Sections~2.1 and~2.3.
\url{https://doi.org/10.1093/imrn/rnae227}.

\bibitem{beta-dlmf}
NIST Digital Library of Mathematical Functions, Section~8.17,
\emph{Incomplete Beta Functions}.
\url{https://dlmf.nist.gov/8.17}.

\bibitem{beta-stats-baker67}
A. Baker, \emph{Linear forms in the logarithms of algebraic numbers
(III)}, Mathematika \textbf{14} (1967), 220--228.
\href{https://doi.org/10.1112/S0025579300003843}{doi:10.1112/S0025579300003843}.

\bibitem{beta-stats-baker66}
A. Baker, \emph{Linear forms in the logarithms of algebraic numbers},
Mathematika \textbf{13} (1966), 204--216.
\href{https://doi.org/10.1112/S0025579300003971}{doi:10.1112/S0025579300003971}.

\bibitem{beta-stats-vdp}
A. J. van der Poorten, \emph{On the arithmetic nature of definite
integrals of rational functions}, Proc. Amer. Math. Soc. \textbf{29}
(1971), 451--456.
\href{https://doi.org/10.1090/S0002-9939-1971-0276180-X}{doi:10.1090/S0002-9939-1971-0276180-X}.

\bibitem{beta-stats-ratesci}
P. Laud, \emph{ratesci: Confidence Intervals for Comparisons of
Binomial or Poisson Rates}, official R package documentation,
vignette ``Single rate'' and reference entry ``jeffreysci''.
\url{https://cran.r-project.org/web/packages/ratesci/vignettes/single_rate.html};
\url{https://search.r-project.org/CRAN/refmans/ratesci/html/jeffreysci.html}.

\bibitem{beta-stats-dlmf}
NIST Digital Library of Mathematical Functions,
\S5.5, \emph{Functional Relations}, and \S5.12, \emph{Beta Function}.
\url{https://dlmf.nist.gov/5.5}; \url{https://dlmf.nist.gov/5.12}.

\bibitem{beta-stats-beukers}
F. Beukers, \emph{A refined version of the Siegel--Shidlovskii theorem},
Ann. of Math. (2) \textbf{163} (2006), 369--379.
\href{https://doi.org/10.4007/annals.2006.163.369}{doi:10.4007/annals.2006.163.369}.

\bibitem{gamma-Sheth}
A. Sheth,
\emph{Real Geometric transcendence for the Gamma function},
arXiv:2605.15117v1 (2026), especially \S4.2.
\url{https://arxiv.org/abs/2605.15117}.

\bibitem{gamma-Rivoal}
T. Rivoal,
\emph{On the arithmetic nature of the values of the Gamma function,
Euler's constant and Gompertz's constant},
Michigan Math. J. \textbf{61} (2012), 239--254,
Theorem 2 and \S5.
\url{https://rivoal.perso.math.cnrs.fr/articles/gammater.pdf}.

\bibitem{gamma-DLMF}
NIST Digital Library of Mathematical Functions,
\S\S5.7, 5.9, 5.12, 6.2, and 18.3:
Gamma and beta identities, Binet's representations,
exponential-integral relations, and Jacobi orthogonality.
\url{https://dlmf.nist.gov/}.

\bibitem{gamma-Schmudgen}
K. Schm\"udgen,
\emph{The Moment Problem},
Graduate Texts in Mathematics 277, Springer, 2017,
Corollary 3.24 (Markov--Luk\'acs interval positivity).
\url{https://doi.org/10.1007/978-3-319-64546-9}.

\bibitem{gamma-Srinivasan}
V. R. Srinivasan,
\emph{Extensions by Antiderivatives, Exponentials of Integrals and
by Iterated Logarithms},
arXiv:0811.3004, Theorem 2.5.
\url{https://arxiv.org/abs/0811.3004}.

\bibitem{gamma-Effective}
S. Fischler and T. Rivoal,
\emph{Effective algebraic independence of values of E-functions},
author manuscript, Theorem 1.
\url{https://rivoal.perso.math.cnrs.fr/articles/iaeff.pdf}.

\bibitem{gamma-Gevrey}
S. Fischler and T. Rivoal,
\emph{Relations between values of arithmetic Gevrey series, and
applications to values of the Gamma function},
J. Number Theory \textbf{261} (2024), 36--54,
Conjectures 1 and 3, Theorem 2, and Proposition 2.
\url{https://doi.org/10.1016/j.jnt.2024.02.016}.
Author version: \url{https://arxiv.org/abs/2301.13518}.

\bibitem{gamma-FresanJossen}
J. Fres\'an and P. Jossen,
\emph{Exponential motives}, author manuscript,
\S12.8, especially Proposition 12.8.7 and Corollary 12.8.8.
\url{https://javier.fresan.perso.math.cnrs.fr/expmot.pdf}.

\bibitem{gamma-Powers}
M. R. Powers,
\emph{Transcendence Results for $\Gamma^{(n)}(1)$ and Related Sequences
of Generalized Constants},
arXiv:2511.01849v7 (13 April 2026).
\url{https://arxiv.org/html/2511.01849v7}.

\end{thebibliography}
\end{document}
